% !TeX program = xelatex
%DIF LATEXDIFF DIFFERENCE FILE
%DIF DEL Public preprint of 29 October 2026, 13:33:54 Paris
%DIF ADD Bibliographic correction of 29 October 2026, 14:29:39 Paris
% Source of the preprint "Deep linear networks are mildly singular at every exact fit".
% The release includes a self-contained flattened .tex file. Compile it with XeLaTeX three times.
% Lines "%%INPUT file" are inlined; lines "%%LEAN path:a-b" are replaced by a verbatim copy
% of lines a..b of that Lean file, followed by its path, line numbers and SHA-256 prefix.
\documentclass[11pt,a4paper]{article}
\usepackage[margin=27mm]{geometry}
\setlength{\parskip}{0.3em}
\usepackage{fontspec}
\setmainfont{Latin Modern Roman}
\setmonofont{DejaVu Sans Mono}[Scale=MatchLowercase]
\usepackage{amsmath,amssymb,amsthm,mathtools}
\usepackage{enumitem,booktabs,array}
\usepackage{xcolor}
\usepackage{fvextra}
\usepackage{tikz}
\usetikzlibrary{arrows.meta}
\usepackage[colorlinks=true,linkcolor=blue!50!black,citecolor=green!35!black,urlcolor=blue!50!black]{hyperref}
\usepackage{xurl}
\hypersetup{pdftitle={Deep linear networks are mildly singular at every exact fit},
  pdfauthor={Niels Feld},pdfsubject={Exact-fit classification and disclosure of AI assistance},
  pdfkeywords={deep linear networks, local volume asymptotics, Lean, AI disclosure}}

% Glyphs that DejaVu Sans Mono lacks (used only inside Lean quotations).
\newfontfamily\leansymbolfont{DejaVu Math TeX Gyre}
\XeTeXinterchartokenstate=1
\newXeTeXintercharclass\leansymbolclass
\XeTeXcharclass"1D7D9=\leansymbolclass
\XeTeXcharclass"226B=\leansymbolclass
\XeTeXcharclass"2A05=\leansymbolclass
\XeTeXcharclass"2A06=\leansymbolclass
\XeTeXinterchartoks 0 \leansymbolclass={\begingroup\leansymbolfont}
\XeTeXinterchartoks \leansymbolclass 0={\endgroup}
\XeTeXinterchartoks 4095 \leansymbolclass={\begingroup\leansymbolfont}
\XeTeXinterchartoks \leansymbolclass 4095={\endgroup}

\DefineVerbatimEnvironment{LeanCode}{Verbatim}{fontsize=\fontsize{8.4}{10.2}\selectfont,
  frame=leftline,framerule=0.8pt,rulecolor=\color{black!45},framesep=4pt,xleftmargin=2pt,
  breaklines=true,breakanywhere=true,breaksymbolleft={},breaksymbolright={}}
\newcommand{\leansource}[5]{{\footnotesize\raggedright\color{black!60}%
  [Q#1] \lean{#2}, lines #3--#4, file SHA-256 \texttt{#5}\ldots\par}\smallskip}
\setlength{\emergencystretch}{2.5em}
\ExplSyntaxOn
\cs_new_protected:Npn \__leanid_char:n #1
  { \str_if_eq:nnTF {#1} {~} { \nobreakspace\allowbreak } { #1 \allowbreak } }
\NewDocumentCommand \lean { m }
  { \texttt{ \str_map_inline:nn {#1} { \__leanid_char:n {##1} } } }
\ExplSyntaxOff
\newcommand{\leanhash}{\texttt{\#}}

\newtheorem{theorem}{Theorem}[section]
\newtheorem{proposition}[theorem]{Proposition}
\newtheorem{lemma}[theorem]{Lemma}
\newtheorem{corollary}[theorem]{Corollary}
\newtheorem{fact}[theorem]{Fact}
\newtheorem{mainthm}{Theorem}
\renewcommand{\themainthm}{\Alph{mainthm}}
\theoremstyle{definition}
\newtheorem{definition}[theorem]{Definition}
\newtheorem{example}[theorem]{Example}
\newtheorem{question}[theorem]{Question}
\theoremstyle{remark}
\newtheorem{remark}[theorem]{Remark}

\numberwithin{equation}{section}

% ---------- notation
\newcommand{\N}{\mathbb{N}}
\newcommand{\Z}{\mathbb{Z}}
\newcommand{\Q}{\mathbb{Q}}
\newcommand{\R}{\mathbb{R}}
\newcommand{\C}{\mathbb{C}}
\newcommand{\ii}[2]{\mathopen{\text{\leansymbolfont⟦}}#1,#2\mathclose{\text{\leansymbolfont⟧}}}
\newcommand{\Rep}{\operatorname{Rep}}
\newcommand{\GL}{\operatorname{GL}}
\newcommand{\rk}{\operatorname{rk}}
\newcommand{\codim}{\operatorname{codim}}
\newcommand{\dimH}{\dim_{\mathrm{H}}}
\newcommand{\vol}{\operatorname{vol}}
\newcommand{\Sch}{\operatorname{Sch}}
\newcommand{\Ext}{\operatorname{Ext}}
\newcommand{\Hom}{\operatorname{Hom}}
\newcommand{\id}{\mathrm{I}}
\newcommand{\Fr}[1]{\lVert #1\rVert_{\mathrm F}}
\newcommand{\norm}[1]{\lVert #1\rVert}
\newcommand{\abs}[1]{\lvert #1\rvert}
\newcommand{\eps}{\varepsilon}
\newcommand{\cout}{c_{\mathrm{out}}}
\newcommand{\creg}{c_{\mathrm{reg}}}
\newcommand{\Zb}{\mathcal Z}
\newcommand{\I}{\mathbb I}
\newcommand{\lin}{\rightarrowtail}
\newcommand{\Sl}{\mathcal S}
\newcommand{\Prog}{\mathsf{Prog}}
\newcommand{\g}{\mathfrak g}
\newcommand{\Lk}{\operatorname{Lk}}
\newcommand{\Int}{\mathrm{Int}}
\newcommand{\one}{\mathbf 1}
%DIF PREAMBLE EXTENSION ADDED BY LATEXDIFF
\RequirePackage[normalem]{ulem} %DIF PREAMBLE
\RequirePackage{xcolor} %DIF PREAMBLE
\providecommand{\DIFaddtex}[1]{{\protect\color{blue}#1}} %DIF PREAMBLE
\providecommand{\DIFdeltex}[1]{{\protect\color{red}\sout{#1}}} %DIF PREAMBLE
\providecommand{\DIFaddbegin}{} %DIF PREAMBLE
\providecommand{\DIFaddend}{} %DIF PREAMBLE
\providecommand{\DIFdelbegin}{} %DIF PREAMBLE
\providecommand{\DIFdelend}{} %DIF PREAMBLE
\providecommand{\DIFmodbegin}{} %DIF PREAMBLE
\providecommand{\DIFmodend}{} %DIF PREAMBLE
\providecommand{\DIFaddFL}[1]{\DIFadd{#1}} %DIF PREAMBLE
\providecommand{\DIFdelFL}[1]{\DIFdel{#1}} %DIF PREAMBLE
\providecommand{\DIFaddbeginFL}{} %DIF PREAMBLE
\providecommand{\DIFaddendFL}{} %DIF PREAMBLE
\providecommand{\DIFdelbeginFL}{} %DIF PREAMBLE
\providecommand{\DIFdelendFL}{} %DIF PREAMBLE
%DIF HYPERREF PREAMBLE %DIF PREAMBLE
\providecommand{\DIFadd}[1]{\texorpdfstring{\DIFaddtex{#1}}{#1}} %DIF PREAMBLE
\providecommand{\DIFdel}[1]{\texorpdfstring{\DIFdeltex{#1}}{}} %DIF PREAMBLE
%DIF AMSMATHULEM PREAMBLE %DIF PREAMBLE
\makeatletter %DIF PREAMBLE
\let\sout@orig\sout %DIF PREAMBLE
\renewcommand{\sout}[1]{\ifmmode\text{\sout@orig{\ensuremath{#1}}}\else\sout@orig{#1}\fi} %DIF PREAMBLE
\makeatother %DIF PREAMBLE
%DIF LISTINGS PREAMBLE %DIF PREAMBLE
\RequirePackage{listings} %DIF PREAMBLE
\lstdefinelanguage{DIFcode}{ %DIF PREAMBLE
%DIF DIFCODE_BOLD %DIF PREAMBLE
  % unfortunately \bfseries cannot be combined with ttfamily without extra packages %DIF PREAMBLE
  % also morecomment=[il] is broken as of v1.5b of listings at least %DIF PREAMBLE
  % workaround: plot in white with tiny font %DIF PREAMBLE
  % morecomment=[il]{\%DIF\ <\ }, %DIF PREAMBLE
  moredelim=[il][\color{white}\tiny]{\%DIF\ <\ }, %DIF PREAMBLE
  moredelim=[il][\sffamily\bfseries]{\%DIF\ >\ } %DIF PREAMBLE
} %DIF PREAMBLE
\lstdefinestyle{DIFverbatimstyle}{ %DIF PREAMBLE
	language=DIFcode, %DIF PREAMBLE
	basicstyle=\ttfamily, %DIF PREAMBLE
	columns=fullflexible, %DIF PREAMBLE
	keepspaces=true %DIF PREAMBLE
} %DIF PREAMBLE
\lstnewenvironment{DIFverbatim}[1][]{\lstset{style=DIFverbatimstyle,#1}}{} %DIF PREAMBLE
\lstnewenvironment{DIFverbatim*}[1][]{\lstset{style=DIFverbatimstyle,showspaces=true,#1}}{} %DIF PREAMBLE
\lstset{extendedchars=\true,inputencoding=utf8}

%DIF END PREAMBLE EXTENSION ADDED BY LATEXDIFF

\begin{document}
\title{Deep linear networks are mildly singular\\ at every exact fit}
\author{Niels Feld\thanks{See the disclosure of AI use in Section~\ref{sec:ai}.}}
\date{29 October 2026, \DIFdelbegin \DIFdel{13}\DIFdelend \DIFaddbegin \DIFadd{14}\DIFaddend :\DIFdelbegin \DIFdel{33}\DIFdelend \DIFaddbegin \DIFadd{29}\DIFaddend :\DIFdelbegin \DIFdel{54 }\DIFdelend \DIFaddbegin \DIFadd{39 }\DIFaddend (Paris, UTC+01:00). Not refereed.}

\maketitle
\begin{center}\small Bibliographic correction compared with the version of 13:33:54 Paris.\\ Red struck-through text: deletions. Blue text: additions.\end{center}
\begin{abstract}
Let $P(A)=A_{D-1}\cdots A_0$ be the product of a chain of real matrices of arbitrary sizes, and let $\Phi$ be a matrix. At every point $A$ of the fibre $P^{-1}(\Phi)$ we determine the volume of the sublevel sets $\{B:\ \tfrac12\Fr{P(B)-\Phi}^2<\eps\}$ inside a small fixed ball around $A$: it is comparable to $\eps^{\lambda}(\log 1/\eps)^{m-1}$ as $\eps\to0$. The exponent $\lambda$ is half the local codimension of the fibre at $A$, and the integer $m\ge1$ is computed by a finite combinatorial procedure from the ranks of the partial products $A_{j-1}\cdots A_i$, equivalently from the barcode of $A$ as a representation of an equioriented quiver of type~$\mathsf A$. For a nonzero loss germ, $\lambda$ and $m$ are the local real log canonical threshold of the squared loss and its multiplicity. The equality $\lambda=\tfrac12\codim$ at every point was conjectured by Pepin Lehalleur and Rimányi, who proved it for the minimum of $\lambda$ over the fibre. The multiplicity $m$ can be strictly smaller than the number of top-dimensional strata of the fibre through~$A$. The main statements have been formalized in the Lean~4 proof assistant; we describe the formal statement and exactly which parts of the mathematical statements it covers.
\end{abstract}

\setcounter{tocdepth}{2}
\tableofcontents
\section{Introduction}\label{sec:intro}

\subsection{The question}
Fix an integer $D\ge1$ (the \emph{depth}) and integers $n_0,\dots,n_D\ge0$ (the \emph{widths}). A \emph{chain} is a tuple $A=(A_0,\dots,A_{D-1})$ of real matrices with $A_k\in\R^{n_{k+1}\times n_k}$, and its \emph{product} is
\[
P(A)=A_{D-1}A_{D-2}\cdots A_0\in\R^{n_D\times n_0}.
\]
In machine learning, $x\mapsto P(A)x$ is the function computed by a \emph{deep linear network} with weights $A$. Given a matrix $\Phi\in\R^{n_D\times n_0}$, the squared loss
\[
L_\Phi(B)=\tfrac12\Fr{P(B)-\Phi}^2
\]
uses the Frobenius norm $\Fr M=(\sum_{i,j}M_{ij}^2)^{1/2}$ and vanishes exactly on the \emph{fibre} $F_\Phi=P^{-1}(\Phi)$; its points are the \emph{exact fits} of $\Phi$. This paper describes, at every exact fit $A$, how the volume
\[
\vol\{B:\ \norm{B-A}<\delta,\ L_\Phi(B)<\eps\}
\]
behaves as $\eps\to0$, for a small fixed radius $\delta$. Here $\norm\cdot$ is the Euclidean norm in all matrix entries and $\vol$ is Lebesgue measure.

\subsection{Main results}
The precise statements are Theorems~\ref{thm:A} and~\ref{thm:B} in Section~\ref{sec:statements}. In words:

\begin{itemize}
\item \textbf{Theorem~\ref{thm:A}.} For every exact fit $A$ there are unique $\lambda\in[0,+\infty[$ and $m\in\N^*$ such that, for every sufficiently small fixed $\delta>0$, the volume above is bounded above and below by constant multiples of $\eps^{\lambda}(\log1/\eps)^{m-1}$ for all small $\eps$. Moreover
\[
\lambda=\tfrac12\,\codim_AF_\Phi,
\]
where $\codim_AF_\Phi$ is the dimension of the space of chains minus the local Hausdorff dimension of $F_\Phi$ at $A$. In particular $2\lambda\in\N$. The pair $(\lambda,m)$ does not change if Lebesgue measure is multiplied by a function that is continuous and positive at $A$.
\item \textbf{Theorem~\ref{thm:B}.} The integers $2\lambda$ and $m$ are given by an explicit finite procedure whose only inputs are the widths and the ranks of the partial products $A_{j-1}\cdots A_i$ ($0\le i\le j\le D$). The procedure is phrased in terms of the \emph{barcode} of $A$: the decomposition of $A$, viewed as a representation of the quiver $0\to1\to\cdots\to D$, into interval representations.
\end{itemize}

The precise formula for the multiplicity is Theorem~\ref{thm:mult}; the finite chains and the integers maximized along them are defined in Section~\ref{sec:statements}. The multiplicity does not, in general, count the components of maximal dimension meeting at the chosen point: Example~\ref{ex:depth4} has $m=3$ and four such components.

\subsection{Context}
In Watanabe's singular learning theory~\cite{watanabe-book,watanabe-2018}, the numbers $\lambda$ and $m$ attached to a nonnegative analytic loss at a zero, when the loss does not vanish identically near that point, are the \emph{local real log canonical threshold} and its \emph{multiplicity}; they govern, for instance, the asymptotics of Laplace integrals $\int e^{-tL}$ and of Bayesian free energies. We use only the volume definition (Definition~\ref{def:order}) and none of these consequences. That definition also covers losses identically zero near the point, for which the pair is $(0,1)$. Real-analytic and ideal-theoretic treatments of these thresholds are also given by Saito~\cite{saito} and Lin~\cite{lin}. For a sum of squares of real-analytic functions one has $\lambda\le\frac{1}{2}\codim$ at every zero, where $\codim$ is the local codimension of the zero set (Remark~\ref{rem:upper}, which uses the stratification of real-analytic sets). We informally call a zero where equality holds \emph{mildly singular}; by Theorem~\ref{thm:A}, every exact fit of a deep linear network is mildly singular in this sense.

For positive widths, the \emph{global} threshold of a deep linear network, the minimum of $\lambda$ over the whole fibre, was computed by Aoyagi and Watanabe for $D=2$~\cite{aoyagi-watanabe} and by Aoyagi for all $D$~\cite{aoyagi-2024}. Pepin Lehalleur and Rimányi~\cite{lehalleur-rimanyi} computed the codimension of the fibres and the number of their top-dimensional components using quiver representations. By comparing with Aoyagi's formula, they showed that the global threshold equals half the codimension of the fibre, for every target~$\Phi$ with nonempty fibre, and they conjectured that the equality holds at every point of the fibre~\cite[Section~8]{lehalleur-rimanyi}. Theorem~\ref{thm:A} proves this conjecture; our proof does not use Aoyagi's formula. For equal widths, the formula for $\lambda$ given by Theorems~\ref{thm:A} and~\ref{thm:B} is stated without proof in~\cite[Appendix~B, Theorem~B.5]{hitchcock-hoogland}, where it serves as ground truth for a benchmark of posterior samplers. The stratification of the fibres of $P$ over $\R$, the dimensions of the strata and their closure relations are studied in detail by Shewchuk and Bhattacharya~\cite{shewchuk-bhattacharya}, and the top-dimensional components of the zero fibre are described by Koncki and Rimányi~\cite{koncki-rimanyi}. We have not found the multiplicity formula of Theorem~\ref{thm:B} in the literature for $D\ge3$ (for $D=2$ see below); this statement reflects a bounded literature search.

For $D=2$, the pointwise problem is part~(a) of the question MAIS-O77~\cite{mais-o77} of the \emph{Math for AI Safety} program~\cite{levine-mais,mais-repo}, and also its question MAIS-O70~\cite{mais-o70} on reduced-rank regression. Candidate solutions giving the depth-two formula for $(\lambda,m)$, and the depth-two case of Theorem~\ref{thm:A}, were posted in August 2026 by Sneiderman~\cite{mais-issue-3} and by Novikov with OpenAI Codex~\cite{mais-issue-12}. A Lean-checked proof is described in~\cite{feld-o77}, and Example~\ref{ex:o77} explains how Theorem~\ref{thm:B} recovers the formula. In concurrent work motivated by low-rank adaptation, Ferry, Monod and Yano~\cite{ferry-monod-yano} study the two-factor case with the Frobenius norm replaced by a positive semidefinite quadratic form on matrices. When that form is positive definite, their loss lies between two positive multiples of $L_\Phi$, so by Lemma~\ref{lem:compare} its local invariants at exact fits are those of Theorem~\ref{thm:B} with $D=2$. When the form is degenerate, the zero set is the preimage of an affine subspace of matrices rather than a fibre of $P$; they analyse this case locally, and our methods do not cover it (see Remark~\ref{rem:covariance}). Estimators of local learning coefficients are studied in~\DIFdelbegin \DIFdel{\mbox{%DIFAUXCMD
\cite{furman-lau,lau-llc,hitchcock-hoogland}}\hskip0pt%DIFAUXCMD
. }\DIFdelend \DIFaddbegin \DIFadd{\mbox{%DIFAUXCMD
\cite{lau-llc,hitchcock-hoogland}}\hskip0pt%DIFAUXCMD
. The earlier preprint of Furman and Lau~\mbox{%DIFAUXCMD
\cite{furman-lau} }\hskip0pt%DIFAUXCMD
was expanded and merged into~\mbox{%DIFAUXCMD
\cite{lau-llc}}\hskip0pt%DIFAUXCMD
. }\DIFaddend Theorem~\ref{thm:B} determines the local pair at an exact fit from its rank table. This pair need not coincide with the global pair; an application to the output of an optimizer therefore requires the ranks at that output. Hitchcock and Hoogland~\cite[Remark~B.6]{hitchcock-hoogland} already emphasize that deep-linear examples satisfy $\lambda=\tfrac12\codim$ and consequently do not represent losses with a strictly smaller threshold, such as $x^4$ at $0$.

Related pointwise computations in other statistical models appear in Lin~\cite[Section~5]{lin}; formulas at nonsingular points of the set of realizable parameters and upper bounds at certain singular points are developed by Kurumadani~\cite{kurumadani,kurumadani-upper}. Combinatorial formulas for a threshold and its multiplicity, including a maximization over chains of subspaces, also occur for real hyperplane arrangements in Kosta and Windisch~\cite[Theorem~4.2 and Corollary~4.4]{kosta-windisch}. That work builds on the resolution constructions of De Concini and Procesi~\cite{deconcini-procesi} and the threshold calculations of Mustaţă~\cite{mustata-arrangements} and Teitler~\cite{teitler}. These are related precedents, rather than inputs to the proofs below.

\subsection{Method}
The proof uses four ingredients.
\begin{enumerate}[label=(\roman*)]
\item The orbits of $\prod_v\GL_{n_v}(\R)$ on chains, classified by barcodes. A transverse slice to the orbit of an interval representation has explicit coordinates, one for each \emph{linked pair} of intervals (Section~\ref{sec:quiver}). The interval classification is the type-$\mathsf A$ case of Gabriel's theorem~\cite{gabriel} and, in persistence-module language, is proved by Crawley-Boevey~\cite{crawley-boevey}; the normal-space calculation is the interval case of Voigt's lemma~\cite{voigt}.
\item A Schur-complement change of coordinates, which splits off the rank of $\Phi$ and reduces the problem to the zero target (Section~\ref{sec:reduction}).
\item For the zero target, a dilation of the slice coordinates under which the loss is homogeneous. Weighted polar coordinates then give a recursion for $(\lambda,m)$ over the finite poset of orbits (Section~\ref{sec:zero}).
\item A stationarity argument that identifies $2\lambda$ with a combinatorial minimum, and an explicit family of submanifolds of the fibre that identifies the same minimum with the local codimension (Sections~\ref{sec:exponent} and~\ref{sec:multiplicity}).
\end{enumerate}

\subsection{Formalization}
The public repository~\cite{feld-deep} contains a Lean~4 formalization built on the Mathlib library. It was produced with AI assistance; see Section~\ref{sec:ai}. Section~\ref{sec:formal} and Appendix~\ref{app:lean} state the formal theorem and translate every definition it depends on, so that a reader can check, without reading the formal proofs, how it implies Theorem~\ref{thm:A} and the formulas of Theorem~\ref{thm:B}; Proposition~\ref{prop:formal-implies} and the remark after Proposition~\ref{prop:prog-translation} list the steps that remain in the text. The mathematical text (Sections~\ref{sec:intro}--\ref{sec:proofs} and Appendix~\ref{app:generation}) does not depend on the formalization.

\subsection{Organization}
\begin{itemize}
\item Section~\ref{sec:notation} fixes notation.
\item Section~\ref{sec:statements} gives the definitions and states the results, with examples.
\item Sections~\ref{sec:volume}--\ref{sec:multiplicity} contain the proofs, which are assembled in Section~\ref{sec:proofs}.
\item Section~\ref{sec:formal} discusses the formalization, and Section~\ref{sec:ai} is the AI use statement.
\item Appendix~\ref{app:lean} contains the formal statement and its translation. Appendix~\ref{app:generation} proves a combinatorial lemma on barcodes.
\end{itemize}
\section{Notation and conventions}\label{sec:notation}

\subsection{Numbers, intervals and matrices}\label{sec:numbers}
$\N=\{0,1,2,\dots\}$ and $\N^*=\N\setminus\{0\}$. We use the French notation for intervals of real numbers: for real numbers $a\le b$,
\begin{gather*}
[a,b]=\{x\in\R:\ a\le x\le b\},\qquad ]a,b[\,=\{x\in\R:\ a<x<b\},\\
[a,b[\,=\{x\in\R:\ a\le x<b\},\qquad ]a,b]=\{x\in\R:\ a<x\le b\},
\end{gather*}
and similarly $]a,+\infty[$ and $[a,+\infty[$. The extended half-line is $[0,+\infty]=[0,+\infty[\,\cup\{+\infty\}$, with its usual order, addition and multiplication (and $0\cdot(+\infty)=0$). For integers $a,b$, the \emph{integer interval} is $\ii ab=\{k\in\Z:\ a\le k\le b\}$; it is empty if $a>b$. We write $\#X$ for the cardinality of a finite set $X$.

For $p,q\in\N$, $\R^{p\times q}$ is the space of real $p\times q$ matrices, identified with the linear maps $\R^q\to\R^p$; if $p=0$ or $q=0$ it is the zero space. We write $\id_p$ for the identity matrix of size $p$, $\rk M$ for the rank of a matrix $M$,
\[
\Fr M=\Bigl(\sum_{i,j}M_{ij}^2\Bigr)^{1/2}\quad\text{and}\quad\norm M_{\mathrm{op}}=\sup\{\norm{Mx}:\ x\in\R^q,\ \norm x\le1\}
\]
for its Frobenius norm and its operator norm, where $\norm x$ is the Euclidean norm of a vector. This definition gives operator norm $0$ when $q=0$, since the closed unit ball of $\R^0$ consists of its zero vector. Recall that $\norm{MN}_{\mathrm F}\le\norm M_{\mathrm{op}}\Fr N$ and $\norm{M}_{\mathrm{op}}\le\Fr M$. For a square matrix $M$, $e^M=\exp(M)=\sum_{k\ge0}M^k/k!$.

\subsection{Comparison of functions}\label{sec:comparison}
Let $\eps_1\in\,]0,+\infty]$ and let $f,g:\,]0,\eps_1[\,\to[0,+\infty]$ be two functions. For $\eps_0\in\,]0,\eps_1]$ we write $f\asymp g$ on $]0,\eps_0[$ if there are real numbers $c,C$ with $0<c\le C$ and $c\,g(\eps)\le f(\eps)\le C\,g(\eps)$ for every $\eps\in\,]0,\eps_0[$. We write $f(\eps)\asymp g(\eps)$ as $\eps\to0$ if $f\asymp g$ on $]0,\eps_0[$ for some $\eps_0\in\,]0,\eps_1]$. Throughout, $\log$ is the natural logarithm, and
\begin{equation}\label{eq:L}
L:\ ]0,1[\,\to\,]0,+\infty[,\qquad L(\eps)=\log(1/\eps)=-\log\eps .
\end{equation}
The function $L$ is decreasing, and $L(\eps)>1$ if and only if $\eps<e^{-1}$. For $\kappa>0$ and $\eps\in\,]0,\min(1,\kappa)[$ we have $L(\eps/\kappa)=L(\eps)+\log\kappa$; hence
\begin{equation}\label{eq:L-shift}
\tfrac12L(\eps)\le L(\eps/\kappa)\le2L(\eps)\qquad\text{whenever }\eps<\min(1,\kappa)\text{ and }L(\eps)\ge2\abs{\log\kappa}.
\end{equation}

\subsection{Hausdorff measures and Hausdorff dimension}\label{sec:hausdorff}
Let $(E,\norm\cdot)$ be a finite-dimensional real normed vector space. The \emph{diameter} of a nonempty subset $U\subseteq E$ is $\operatorname{diam}U=\sup\{\norm{x-y}:\ x,y\in U\}\in[0,+\infty]$, and $\operatorname{diam}\emptyset=0$. For $s\in[0,+\infty[$ put $\gamma_s(U)=(\operatorname{diam}U)^s$ if $U\neq\emptyset$, with the convention $0^0=1$, and $\gamma_s(\emptyset)=0$. For $S\subseteq E$, $s\in[0,+\infty[$ and $\eta\in\,]0,+\infty[$ put
\[
\mathcal H^s_\eta(S)=\inf\Bigl\{\sum_{i\in\N}\gamma_s(U_i):\ (U_i)_{i\in\N}\in\mathcal C_\eta(S)\Bigr\}\in[0,+\infty],
\]
where $\mathcal C_\eta(S)$ is the set of sequences $(U_i)_{i\in\N}$ of subsets of $E$ with $S\subseteq\bigcup_{i\in\N}U_i$ and $\operatorname{diam}U_i\le\eta$ for all $i$ (\emph{$\eta$-covers} of $S$).
The function $\eta\mapsto\mathcal H^s_\eta(S)$ is nonincreasing. The \emph{$s$-dimensional Hausdorff measure} of $S$ (for the norm $\norm\cdot$) is
\[
\mathcal H^s(S)=\sup_{\eta>0}\mathcal H^s_\eta(S)=\lim_{\eta\to0}\mathcal H^s_\eta(S)\in[0,+\infty],
\]
and the \emph{Hausdorff dimension} of $S$ is
\[
\dimH S=\sup\bigl\{s\in[0,+\infty[\ :\ \mathcal H^s(S)=+\infty\bigr\},\qquad\text{with }\sup\emptyset=0.
\]
We use the following standard facts~\cite[Chapter~4]{mattila},~\cite{falconer}.
\begin{enumerate}[label=(H\arabic*),leftmargin=3em,start=0]
\item\label{H:def} If $0\le s<t$ and $\mathcal H^s(S)<+\infty$, then $\mathcal H^t(S)=0$. Consequently $\dimH S=\inf\{t\in[0,+\infty[\ :\ \mathcal H^t(S)=0\}$. Moreover $\mathcal H^t(S)=0$ for every $t>\dim E$, so $\dimH S\le\dim E$.
\item\label{H:mono} If $S\subseteq S'$ then $\dimH S\le\dimH S'$. If $S=\bigcup_{i\in\N}S_i$ then $\dimH S=\sup_{i\in\N}\dimH S_i$.
\item\label{H:lip} If $E'$ is another finite-dimensional real normed vector space, $S\subseteq E$, and $f:S\to E'$ is Lipschitz, then $\dimH f(S)\le\dimH S$. Consequently $\dimH S$ does not depend on the choice of the norm on $E$ (all norms are equivalent, so the identity map is bi-Lipschitz), and $\dimH$ is invariant under bi-Lipschitz maps, for instance under the restriction of a $C^1$ diffeomorphism between open sets to a compact subset (a $C^1$ map is Lipschitz on every compact subset of its open domain).
\item\label{H:mfd} A nonempty open subset of a $C^1$ embedded submanifold of $E$ of dimension $k$ has Hausdorff dimension $k$.
\end{enumerate}

\begin{definition}[Local codimension]\label{def:codim}
Let $(E,\norm\cdot)$ be as above, $S\subseteq E$ and $a\in E$. For $\delta\in\,]0,+\infty[$ let $B_\delta(a)=\{x\in E:\ \norm{x-a}<\delta\}$ be the open ball. The \emph{local Hausdorff dimension} of $S$ at $a$ and the \emph{local codimension} of $S$ at $a$ are
\[
d_a(S)=\inf_{\delta>0}\dimH\bigl(S\cap B_\delta(a)\bigr)\in[0,\dim E],\qquad\codim_aS=\dim E-d_a(S)\in[0,\dim E].
\]
\end{definition}
The function $\delta\mapsto\dimH(S\cap B_\delta(a))$ is nondecreasing by~\ref{H:mono}, so the infimum is a limit as $\delta\to0$. If $\norm\cdot'$ is another norm on $E$, with balls $B'_\delta(a)$, there is $K\ge1$ with $B'_{\delta/K}(a)\subseteq B_\delta(a)\subseteq B'_{K\delta}(a)$ for all $\delta>0$; with~\ref{H:mono} and~\ref{H:lip} this shows that $d_a(S)$ does not depend on the norm. If $S$ is a finite union of $C^1$ embedded submanifolds $S_1,\dots,S_N$ of $E$ and $a\in\overline S$, then
\begin{equation}\label{eq:localdim-strata}
d_a(S)=\max\{\dim S_i:\ i\in\ii1N,\ a\in\overline{S_i}\}.
\end{equation}
Indeed, if $a\notin\overline{S_i}$, then $S_i\cap B_\delta(a)=\emptyset$ for small $\delta$; if $a\in\overline{S_i}$, then $S_i\cap B_\delta(a)$ is a nonempty open subset of $S_i$ for every $\delta>0$. Apply~\ref{H:mono} and~\ref{H:mfd}.

\subsection{Measures and integrals}\label{sec:measures}
Let $E$ be a finite-dimensional real vector space, with its Borel $\sigma$-algebra. A \emph{Haar measure} on $E$ is a positive multiple of the image of Lebesgue measure under a linear isomorphism $\R^{\dim E}\to E$; this does not depend on the isomorphism, and any two Haar measures differ by a positive constant factor. When $E=\R^N$ we write $\vol$ for Lebesgue measure. For a function $\varphi:E\to\R$ we write $\varphi^+=\max(\varphi,0)$. For a Borel set $S\subseteq E$, a Haar measure $\mu$ and a function $g:E\to[0,+\infty]$, the \emph{lower integral}
\[
\int^-_Sg\,d\mu=\sup\Bigl\{\int_Sg'\,d\mu:\ g':E\to[0,+\infty]\text{ Borel measurable},\ g'\le g\Bigr\}\in[0,+\infty]
\]
equals the Lebesgue integral when $g$ is measurable. It is used only so that no measurability hypothesis on the functions $\varphi$ below is needed. It is nondecreasing in $S$ and in $g$, and if $c_1,c_2\in[0,+\infty[$ satisfy $c_1\le\varphi^+(x)\le c_2$ for all $x\in S$, then
\begin{equation}\label{eq:lower-integral}
c_1\,\mu(S)\le\int^-_S\varphi^+\,d\mu\le c_2\,\mu(S).
\end{equation}
(For the lower bound take $g'$ equal to $c_1$ on $S$ and to $0$ outside $S$; this function is Borel measurable and lies below $\varphi^+$. For the upper bound note that every admissible $g'$ is at most $c_2$ on $S$.)

\subsection{Chains}\label{sec:chains}
Let $D\in\N$. The \emph{quiver} of type $\mathsf A_{D+1}$ considered here has the \emph{vertices} $v\in\ii0D$ and the \emph{arrows} $k\in\ii0{D-1}$, the arrow $k$ going from the vertex $k$ to the vertex $k+1$:
\[
0\longrightarrow1\longrightarrow\cdots\longrightarrow D .
\]
A \emph{width vector} is an element $n=(n_0,\dots,n_D)\in\N^{D+1}$; it determines $D$. For a width vector $n$ put
\[
\Rep(n)=\prod_{k=0}^{D-1}\R^{n_{k+1}\times n_k},\qquad N(n)=\dim\Rep(n)=\sum_{k=0}^{D-1}n_kn_{k+1}.
\]
The elements of $\Rep(n)$ are the \emph{chains} $A=(A_0,\dots,A_{D-1})$; the matrix $A_k$ is attached to the arrow $k$. For $D=0$, $\Rep(n)$ is the zero space. Unless stated otherwise, $\Rep(n)$ carries the Euclidean norm $\norm A=(\sum_k\Fr{A_k}^2)^{1/2}$, and $B_\delta(A)$ denotes its open ball of radius $\delta$ around $A$. For $0\le i\le j\le D$ the \emph{partial products} of a chain $A$ are
\[
P_{ij}(A)=A_{j-1}A_{j-2}\cdots A_i\in\R^{n_j\times n_i}\quad(i<j),\qquad P_{ii}(A)=\id_{n_i},
\]
and its \emph{product} is $P(A)=P_{0D}(A)$. The \emph{rank table} of $A$ is the family $R(A)=(R_{ij}(A))_{0\le i\le j\le D}$ of integers $R_{ij}(A)=\rk P_{ij}(A)$; note that $R_{vv}(A)=n_v$ for every vertex $v$.

For $\Phi\in\R^{n_D\times n_0}$ we define the \emph{fibre} $F_{n,\Phi}$ and the \emph{loss} $L_{n,\Phi}:\Rep(n)\to\R$ by
\[
F_{n,\Phi}=\{A\in\Rep(n):\ P(A)=\Phi\},\qquad L_{n,\Phi}(A)=\tfrac12\Fr{P(A)-\Phi}^2 .
\]
The \emph{zero fibre} is $Z(n)=F_{n,0}$, and we put $K_n(A)=\Fr{P(A)}^2=2L_{n,0}(A)$. When only one width vector $n$ is involved, we write $F_\Phi$ and $L_\Phi$ for $F_{n,\Phi}$ and $L_{n,\Phi}$. The points of $F_\Phi$ are the \emph{exact fits} of $\Phi$.

The group $G(n)=\prod_{v=0}^D\GL_{n_v}(\R)$ acts linearly on $\Rep(n)$ by
\[
(g\cdot A)_k=g_{k+1}A_kg_k^{-1}\qquad(g=(g_v)_{v\in\ii0D}\in G(n),\ k\in\ii0{D-1}).
\]
Then $P_{ij}(g\cdot A)=g_jP_{ij}(A)g_i^{-1}$, so the rank table is constant on the orbits of $G(n)$. The group $G(n)$ is a Lie group, open in the real vector space
\[
\g(n)=\prod_{v=0}^D\R^{n_v\times n_v},
\]
which is its Lie algebra. Its exponential map is $\exp:\g(n)\to G(n)$, $\xi=(\xi_v)_v\mapsto(e^{\xi_v})_v$. We use only that $t\mapsto\exp(t\xi)$ is a one-parameter subgroup of $G(n)$ with derivative $\xi$ at $t=0$, and that $\exp$ restricts to a diffeomorphism from a neighbourhood of $0$ in $\g(n)$ onto a neighbourhood of the identity in $G(n)$.
\section{Definitions and statements}\label{sec:statements}

\subsection{Local volume order}

\begin{definition}[Local volume order]\label{def:order}
Let $(E,\norm\cdot)$ be a finite-dimensional real normed vector space, $U\subseteq E$ an open subset, $a\in U$, and $f:U\to\R$ a continuous function. Let $\mu$ be a Haar measure on $E$, and let $\varphi:E\to\R$ be a function that is continuous at $a$ and satisfies $\varphi(a)>0$ (a \emph{prior}). For $\delta\in\,]0,+\infty[$ with $B_\delta(a)\subseteq U$ and for $\eps\in\,]0,+\infty[$ put
\[
S_{f,a}(\delta,\eps)=\{x\in B_\delta(a):\ \abs{f(x)-f(a)}<\eps\},\qquad V^{\mu,\varphi}_{f,a}(\delta,\eps)=\int^-_{S_{f,a}(\delta,\eps)}\varphi^+\,d\mu\ \in[0,+\infty].
\]
Let $\lambda\in\R$ and $m\in\N$. We say that $f$ has \emph{local volume order $(\lambda,m)$ at $a$} (for $\mu$ and $\varphi$) if $\lambda\ge0$, $m\ge1$, and there is $\delta_0\in\,]0,+\infty[$ with $B_{\delta_0}(a)\subseteq U$ such that for every $\delta\in\,]0,\delta_0[$ there are real numbers $c,C,\eps_0$ with $0<c\le C$ and $0<\eps_0<e^{-1}$ such that
\begin{equation}\label{eq:order}
c\,\eps^\lambda L(\eps)^{m-1}\ \le\ V^{\mu,\varphi}_{f,a}(\delta,\eps)\ \le\ C\,\eps^\lambda L(\eps)^{m-1}\qquad\text{for all }\eps\in\,]0,\eps_0[,
\end{equation}
where $L(\eps)=\log(1/\eps)$ is the function~\eqref{eq:L} of Section~\ref{sec:comparison}.
\end{definition}

The set $S_{f,a}(\delta,\eps)$ is open, since $f$ is continuous on the open set $U$; in particular it is a Borel set. The order of the quantifiers is: first $\delta_0$; then every fixed radius $\delta\in\,]0,\delta_0[$; then constants $c,C,\eps_0$, which may depend on $\delta$; then all $\eps\in\,]0,\eps_0[$. Nothing uniform in $\delta$ is required. If $f\ge0$ and $f(a)=0$, then $S_{f,a}(\delta,\eps)$ is the sublevel set $\{x\in B_\delta(a):\ f(x)<\eps\}$.

\begin{lemma}\label{lem:order-basic}
Keep the notation and hypotheses of Definition~\ref{def:order}.
\begin{enumerate}[label=\textup{(\alph*)}]
\item There is at most one pair $(\lambda,m)\in\R\times\N$ such that $f$ has local volume order $(\lambda,m)$ at $a$ for $\mu$ and $\varphi$.
\item Let $\mu'$ be a Haar measure on $E$, $\varphi':E\to\R$ a function continuous at $a$ with $\varphi'(a)>0$, and $\norm\cdot'$ a norm on $E$. If $f$ has local volume order $(\lambda,m)$ at $a$ for $\mu$, $\varphi$ and the balls of $\norm\cdot$, then it has local volume order $(\lambda,m)$ at $a$ for $\mu'$, $\varphi'$ and the balls of $\norm\cdot'$.
\item If $f$ is constant on a neighbourhood of $a$, then $f$ has local volume order $(0,1)$ at $a$.
\end{enumerate}
\end{lemma}
The proof is in Section~\ref{sec:volume}. By (a) and (b), when $f$ has a local volume order at $a$ for some choice of $(\mu,\varphi,\norm\cdot)$, it has the same one for every choice; we then denote it by $(\lambda_a(f),m_a(f))$.

\begin{remark}[Terminology]\label{rem:rlct}
Let $f\ge0$ be real-analytic on a neighbourhood of $a$ with $f(a)=0$. If $f$ is identically zero on a neighbourhood of $a$, its local volume order is $(0,1)$ by Lemma~\ref{lem:order-basic}(c). Otherwise, Hironaka's resolution of singularities~\cite{hironaka} gives a local volume order at $a$; in this case $\lambda_a(f)$ and $m_a(f)$ are the local real log canonical threshold of $f$ at $a$ and its multiplicity (\cite[Chapters~2 and~6]{watanabe-book}). We do not use resolution of singularities: all existence statements below are proved directly. The losses $L_\Phi$ are polynomials, so this remark applies to them; the content of the theorems below is the value of the pair, not its existence.
\end{remark}

\begin{remark}[An upper bound]\label{rem:upper}
Let $f=\sum_{j=1}^Jf_j^2$, where the functions $f_j:U\to\R$ are of class $C^1$ and vanish at $a$. Suppose that the zero set $f^{-1}(0)$ contains a $C^1$ embedded submanifold $S$ of $E$ of codimension $c_S$ with $a\in\overline S$, and that $f$ has a local volume order $(\lambda,m)$ at $a$. Then $\lambda\le c_S/2$. Indeed, identify $E$ with $\R^{\dim E}$ by a linear isomorphism. By Lemma~\ref{lem:order-basic}(b), $f$ has local volume order $(\lambda,m)$ at $a$ for Lebesgue measure $\vol$, the prior $1$ and the Euclidean norm, with some radius $\delta_0$ as in Definition~\ref{def:order}. Let $\delta\in\,]0,\delta_0[$, choose $x_0\in S\cap B_{\delta/2}(a)$, and choose a $C^1$ diffeomorphism $\Psi$ from an open neighbourhood $Q'$ of $0$ in $\R^{\dim E-c_S}\times\R^{c_S}$ onto an open neighbourhood $Q\subseteq B_\delta(a)$ of $x_0$ with $\Psi(0)=x_0$ and $\Psi^{-1}(S\cap Q)=Q'\cap(\R^{\dim E-c_S}\times\{0\})$. Let $Q_0=[-t,t]^{\dim E}\subseteq Q'$ be a cube, and write its points $(y,w)$ with $y\in\R^{\dim E-c_S}$ and $w\in\R^{c_S}$. Each $f_j\circ\Psi$ is of class $C^1$ and vanishes on $Q_0\cap\{w=0\}$, so by the mean value inequality on the segments $\{(y,\sigma w):\ \sigma\in[0,1]\}\subseteq Q_0$ there is $K>0$ with $f(\Psi(y,w))\le K\norm w^2$ for $(y,w)\in Q_0$. Let $k_0>0$ be a lower bound of $\abs{\det d\Psi}$ on $Q_0$. Since $f\ge0$ and $f(a)=0$, the set $\Psi(\{(y,w)\in Q_0:\ K\norm w^2<\eps\})$ is contained in $S_{f,a}(\delta,\eps)$, so
\begin{align*}
V^{\vol,1}_{f,a}(\delta,\eps)&\ge\vol\,\Psi\bigl(\{(y,w)\in Q_0:\ K\norm w^2<\eps\}\bigr)\\
&\ge k_0\,(2t)^{\dim E-c_S}\,\vol\{w\in[-t,t]^{c_S}:\ \norm w^2<\eps/K\},
\end{align*}
which is at least a positive constant times $\eps^{c_S/2}$ for small $\eps$. Since $\eps^\lambda L(\eps)^{m-1}/\eps^{c_S/2}\to0$ as $\eps\to0$ when $\lambda>c_S/2$, the upper bound in~\eqref{eq:order} forces $\lambda\le c_S/2$.

If $f$ is real-analytic, its zero set is, near $a$, a finite union of real-analytic embedded submanifolds of $E$~\cite{bierstone-milman}. By~\eqref{eq:localdim-strata}, one of them, $S$, satisfies $a\in\overline S$ and has codimension $c_S=\codim_af^{-1}(0)$. Hence $\lambda_a(f)\le\frac{1}{2}\codim_af^{-1}(0)$ whenever $f$ has a local volume order at $a$. For the losses $L_\Phi$, the submanifolds are given directly by Lemma~\ref{lem:fibre-strata}. Theorem~\ref{thm:A} below asserts equality in this upper bound for $L_\Phi$ at every exact fit.
\end{remark}

\subsection{The first main theorem}

\begin{mainthm}\label{thm:A}
Let $D\in\N$, let $n\in\N^{D+1}$ be a width vector, let $\Phi\in\R^{n_D\times n_0}$, and let $A\in F_\Phi$. Then the loss $L_\Phi:\Rep(n)\to\R$ has a local volume order $(\lambda_A(L_\Phi),m_A(L_\Phi))$ at $A$, and
\[
\lambda_A(L_\Phi)=\tfrac12\codim_AF_\Phi .
\]
In particular $2\lambda_A(L_\Phi)\in\N$, and the pair does not depend on the Haar measure, on the prior or on the norm used in Definition~\ref{def:order}.
\end{mainthm}

For $D=0$ the space $\Rep(n)$ is a point and $P$ is the constant map with value $\id_{n_0}$. The statement is then trivial: if $A\in F_\Phi$, then $L_\Phi$ is identically $0$, its local volume order is $(0,1)$ by Lemma~\ref{lem:order-basic}(c), and the codimension of a point in a point is $0$. The content is the case $D\ge1$.

\subsection{Barcodes}\label{sec:barcodes}
In this subsection $D\in\N^*$ is fixed.

\begin{definition}[Intervals, barcodes, slots]\label{def:barcode}
\begin{enumerate}[label=\textup{(\alph*)}]
\item An \emph{interval} (of length $D$) is a subset of $\ii0D$ of the form $\ii be$ with $0\le b\le e\le D$. The integers $b$ and $e$ are determined by the interval $I$; we write $b(I)=b$ and $e(I)=e$. The set of intervals is denoted $\Int_D$.
\item A \emph{barcode} (of length $D$) is a map $h:\Int_D\to\N$. For $I\in\Int_D$, $\one_I$ is the barcode with $\one_I(I)=1$ and $\one_I(I')=0$ for $I'\neq I$. Barcodes are added, subtracted (when the result has values in $\N$) and compared pointwise, and $h=\sum_{I\in\Int_D}h(I)\,\one_I$ for every barcode $h$. In examples we write a barcode as the formal sum $\sum_Ih(I)\,I$ of its intervals; for instance $2\ii00+\ii02$ is the barcode $2\cdot\one_{\ii00}+\one_{\ii02}$.
\item The set of \emph{slots} of a barcode $h$ is the finite set
\[
\Sl(h)=\{(I,j)\in\Int_D\times\N^*:\ j\le h(I)\}.
\]
For a slot $s=(I,j)$ we put $I_s=I$, $b_s=b(I)$ and $e_s=e(I)$, so that $I_s=\ii{b_s}{e_s}$. Thus $h(I)$ is the number of slots $s$ with $I_s=I$, and distinct slots may have the same interval.
\end{enumerate}
\end{definition}

\begin{definition}[Dimension vector and rank table of a barcode]\label{def:barcode-dims}
Let $h$ be a barcode of length $D$. For a vertex $v\in\ii0D$ put
\[
\Sl_v(h)=\{s\in\Sl(h):\ v\in I_s\},\qquad\dim_v(h)=\#\Sl_v(h)=\sum_{I\ni v}h(I).
\]
The \emph{dimension vector} of $h$ is $\dim(h)=(\dim_0(h),\dots,\dim_D(h))\in\N^{D+1}$. For $0\le i\le j\le D$ put
\[
\rho_{ij}(h)=\#\{s\in\Sl(h):\ \ii ij\subseteq I_s\}=\sum_{I\supseteq\ii ij}h(I).
\]
The family $\rho(h)=(\rho_{ij}(h))_{0\le i\le j\le D}$ is the \emph{rank table} of $h$. Note that $\rho_{vv}(h)=\dim_v(h)$ for every vertex $v$.
\end{definition}

A slot $s$ of $h$ is \emph{full-span} if $I_s=\ii0D$, an \emph{input slot} if $b_s=0$, an \emph{output slot} if $e_s=D$, and \emph{interior} if $0<b_s$ and $e_s<D$. A full-span slot is both an input and an output slot; every other slot is of exactly one of the three kinds input, output, interior. The barcode $h$ is \emph{zero-output} if $h(\ii0D)=0$, i.e.\ if it has no full-span slot. Its \emph{zero-output part} is the barcode
\[
h^\circ=h-h(\ii0D)\,\one_{\ii0D},
\]
which agrees with $h$ except that $h^\circ(\ii0D)=0$; thus $\Sl(h^\circ)$ is the set of slots of $h$ that are not full-span. Since every vertex lies in $\ii0D$,
\begin{equation}\label{eq:zero-output-part}
\dim_v(h^\circ)=\dim_v(h)-h(\ii0D),\qquad\rho_{ij}(h^\circ)=\rho_{ij}(h)-h(\ii0D)
\end{equation}
for every vertex $v$ and all $0\le i\le j\le D$.
For a width vector $n\in\N^{D+1}$ we denote by $\Zb(n)$ the set of zero-output barcodes $h$ of length $D$ with $\dim(h)=n$, i.e.\ with $\dim_v(h)=n_v$ for every vertex $v\in\ii0D$. It is a finite set, since $h(I)\le n_{b(I)}$ for every $h\in\Zb(n)$ and $I\in\Int_D$. For an integer $r\in\N$ with $r\le n_v$ for all $v$, we write $n-r=(n_0-r,\dots,n_D-r)$.

\begin{definition}[Interval chain]\label{def:interval-chain}
Let $h$ be a barcode of length $D$, put $n=\dim(h)$, and fix a total order on the finite set $\Sl(h)$. For each vertex $v$ the set $\Sl_v(h)$ has exactly $n_v$ elements (Definition~\ref{def:barcode-dims}). Let $\sigma_v:\ii1{n_v}\to\Sl_v(h)$ be the unique increasing bijection, and for $s\in\Sl_v(h)$ let $\epsilon^v_s$ be the $\sigma_v^{-1}(s)$-th vector of the standard basis of $\R^{n_v}$. Thus $(\epsilon^v_s)_{s\in\Sl_v(h)}$ is the standard basis of $\R^{n_v}$, indexed by the slots present at $v$. The \emph{interval chain} $\I(h)\in\Rep(n)$ is defined, for each arrow $k\in\ii0{D-1}$, by
\[
\I(h)_k\,\epsilon^k_s=\begin{cases}\epsilon^{k+1}_s&\text{if }k+1\in I_s,\\0&\text{if }k+1\notin I_s,\end{cases}\qquad s\in\Sl_k(h).
\]
\end{definition}
So each matrix $\I(h)_k$ is a partial permutation matrix: a slot is carried from one vertex to the next as long as its interval lasts. Changing the total order on $\Sl(h)$ replaces $\I(h)$ by $g\cdot\I(h)$, where $g\in G(n)$ consists of permutation matrices; whenever $\I(h)$ is used below, a total order is fixed, and no statement depends on it.

\begin{fact}[Interval decomposition]\label{fact:gabriel}
\begin{enumerate}[label=\textup{(\alph*)}]
\item For every width vector $n\in\N^{D+1}$ and every chain $A\in\Rep(n)$ there are a barcode $h$ of length $D$ with $\dim(h)=n$ and an element $g\in G(n)$ with $g\cdot A=\I(h)$.
\item For every barcode $h$ of length $D$, the rank table of the chain $\I(h)\in\Rep(\dim(h))$ (Section~\ref{sec:chains}) is the rank table of the barcode $h$: $R(\I(h))=\rho(h)$.
\item For every barcode $h$ of length $D$ and all $0\le b\le e\le D$,
\[
h(\ii be)=\rho_{be}(h)-\rho_{b-1,e}(h)-\rho_{b,e+1}(h)+\rho_{b-1,e+1}(h),
\]
where $\rho_{ij}(h)$ is read as $0$ when $i<0$ or $j>D$. In particular the rank table $\rho(h)$ determines $h$.
\end{enumerate}
\end{fact}
Part (a) is Gabriel's theorem for the quiver $0\to1\to\cdots\to D$ of type $\mathsf A_{D+1}$~\cite{gabriel}; for this quiver it is also the structure theorem for persistence modules~\cite{crawley-boevey}. Part (b) is a direct computation: for $s\in\Sl_i(h)$, $P_{ij}(\I(h))$ maps $\epsilon^i_s$ to $\epsilon^j_s$ if $\ii ij\subseteq I_s$ and to $0$ otherwise, so its rank is $\rho_{ij}(h)$. For (c), $\rho_{be}(h)-\rho_{b-1,e}(h)$ is the number of slots $s$ with $b_s=b$ and $e_s\ge e$, and $\rho_{b,e+1}(h)-\rho_{b-1,e+1}(h)$ is the number of slots with $b_s=b$ and $e_s\ge e+1$.

\begin{lemma}\label{lem:barcode-of-chain}
Let $n\in\N^{D+1}$ and $A\in\Rep(n)$.
\begin{enumerate}[label=\textup{(\roman*)}]
\item There is exactly one barcode $h$ of length $D$ with $\rho(h)=R(A)$. It satisfies $\dim(h)=n$ and $h(\ii0D)=\rk P(A)$, and there is $g\in G(n)$ with $g\cdot A=\I(h)$.
\item Two chains $A,A'\in\Rep(n)$ lie in the same $G(n)$-orbit if and only if $R(A)=R(A')$.
\end{enumerate}
\end{lemma}
\begin{proof}
(i) By Fact~\ref{fact:gabriel}(a) there are a barcode $h$ of length $D$ with $\dim(h)=n$ and $g\in G(n)$ with $g\cdot A=\I(h)$. Since the rank table is constant on orbits, $R(A)=R(\I(h))=\rho(h)$ by Fact~\ref{fact:gabriel}(b). A barcode $h'$ with $\rho(h')=R(A)=\rho(h)$ equals $h$ by Fact~\ref{fact:gabriel}(c). Finally $h(\ii0D)=\rho_{0D}(h)=R_{0D}(A)=\rk P(A)$, because $\ii0D$ is the only interval containing $\ii0D$.

(ii) The rank table is constant on orbits. Conversely, if $R(A)=R(A')$, then (i) gives the same barcode $h$ for $A$ and $A'$, and both lie in the orbit of $\I(h)$.
\end{proof}

\begin{definition}[Barcode of a chain; residual barcode]\label{def:barcode-of-chain}
Let $n\in\N^{D+1}$ and $A\in\Rep(n)$.
\begin{enumerate}[label=\textup{(\alph*)}]
\item The \emph{barcode of $A$} is the barcode $\beta(A)$ of Lemma~\ref{lem:barcode-of-chain}(i): the unique barcode of length $D$ with $\rho(\beta(A))=R(A)$.
\item The \emph{residual barcode} of $A$ is the zero-output part $h_A=\beta(A)^\circ$ of $\beta(A)$.
\end{enumerate}
\end{definition}
Put $r=\rk P(A)$. By Lemma~\ref{lem:barcode-of-chain}(i), $\beta(A)$ has exactly $r$ full-span slots, and each of them contains every vertex, so $r\le\dim_v(\beta(A))=n_v$ for every vertex $v$. By~\eqref{eq:zero-output-part}, $h_A\in\Zb(n-r)$. By Lemma~\ref{lem:barcode-of-chain}, two chains $A,A'\in\Rep(n)$ lie in the same $G(n)$-orbit if and only if $\beta(A)=\beta(A')$, and $A$ lies in the orbit of $\I(\beta(A))$.

\begin{definition}[Links]\label{def:links}
Let $h$ be a barcode of length $D$. An ordered pair $(s,s')\in\Sl(h)\times\Sl(h)$ is \emph{linked}, written $s\lin s'$, if
\[
b_s<b_{s'}\le e_s+1\le e_{s'} .
\]
We denote by $\Lk(h)\subseteq\Sl(h)\times\Sl(h)$ the set of linked pairs. For a slot $s$ let $\mathrm{in}(s)=\#\{s'\in\Sl(h):\ s'\lin s\}$ and $\mathrm{out}(s)=\#\{s'\in\Sl(h):\ s\lin s'\}$. The \emph{output count} of $h$ is
\[
\cout(h)=\#\{(s,s')\in\Lk(h):\ s'\text{ is an output slot}\}.
\]
The barcode $h$ is \emph{balanced} if $\mathrm{in}(s)=\mathrm{out}(s)$ for every interior slot $s$.
\end{definition}
In words: $s\lin s'$ means that $I_{s'}$ begins strictly after $I_s$ begins and at most one vertex after $I_s$ ends, and that $I_{s'}$ ends strictly after $I_s$ ends; equivalently, $I_s\cup I_{s'}$ is an interval, and $I_s$ begins and ends strictly before $I_{s'}$. No pair involving a full-span slot $s_0$ is linked: $s_0\lin s'$ would require $e_{s'}>D$, and $s\lin s_0$ would require $b_s<0$. If $s\lin s'$, then $s$ is not an output slot (since $e_s<e_{s'}\le D$) and $s'$ is not an input slot (since $b_{s'}>b_s\ge0$). By Lemma~\ref{lem:ext} below, $\#\Lk(h)$ is the codimension of the orbit of $\I(h)$.

\begin{definition}[Order and combinatorial invariants]\label{def:invariants}
For barcodes $h,k$ of length $D$ we write $h\preceq k$ if $\dim(h)=\dim(k)$ and $\rho_{ij}(h)\le\rho_{ij}(k)$ for all $0\le i\le j\le D$, and $h\prec k$ if moreover $h\neq k$. Let $n\in\N^{D+1}$ and $h\in\Zb(n)$.
\begin{enumerate}[label=\textup{(\roman*)}]
\item $C(h)=\min\{\cout(k):\ k\in\Zb(n),\ h\preceq k,\ k\text{ balanced}\}$. (Lemma~\ref{lem:nonempty} shows that this set is nonempty.)
\item $h$ is \emph{critical} if it is balanced and $\cout(h)=C(h)$.
\item $F(h)$ is the maximum of $\#\{i\in\ii0{\ell-1}:\ h_i\text{ is critical}\}$ over all $\ell\in\N$ and all finite sequences $h=h_0\prec h_1\prec\cdots\prec h_\ell$ of elements of $\Zb(n)$ with $C(h_i)=C(h)$ for all $i\in\ii0\ell$.
\end{enumerate}
\end{definition}

The relation $\preceq$ is a partial order on the barcodes of length $D$; it is antisymmetric by Fact~\ref{fact:gabriel}(c). Because $\Zb(n)$ is finite and $\prec$ is a strict order, the maximum in (iii) is taken over a finite nonempty set (the sequence with $\ell=0$ is allowed), and $F(h)\le\#\Zb(n)$. The last term $h_\ell$ of a sequence is never counted.

\begin{lemma}\label{lem:nonempty}
Let $n\in\N^{D+1}$ and $h\in\Zb(n)$, and let $k$ be a maximal element, for the order $\preceq$, of the finite set $\{k'\in\Zb(n):\ h\preceq k'\}$. Then every linked pair of $k$ goes from an input slot to an output slot, and $k$ is balanced. In particular the set minimized in the definition of $C(h)$ is nonempty.
\end{lemma}
\begin{proof}
The set contains $h$, so it has maximal elements. Suppose that $k$ has a linked pair $\pi=(s,s')$ such that $s$ is not an input slot or $s'$ is not an output slot, and let $k^\pi$ be the barcode of Lemma~\ref{lem:switch}; its definition, the equality $\dim(k^\pi)=n$ and the formula for its rank table proved there are purely combinatorial. The intervals of $k^\pi$ are those of $k$, except that $I_s,I_{s'}$ are replaced by $\ii{b_s}{e_{s'}}$ and (if $b_{s'}\le e_s$) $\ii{b_{s'}}{e_s}$. The first equals $\ii0D$ only if $b_s=0$ and $e_{s'}=D$, which is excluded; the second does not contain $0$, since $b_{s'}>b_s\ge0$. So $k^\pi\in\Zb(n)$. By the formula of Lemma~\ref{lem:switch}, $\rho_{ij}(k^\pi)\ge\rho_{ij}(k)$ for all $i\le j$, with strict inequality for $(i,j)=(b_s,e_{s'})$. Hence $h\preceq k\prec k^\pi$, which contradicts the maximality of $k$. So every linked pair of $k$ goes from an input slot to an output slot. An interior slot is neither, so it belongs to no linked pair, and $\mathrm{in}(s)=\mathrm{out}(s)=0$ for every interior slot $s$: $k$ is balanced. It belongs to the set minimized in $C(h)$.
\end{proof}

\subsection{The second main theorem}
Let $D\in\N^*$ and $n\in\N^{D+1}$. For $r\in\N$ with $r\le\min(n_0,n_D)$ put
\[
\creg(n,r)=r\,(n_0+n_D-r);
\]
it is the dimension of the manifold of $n_D\times n_0$ real matrices of rank $r$.

\begin{mainthm}\label{thm:B}
Let $D\in\N^*$, let $n\in\N^{D+1}$ be a width vector, let $\Phi\in\R^{n_D\times n_0}$, let $A\in F_\Phi$, and put $r=\rk\Phi$. Let $h_A\in\Zb(n-r)$ be the residual barcode of $A$ (Definition~\ref{def:barcode-of-chain}). The set minimized in the definition of $C(h_A)$ is nonempty, and
\[
\codim_AF_\Phi=2\lambda_A(L_\Phi)=\creg(n,r)+C(h_A),\qquad m_A(L_\Phi)=1+F(h_A).
\]
Moreover $C(h_A)=\min\{\#\Lk(k):\ k\in\Zb(n-r),\ h_A\preceq k\}$.
\end{mainthm}

Since $r=R_{0D}(A)$ and $h_A=\beta(A)^\circ$ with $\rho(\beta(A))=R(A)$, the right-hand sides depend only on $n$ and on the rank table $R(A)$. The rank table is computable from $A$, and the sets involved are finite, so Theorem~\ref{thm:B} gives an algorithm. For a width vector $n$ and the rank table $R=R(A)$ of a chain $A\in\Rep(n)$ we put
\[
\Prog(n,R)=\{(\creg(n,r)+C(h_A),\,1+F(h_A))\}\ \text{ if }D\ge1,\qquad\Prog(n,R)=\{(0,1)\}\ \text{ if }D=0,
\]
a one-element set, which depends only on $(n,R)$; the value for $D=0$ agrees with the discussion after Theorem~\ref{thm:A}. Appendix~\ref{app:lean} compares $\Prog$ with the formal definition.

The geometric meaning of $F$ is given in Theorem~\ref{thm:mult}: by Proposition~\ref{prop:closure}, the relation $h\preceq k$ means that the orbit of $\I(h)$ lies in the closure of the orbit of $\I(k)$. Thus $F(h_A)$ counts critical strata along chains of strata of the residual zero fibre that degenerate to $\I(h_A)$ and along which the local codimension stays constant.

\subsection{Examples}\label{sec:examples}
Barcodes are written as formal sums of intervals (Definition~\ref{def:barcode}(b)).

\begin{example}[One matrix]\label{ex:one}
Let $D=1$, $n=(n_0,n_1)$, and $A=(A_0)$ with $A_0=\Phi$ of rank $r$. Then $\beta(A)=r\,\ii01+(n_0-r)\ii00+(n_1-r)\ii11$ and $h_A=(n_0-r)\ii00+(n_1-r)\ii11$. The linked pairs of $h_A$ are the pairs $(s,s')$ with $I_s=\ii00$ and $I_{s'}=\ii11$, so $\cout(h_A)=\#\Lk(h_A)=(n_0-r)(n_1-r)$. There is no interior slot, so $h_A$ is balanced. A barcode of length $1$ without full-span slot is determined by its dimension vector, so $\Zb(n-r)=\{h_A\}$. Hence $C(h_A)=(n_0-r)(n_1-r)$ and $F(h_A)=0$, and Theorem~\ref{thm:B} gives $2\lambda=r(n_0+n_1-r)+(n_0-r)(n_1-r)=n_0n_1$ and $m=1$. This is consistent with the fact that $P$ is the identity map of $\Rep(n)=\R^{n_1\times n_0}$, so that $F_\Phi=\{\Phi\}$ is a point.
\end{example}

\begin{example}[Scalar chains]\label{ex:scalar}
Let $n_v=1$ for all $v\in\ii0D$. A barcode with dimension vector $(1,\dots,1)$ is a partition of $\ii0D$ into consecutive intervals, and $h\preceq k$ means that every interval of $h$ is contained in an interval of $k$. If $A_k\neq0$ for all $k$, then $r=1$, $\beta(A)=\ii0D$, $h_A=0$ is the empty barcode, and Theorem~\ref{thm:B} gives $(\lambda,m)=(\tfrac12,1)$. If exactly $p\ge1$ of the matrices vanish, say $A_{k_1}=\cdots=A_{k_p}=0$ with $k_1<\dots<k_p$, then $\Phi=0$, $r=0$ and
\[
h_A=\ii0{k_1}+\ii{k_1+1}{k_2}+\cdots+\ii{k_p+1}D .
\]
Consecutive intervals are linked and no others, so $\Lk(h_A)$ consists of $p$ pairs forming a path. Every interior slot has one incoming and one outgoing linked pair, and $\cout(h_A)=1$. The barcodes $k\in\Zb(n)$ with $h_A\preceq k$ are obtained by merging consecutive intervals of $h_A$ while keeping at least two intervals. They all have the same shape, so they are all balanced with $\cout=1$, all critical, and $C=1$ on all of them. A longest sequence $h_A=h_0\prec\cdots\prec h_\ell$ merges one pair at each step and has $\ell=p-1$, so $F(h_A)=p-1$ and $(\lambda,m)=(\tfrac12,p)$. Directly: near $A$, the loss $L_0$ is comparable to the square of a product of $p$ independent coordinates.
\end{example}

\begin{example}[Depth two; the O77 table]\label{ex:o77}
Let $D=2$ and $n=(I,H,O)$; in this example only, the letters $I,H,O$ denote these widths (input, hidden, output), and $a,b,p,q,s,\eta,c(t)$ denote the quantities defined below. Let $a=\rk A_0$, $b=\rk A_1$ and $r=\rk\Phi$. For $I=H=O=2$ the program gives:
\begin{center}\small
\begin{tabular}{@{}ccclc@{}}
\toprule
$a$&$b$&$r$&$\beta(A)$&$(2\lambda,m)$\\\midrule
$0$&$0$&$0$&$2\ii00+2\ii11+2\ii22$&$(3,1)$\\
$1$&$1$&$0$&$\ii00+\ii01+\ii12+\ii22$&$(3,1)$\\
$2$&$0$&$0$&$2\ii01+2\ii22$&$(4,1)$\\
$1$&$1$&$1$&$\ii00+\ii02+\ii11+\ii22$&$(4,2)$\\
$2$&$1$&$1$&$\ii01+\ii02+\ii22$&$(4,1)$\\
$2$&$2$&$2$&$2\ii02$&$(4,1)$\\\bottomrule
\end{tabular}
\end{center}
In depth two, a closed formula for the pair was proposed by Sneiderman~\cite{mais-issue-3} (for the problem MAIS-O70) and by Novikov~\cite{mais-issue-12} (for MAIS-O77), and proved in Lean in~\cite{feld-o77}. For $I,O>0$ and $r<H$ it reads
\[
2\lambda=s+\min_{0\le t\le\min(\eta,q)}c(t),\qquad m=\#\{t\in\ii0{\min(\eta,q)}:\ c(t)=\min c\},
\]
where $p=O-b$, $q=I-a$, $\eta=H+r-a-b$, $s=Oa+b(I-a)$ and $c(t)=(\eta-t)(q-t)+pt$. Since a local volume order is unique (Lemma~\ref{lem:order-basic}), Theorem~\ref{thm:B} and this formula agree under these hypotheses. The last row of the table, with $r=H$, is outside them: there $P$ is a submersion at $A$, so $F_\Phi$ is a submanifold of codimension $O\cdot I=4$ near $A$.
\end{example}

\begin{example}[A squared logarithm in depth three]\label{ex:depth3}
Let $n=(2,2,2,2)$ and $\beta(A)=\ii00+\ii03+\ii11+\ii22+\ii33$, so that $r=1$. The residual barcode is the barcode $h_A=\ii00+\ii11+\ii22+\ii33\in\Zb(1,1,1,1)$ of Example~\ref{ex:scalar} with $p=3$. Hence $2\lambda=\creg(n,1)+1=3+1=4$ and $m=3$.
\end{example}

\begin{example}[Multiplicity smaller than the number of top strata]\label{ex:depth4}
Let $n=(2,2,2,2,2)$, $\Phi=0$, and let $h\in\Zb(n)$ be the barcode with six slots $s_1,\dots,s_6$:
\[
h=\underbrace{\ii00}_{s_1}+\underbrace{\ii02}_{s_2}+\underbrace{\ii11}_{s_3}+\underbrace{\ii24}_{s_4}+\underbrace{\ii33}_{s_5}+\underbrace{\ii44}_{s_6}.
\]
The linked pairs are $s_1\lin s_3$, $s_2\lin s_4$, $s_2\lin s_5$, $s_3\lin s_4$ and $s_5\lin s_6$. Give them the coordinates $x_1,x_2,x_3,x_4,x_5$ respectively. With the bases ordered $(s_1,s_2)$, $(s_2,s_3)$, $(s_2,s_4)$, $(s_4,s_5)$, $(s_4,s_6)$ at the vertices $0,\dots,4$, the chains $\I(h)+\sum_{\pi}x_\pi J_\pi$ of the slice chart of Lemma~\ref{lem:slice} are
\begin{gather*}
A_0=\begin{pmatrix}0&1\\x_1&0\end{pmatrix},\quad A_1=\begin{pmatrix}1&0\\x_2&x_4\end{pmatrix},\quad A_2=\begin{pmatrix}0&1\\x_3&0\end{pmatrix},\quad A_3=\begin{pmatrix}1&0\\0&x_5\end{pmatrix},\\
\text{so that}\quad A_3A_2A_1A_0=\begin{pmatrix}x_1x_4&x_2\\0&x_3x_5\end{pmatrix}.
\end{gather*}
By Section~\ref{sec:zero}, near $\I(h)$ the loss is comparable to $x_2^2+(x_1x_4)^2+(x_3x_5)^2$ in suitable coordinates, up to spectator variables. A direct computation gives volume $\asymp\eps^{3/2}L(\eps)^2$: the estimate $\vol\{(x,x')\in[-1,1]^2:\ \abs{xx'}<\eta\}\asymp\eta L(\eta)$, applied to $x_1x_4$ and to $x_3x_5$, gives the upper bound, and requiring each of the three terms to be $<\eps/3$ gives the lower bound. So $(\lambda,m)=(\tfrac32,3)$, as Theorem~\ref{thm:B} predicts.

The zero set of $(x_2,x_1x_4,x_3x_5)$ has four irreducible components $\{x_2=0\}\cap\{x_1=0\text{ or }x_4=0\}\cap\{x_3=0\text{ or }x_5=0\}$, all of codimension $3$. They are the closures of the four orbits of the barcodes $\ii00+\ii02+\ii14+\ii34$, $\ii00+\ii03+\ii14+\ii44$, $\ii01+\ii02+\ii24+\ii34$ and $\ii01+\ii03+\ii24+\ii44$. Each of these barcodes has exactly three linked pairs, all from an input slot to an output slot. So four top-dimensional strata pass through $\I(h)$, while $m=3$. The logarithmic exponent $m-1=2$ counts the two independent ``product crossings'' $x_1x_4$ and $x_3x_5$, whereas the number of components counts the $2\times2$ ways of resolving them.
\end{example}

\begin{corollary}[Global minimum]\label{cor:global}
Let $D\in\N$, $n\in\N^{D+1}$ and $\Phi\in\R^{n_D\times n_0}$ with $F_\Phi\neq\emptyset$. Then
\[
\min_{A\in F_\Phi}\lambda_A(L_\Phi)=\tfrac12\bigl(N(n)-\dimH F_\Phi\bigr).
\]
\end{corollary}
\begin{proof}
For $D=0$ both sides are $0$. For $D\ge1$, $F_\Phi$ is a finite union of embedded submanifolds (Lemma~\ref{lem:fibre-strata}), so by~\ref{H:mono} and~\eqref{eq:localdim-strata} $\dimH F_\Phi=\max_{A\in F_\Phi}d_A(F_\Phi)$. Apply Theorem~\ref{thm:A} at each point.
\end{proof}
Pepin Lehalleur and Rimányi~\cite{lehalleur-rimanyi} proved this global equality for the real log canonical threshold defined through zeta functions, by comparing their codimension formula with Aoyagi's computation~\cite{aoyagi-2024}.
\section{Local volume orders}\label{sec:volume}

Throughout this section $(E,\norm\cdot)$ is a finite-dimensional real normed vector space, $U\subseteq E$ is an open subset, $a\in U$, and $f:U\to\R$ is a continuous function. A \emph{prior} is a function $\varphi:E\to\R$ that is continuous at $a$ with $\varphi(a)>0$. We put $\delta_U=\sup\{\delta>0:\ B_\delta(a)\subseteq U\}\in\,]0,+\infty]$. It is convenient to allow more general domains of integration than balls.

\begin{definition}[Window systems]\label{def:windows}
A \emph{window system at $a$ in $U$} is a family $\mathcal W=(W_\delta)_{\delta\in\,]0,\delta_{\mathcal W}[}$ of Borel subsets of $U$, where $\delta_{\mathcal W}\in\,]0,+\infty]$, which \emph{shrinks to $a$}: for every $\eta>0$ there is $\theta>0$ such that $W_\delta\subseteq B_\eta(a)$ for every $\delta\in\,]0,\min(\theta,\delta_{\mathcal W})[$. For a Haar measure $\mu$ on $E$, a prior $\varphi$, $\delta\in\,]0,\delta_{\mathcal W}[$ and $\eps\in\,]0,+\infty[$ we put
\[
V^{\mathcal W,\mu,\varphi}_{f,a}(\delta,\eps)=\int^-_{\{x\in W_\delta:\ \abs{f(x)-f(a)}<\eps\}}\varphi^+\,d\mu\ \in[0,+\infty].
\]
We say that $f$ has \emph{local volume order $(\lambda,m)$ at $a$ with respect to $\mathcal W$} (for $\mu$ and $\varphi$) if the condition of Definition~\ref{def:order} holds with $V^{\mu,\varphi}_{f,a}$ replaced by $V^{\mathcal W,\mu,\varphi}_{f,a}$ and the condition $B_{\delta_0}(a)\subseteq U$ replaced by $\delta_0\le\delta_{\mathcal W}$. Two window systems $\mathcal W$ and $\mathcal W'$ at $a$ in $U$ are \emph{equivalent} if for every $\eta>0$ there is $\theta>0$ such that for every $\delta\in\,]0,\min(\theta,\delta_{\mathcal W})[$ there are $\delta_1,\delta_2\in\,]0,\min(\eta,\delta_{\mathcal W'})[$ with $W'_{\delta_1}\subseteq W_\delta\subseteq W'_{\delta_2}$, and if the same holds with $\mathcal W$ and $\mathcal W'$ exchanged.
\end{definition}

The \emph{ball system} $\mathcal B=(B_\delta(a))_{\delta\in\,]0,\delta_U[}$ is a window system at $a$ in $U$, and $V^{\mathcal B,\mu,\varphi}_{f,a}=V^{\mu,\varphi}_{f,a}$, so local volume orders with respect to $\mathcal B$ are those of Definition~\ref{def:order}. The set $\{x\in W_\delta:\ \abs{f(x)-f(a)}<\eps\}$ is the intersection of the Borel set $W_\delta$ with an open set, hence a Borel set. The function $V^{\mathcal W,\mu,\varphi}_{f,a}(\delta,\eps)$ is nondecreasing in $\eps$, and nondecreasing in the window: if $\mathcal W'$ is another window system and $W'_{\delta'}\subseteq W_\delta$, then $V^{\mathcal W',\mu,\varphi}_{f,a}(\delta',\eps)\le V^{\mathcal W,\mu,\varphi}_{f,a}(\delta,\eps)$ for all $\eps$.

\begin{lemma}[Sandwich]\label{lem:sandwich}
Let $\mathcal W$ and $\mathcal W'$ be equivalent window systems at $a$ in $U$, let $\mu$ be a Haar measure on $E$ and $\varphi$ a prior. For all $(\lambda,m)\in\R\times\N$, $f$ has local volume order $(\lambda,m)$ at $a$ with respect to $\mathcal W$ if and only if it does with respect to $\mathcal W'$ (both for $\mu$ and $\varphi$).
\end{lemma}
\begin{proof}
In this proof write $V=V^{\mathcal W,\mu,\varphi}_{f,a}$ and $V'=V^{\mathcal W',\mu,\varphi}_{f,a}$. Suppose that $f$ has local volume order $(\lambda,m)$ with respect to $\mathcal W'$, with radius $\delta_0'$. Apply the definition of equivalence with $\eta=\delta_0'$, and let $\theta$ be the corresponding number; put $\delta_0=\min(\theta,\delta_{\mathcal W})$. Let $\delta\in\,]0,\delta_0[$, and choose $\delta_1,\delta_2\in\,]0,\delta'_0[$ with $W'_{\delta_1}\subseteq W_\delta\subseteq W'_{\delta_2}$. Then $V'(\delta_1,\eps)\le V(\delta,\eps)\le V'(\delta_2,\eps)$ for all $\eps>0$. Let $(c_i,C_i,\eps_i)$ be the constants of~\eqref{eq:order} for $V'$ at the radius $\delta_i$ ($i=1,2$). Then~\eqref{eq:order} holds for $V$ at the radius $\delta$ with the constants $c_1$, $\max(C_2,c_1)$ and $\min(\eps_1,\eps_2)$. The converse is the same argument with $\mathcal W$ and $\mathcal W'$ exchanged.
\end{proof}

\begin{example}\label{ex:windows}
The following window systems at $a$ in $U$ are equivalent to the ball system $\mathcal B$ of $\norm\cdot$.
\begin{enumerate}[label=\textup{(\alph*)}]
\item The ball system of another norm $\norm\cdot'$ on $E$. Indeed, there is $K\ge1$ with $B'_{\delta/K}(a)\subseteq B_\delta(a)\subseteq B'_{K\delta}(a)$ for all $\delta>0$, where $B'$ denotes the balls of $\norm\cdot'$.
\item When $E=E_1\times E_2$ is the product of two finite-dimensional real normed vector spaces $(E_1,\norm\cdot_1)$ and $(E_2,\norm\cdot_2)$ and $a=(a_1,a_2)$: the family of products $B^{E_1}_\delta(a_1)\times B^{E_2}_\delta(a_2)$ (for small $\delta$), since these are the balls of the norm $\max(\norm{x_1}_1,\norm{x_2}_2)$; then apply (a).
\item Let $(E',\norm\cdot')$ be a finite-dimensional real normed vector space, $U'\subseteq E'$ open, $a'\in U'$, and $\Psi:U'\to U$ a $C^1$ diffeomorphism onto $U$ with $\Psi(a')=a$. Choose $r_0>0$ such that the closed ball $\overline{B'}_{r_0}(a')$ is contained in $U'$. Then the family $(\Psi(B'_\delta(a')))_{\delta\in\,]0,r_0[}$ is a window system at $a$ in $U$ equivalent to $\mathcal B$. Indeed, its members are open, since $\Psi$ is a homeomorphism onto the open set $U$. The map $\Psi$ is Lipschitz on the compact convex set $\overline{B'}_{r_0}(a')$, by the mean value inequality, and $\Psi^{-1}$ is Lipschitz on a closed ball $\overline B_{r_1}(a)\subseteq U$; let $K\ge1$ be a common Lipschitz constant. Then $B_{\delta/K}(a)\subseteq\Psi(B'_\delta(a'))\subseteq B_{K\delta}(a)$ for every $\delta\in\,]0,\min(r_0,r_1)[$.
\end{enumerate}
\end{example}

\begin{lemma}[Uniqueness]\label{lem:unique}
Let $\mathcal W$ be a window system at $a$ in $U$, $\mu$ a Haar measure on $E$ and $\varphi$ a prior. There is at most one pair $(\lambda,m)\in\R\times\N$ such that $f$ has local volume order $(\lambda,m)$ at $a$ with respect to $\mathcal W$, for $\mu$ and $\varphi$.
\end{lemma}
\begin{proof}
Write $V=V^{\mathcal W,\mu,\varphi}_{f,a}$. Suppose $(\lambda,m)$ and $(\lambda',m')$ are both local volume orders, and fix $\delta$ smaller than both radii $\delta_0$. Let $(c,C,\eps_0)$ and $(c',C',\eps'_0)$ be the corresponding constants. For $\eps<\min(\eps_0,\eps'_0)$, writing $L=L(\eps)>1$,
\[
c\,\eps^\lambda L^{m-1}\le V(\delta,\eps)\le C'\eps^{\lambda'}L^{m'-1}.
\]
Taking logarithms, $(\lambda'-\lambda)\,L\le\log(C'/c)+(m'-m)\log L$. Dividing by $L$ and letting $\eps\to0$, so that $L\to+\infty$, gives $\lambda'\le\lambda$; by symmetry $\lambda=\lambda'$. Then $c\,L^{m-1}\le C'L^{m'-1}$ for all small $\eps$, so $m\le m'$, and by symmetry $m=m'$.
\end{proof}

\begin{lemma}[Measure and prior]\label{lem:prior}
Let $\mathcal W$ be a window system at $a$ in $U$, let $\mu,\mu'$ be Haar measures on $E$, and let $\varphi,\varphi'$ be priors. For all $(\lambda,m)\in\R\times\N$, $f$ has local volume order $(\lambda,m)$ at $a$ with respect to $\mathcal W$ for $(\mu,\varphi)$ if and only if it does for $(\mu',\varphi')$.
\end{lemma}
\begin{proof}
Write $V=V^{\mathcal W,\mu,\varphi}_{f,a}$ and $V'=V^{\mathcal W,\mu',\varphi'}_{f,a}$.

\emph{Step 1: the measures.} There is a real number $\kappa>0$ with $\mu'=\kappa\mu$ (Section~\ref{sec:measures}).

\emph{Step 2: the priors near $a$.} Since $\varphi$ is continuous at $a$ and $\varphi(a)>0$, there is $\eta>0$ with $\frac{1}{2}\varphi(a)\le\varphi(x)\le2\varphi(a)$ for all $x\in B_\eta(a)$; in particular $\varphi^+=\varphi$ on $B_\eta(a)$. Likewise there is $\eta'>0$ with $\frac{1}{2}\varphi'(a)\le\varphi'(x)\le2\varphi'(a)$ for all $x\in B_{\eta'}(a)$. Put $\eta''=\min(\eta,\eta')$.

\emph{Step 3: comparison of the volume functions.} Since $\mathcal W$ shrinks to $a$, there is $\theta>0$ with $W_\delta\subseteq B_{\eta''}(a)$ for all $\delta\in\,]0,\min(\theta,\delta_{\mathcal W})[$. For such $\delta$ and every $\eps>0$, the set $S=\{x\in W_\delta:\ \abs{f(x)-f(a)}<\eps\}$ is contained in $B_{\eta''}(a)$, so by~\eqref{eq:lower-integral}
\[
\tfrac12\varphi(a)\,\mu(S)\le V(\delta,\eps)\le2\varphi(a)\,\mu(S),\qquad\tfrac12\varphi'(a)\,\kappa\,\mu(S)\le V'(\delta,\eps)\le2\varphi'(a)\,\kappa\,\mu(S).
\]
Hence, with $k_1=\kappa\varphi'(a)/(4\varphi(a))$ and $k_2=4\kappa\varphi'(a)/\varphi(a)$,
\begin{equation}\label{eq:prior-compare}
k_1\,V(\delta,\eps)\le V'(\delta,\eps)\le k_2\,V(\delta,\eps)\qquad\text{for all }\delta\in\,]0,\min(\theta,\delta_{\mathcal W})[\text{ and }\eps>0 .
\end{equation}

\emph{Step 4: the quantifiers.} Suppose that $f$ has local volume order $(\lambda,m)$ for $(\mu,\varphi)$, with radius $\delta_0$. Put $\delta'_0=\min(\delta_0,\theta)$. Let $\delta\in\,]0,\delta'_0[$, and let $(c,C,\eps_0)$ be the constants of~\eqref{eq:order} for $V$ at the radius $\delta$. By~\eqref{eq:prior-compare}, $V'$ satisfies~\eqref{eq:order} at the radius $\delta$ with the constants $(k_1c,\,k_2C,\,\eps_0)$; note $k_1c\le k_2C$. So $f$ has local volume order $(\lambda,m)$ for $(\mu',\varphi')$, with radius $\delta'_0$. The converse follows by exchanging the roles of $(\mu,\varphi)$ and $(\mu',\varphi')$.
\end{proof}

\begin{lemma}[Comparison]\label{lem:compare}
Let $g_1,g_2:U\to[0,+\infty[$ be continuous with $g_1(a)=g_2(a)=0$. Suppose that there are $\eta>0$ and real numbers $0<\kappa_1\le\kappa_2$ such that $B_\eta(a)\subseteq U$ and
\[
\kappa_1\,g_1(x)\le g_2(x)\le\kappa_2\,g_1(x)\qquad\text{for all }x\in B_\eta(a).
\]
Let $\mathcal W$ be a window system at $a$ in $U$, $\mu$ a Haar measure and $\varphi$ a prior. For all $(\lambda,m)\in\R\times\N$, $g_1$ has local volume order $(\lambda,m)$ at $a$ with respect to $\mathcal W$ if and only if $g_2$ does.
\end{lemma}
\begin{proof}
Write $V_i=V^{\mathcal W,\mu,\varphi}_{g_i,a}$ for $i=1,2$. Let $\theta>0$ be such that $W_\delta\subseteq B_\eta(a)$ for $\delta<\min(\theta,\delta_{\mathcal W})$. For such $\delta$ and every $\eps>0$,
\[
\{x\in W_\delta:\ g_1(x)<\eps/\kappa_2\}\subseteq\{x\in W_\delta:\ g_2(x)<\eps\}\subseteq\{x\in W_\delta:\ g_1(x)<\eps/\kappa_1\},
\]
hence $V_1(\delta,\eps/\kappa_2)\le V_2(\delta,\eps)\le V_1(\delta,\eps/\kappa_1)$. Suppose that $g_1$ has local volume order $(\lambda,m)$ with radius $\delta_0$, and put $\delta'_0=\min(\delta_0,\theta)$. Let $\delta\in\,]0,\delta'_0[$, let $(c,C,\eps_0)$ be the constants of~\eqref{eq:order} for $V_1$ at the radius $\delta$, and put
\[
\eps'_0=\min\bigl(\kappa_1\eps_0,\ \eps_0,\ e^{-2\abs{\log\kappa_1}},\ e^{-2\abs{\log\kappa_2}}\bigr)\in\,]0,e^{-1}[ .
\]
Let $\eps\in\,]0,\eps'_0[$. Then $\eps<\kappa_1\eps_0<\kappa_1\le\kappa_2$ and $\eps<1$, and $L(\eps)>2\abs{\log\kappa_i}$ because $\eps<e^{-2\abs{\log\kappa_i}}$ ($i=1,2$); so the hypotheses of~\eqref{eq:L-shift} hold for $\kappa=\kappa_1$ and $\kappa=\kappa_2$; moreover $\eps/\kappa_2\le\eps/\kappa_1<\eps_0$, so the bounds~\eqref{eq:order} for $V_1$ apply at $\eps/\kappa_1$ and $\eps/\kappa_2$, and $\frac{1}{2}L(\eps)\le L(\eps/\kappa_i)\le2L(\eps)$ for $i=1,2$ by~\eqref{eq:L-shift}. Therefore
\begin{gather*}
V_2(\delta,\eps)\le C\Bigl(\frac{\eps}{\kappa_1}\Bigr)^\lambda L(\eps/\kappa_1)^{m-1}\le C\kappa_1^{-\lambda}2^{m-1}\,\eps^\lambda L(\eps)^{m-1},\\
V_2(\delta,\eps)\ge c\Bigl(\frac{\eps}{\kappa_2}\Bigr)^\lambda L(\eps/\kappa_2)^{m-1}\ge c\,\kappa_2^{-\lambda}2^{1-m}\,\eps^\lambda L(\eps)^{m-1},
\end{gather*}
and $c\kappa_2^{-\lambda}2^{1-m}\le C\kappa_1^{-\lambda}2^{m-1}$. So $g_2$ has local volume order $(\lambda,m)$, with radius $\delta'_0$. For the converse apply the same argument to $\kappa_2^{-1}g_2\le g_1\le\kappa_1^{-1}g_2$.
\end{proof}

\begin{lemma}[Change of variables]\label{lem:diffeo}
Let $(E',\norm\cdot')$ be a finite-dimensional real normed vector space, $U'\subseteq E'$ an open subset, $a'\in U'$, and $\Psi:U'\to U$ a $C^1$ diffeomorphism from $U'$ onto $U$ with $\Psi(a')=a$. For all $(\lambda,m)\in\R\times\N$, $f$ has local volume order $(\lambda,m)$ at $a$ if and only if $f\circ\Psi:U'\to\R$ has local volume order $(\lambda,m)$ at $a'$.
\end{lemma}
\begin{proof}
Since $\Psi$ is a diffeomorphism, $\dim E'=\dim E=:N$. Choose linear isomorphisms $\R^N\to E$ and $\R^N\to E'$, and let $\mu$ and $\mu'$ be the images of Lebesgue measure; we use the constant prior $1$ on $E$ and on $E'$. By Lemma~\ref{lem:order-basic}(b) (whose proof below uses only Lemmas~\ref{lem:sandwich},~\ref{lem:unique} and~\ref{lem:prior}), it suffices to prove the claim for these measures and priors. Write $J(x)=\abs{\det d\Psi(x)}$, the absolute value of the Jacobian determinant of $\Psi$ at $x\in U'$, computed in the chosen coordinates.

Choose $r_0>0$ with $\overline{B'}_{r_0}(a')\subseteq U'$. The function $J$ is continuous and positive on this compact ball, so there is $k\ge1$ with $k^{-1}\le J\le k$ on it. Let $\mathcal W=(\Psi(B'_\delta(a')))_{\delta\in\,]0,r_0[}$; by Example~\ref{ex:windows}(c) it is a window system at $a$ in $U$, equivalent to the ball system. For $\delta\in\,]0,r_0[$ and $\eps>0$ put $S'=\{x\in B'_\delta(a'):\ \abs{f(\Psi(x))-f(a)}<\eps\}$, an open set; then $\Psi(S')=\{y\in\Psi(B'_\delta(a')):\ \abs{f(y)-f(a)}<\eps\}$. The change of variables formula gives $\mu(\Psi(S'))=\int_{S'}J\,d\mu'$, so
\[
k^{-1}\,V^{\mu',1}_{f\circ\Psi,a'}(\delta,\eps)\ \le\ V^{\mathcal W,\mu,1}_{f,a}(\delta,\eps)\ \le\ k\,V^{\mu',1}_{f\circ\Psi,a'}(\delta,\eps)\qquad(\delta\in\,]0,r_0[,\ \eps>0).
\]
As in Step~4 of the proof of Lemma~\ref{lem:prior}, the constant $k$ is absorbed into the constants $c$ and $C$: $f\circ\Psi$ has local volume order $(\lambda,m)$ at $a'$ if and only if $f$ has local volume order $(\lambda,m)$ at $a$ with respect to $\mathcal W$. By Lemma~\ref{lem:sandwich}, the latter holds if and only if $f$ has local volume order $(\lambda,m)$ at $a$.
\end{proof}

\begin{lemma}[Spectator variables]\label{lem:spectator}
Let $(E_1,\norm\cdot_1)$ and $(E_2,\norm\cdot_2)$ be finite-dimensional real normed vector spaces, let $U_1\subseteq E_1$ be open, $a_1\in U_1$, $a_2\in E_2$, and $f_1:U_1\to\R$ continuous. Define $\tilde f:U_1\times E_2\to\R$ by $\tilde f(x_1,x_2)=f_1(x_1)$. For all $(\lambda,m)\in\R\times\N$, $\tilde f$ has local volume order $(\lambda,m)$ at $(a_1,a_2)$ if and only if $f_1$ has local volume order $(\lambda,m)$ at $a_1$.
\end{lemma}
\begin{proof}
By Lemma~\ref{lem:order-basic}(b) we may choose the norm, the measures and the priors. Give $E_1\times E_2$ the norm $\norm{(x_1,x_2)}=\max(\norm{x_1}_1,\norm{x_2}_2)$, whose balls are $B_\delta((a_1,a_2))=B^{E_1}_\delta(a_1)\times B^{E_2}_\delta(a_2)$. Let $\mu_1,\mu_2$ be Haar measures on $E_1,E_2$; then $\mu_1\otimes\mu_2$ is a Haar measure on $E_1\times E_2$. Use the priors $1$. For $\delta>0$ with $B^{E_1}_\delta(a_1)\subseteq U_1$ (equivalently, $B_\delta((a_1,a_2))\subseteq U_1\times E_2$) and $\eps>0$,
\[
\{x\in B_\delta((a_1,a_2)):\ \abs{\tilde f(x)-\tilde f(a_1,a_2)}<\eps\}=S_{f_1,a_1}(\delta,\eps)\times B^{E_2}_\delta(a_2),
\]
so by the definition of the product measure
\[
V^{\mu_1\otimes\mu_2,1}_{\tilde f,(a_1,a_2)}(\delta,\eps)=V^{\mu_1,1}_{f_1,a_1}(\delta,\eps)\cdot\mu_2\bigl(B^{E_2}_\delta(a_2)\bigr).
\]
The factor $\mu_2(B^{E_2}_\delta(a_2))$ is a positive real number that depends on $\delta$ but not on $\eps$. Since the constants $c$ and $C$ in~\eqref{eq:order} may depend on $\delta$, the bounds~\eqref{eq:order} hold for $\tilde f$ at the radius $\delta$ if and only if they hold for $f_1$ at the radius $\delta$ (with the constants multiplied, respectively divided, by this factor), and the radii $\delta_0$ correspond.
\end{proof}

\begin{lemma}[A nondegenerate quadratic block]\label{lem:gauss}
Let $p,k\in\N$, let $U\subseteq\R^p$ be open, $a\in U$, and $g:U\to[0,+\infty[$ continuous with $g(a)=0$. Define $G:U\times\R^k\to[0,+\infty[$ by $G(x,y)=g(x)+\norm y^2$. For all $(\lambda,m)\in\R\times\N$, $g$ has local volume order $(\lambda,m)$ at $a$ if and only if $G$ has local volume order $(\lambda+k/2,m)$ at $(a,0)$.
\end{lemma}
\begin{proof}
If $k=0$ there is nothing to prove; let $k\ge1$. We use Lebesgue measures, the priors $1$, and on $\R^p\times\R^k$ the norm $\max(\norm x,\norm y)$, whose balls are $W_\delta=B_\delta(a)\times\{y:\ \norm y<\delta\}$ (Lemma~\ref{lem:order-basic}(b)). Let $\delta\in\,]0,\delta_U[$, write $v_\delta(t)=V^{\vol,1}_{g,a}(\delta,t)=\vol\{x\in B_\delta(a):\ g(x)<t\}$ for $t>0$ (a nondecreasing function), and let $\omega_k$ be the volume of the Euclidean unit ball of $\R^k$. For $0<\eps<\delta^2$, the condition $\norm y^2<\eps$ implies $\norm y<\delta$, so Fubini's theorem gives
\[
V^{\vol,1}_{G,(a,0)}(\delta,\eps)=\int_{\{y\in\R^k:\ \norm y^2<\eps\}}v_\delta\bigl(\eps-\norm y^2\bigr)\,dy .
\]
Bounding the integrand above by $v_\delta(\eps)$, and below by $v_\delta(\eps/2)$ on $\{\norm y^2<\eps/2\}$, we get, for $0<\eps<\delta^2$,
\begin{equation}\label{eq:gauss}
\omega_k(\eps/2)^{k/2}\,v_\delta(\eps/2)\ \le\ V^{\vol,1}_{G,(a,0)}(\delta,\eps)\ \le\ \omega_k\,\eps^{k/2}\,v_\delta(\eps).
\end{equation}
Suppose $c\,\eps^\lambda L(\eps)^{m-1}\le v_\delta(\eps)\le C\,\eps^\lambda L(\eps)^{m-1}$ for $\eps\in\,]0,\eps_0[$. For $\eps<\min(\eps_0,\delta^2,\frac{1}{4})$ we have $\frac{1}{2}L(\eps)\le L(\eps/2)\le2L(\eps)$ by~\eqref{eq:L-shift}, and~\eqref{eq:gauss} gives
\[
\omega_k2^{-k/2-\lambda-m+1}c\ \eps^{\lambda+k/2}L(\eps)^{m-1}\le V^{\vol,1}_{G,(a,0)}(\delta,\eps)\le\omega_kC\ \eps^{\lambda+k/2}L(\eps)^{m-1}.
\]
Conversely, suppose $c\,\eps^{\lambda+k/2}L(\eps)^{m-1}\le V^{\vol,1}_{G,(a,0)}(\delta,\eps)\le C\,\eps^{\lambda+k/2}L(\eps)^{m-1}$ for $\eps\in\,]0,\eps_0[$. The right inequality of~\eqref{eq:gauss} gives $v_\delta(\eps)\ge\omega_k^{-1}c\,\eps^\lambda L(\eps)^{m-1}$; since $v_\delta(\eps)\le\vol B_\delta(a)<+\infty$ for all $\eps$, this forces $\lambda\ge0$. Next, the left inequality of~\eqref{eq:gauss} at $2\eps$ gives $v_\delta(\eps)\le\omega_k^{-1}\eps^{-k/2}V^{\vol,1}_{G,(a,0)}(\delta,2\eps)\le\omega_k^{-1}2^{\lambda+k/2}C\,\eps^\lambda L(\eps)^{m-1}$, for $\eps<\min(\eps_0/2,\delta^2/2,\frac{1}{4})$, using $L(2\eps)\le L(\eps)$. In both directions the radii $\delta$ correspond, since $W_\delta\subseteq U\times\R^k$ if and only if $B_\delta(a)\subseteq U$.
\end{proof}

\begin{lemma}[Homogeneity]\label{lem:homog-order}
Let $\Lambda:E\to E$ be a linear isomorphism, $\kappa>0$ a real number, and $H:E\to\R$ a continuous function with $H(\Lambda z)=\kappa H(z)$ for all $z\in E$. Then for every $z\in E$ and all $(\lambda,m)\in\R\times\N$, $H$ has local volume order $(\lambda,m)$ at $z$ if and only if it has local volume order $(\lambda,m)$ at $\Lambda z$.
\end{lemma}
\begin{proof}
By Lemma~\ref{lem:diffeo} applied to $\Psi=\Lambda$, $H$ has local volume order $(\lambda,m)$ at $\Lambda z$ if and only if $H\circ\Lambda=\kappa H$ has local volume order $(\lambda,m)$ at $z$. By Lemma~\ref{lem:compare}, applied to $g_1=\abs{H-H(z)}$ and $g_2=\kappa\abs{H-H(z)}$, this holds if and only if $H$ has local volume order $(\lambda,m)$ at $z$; note that $\abs{H(x)-H(z)}<\eps$ is the condition defining the volume functions of $H$ and of $\abs{H-H(z)}$ alike.
\end{proof}

\begin{proof}[Proof of Lemma~\ref{lem:order-basic}]
(a) is Lemma~\ref{lem:unique} for the ball system. (b) By Lemma~\ref{lem:prior} we may replace $(\mu,\varphi)$ by $(\mu',\varphi')$; by Example~\ref{ex:windows}(a) and Lemma~\ref{lem:sandwich} we may replace the balls of $\norm\cdot$ by those of $\norm\cdot'$. (c) Suppose $f$ is constant on $B_{\delta_1}(a)\subseteq U$. For $\delta\in\,]0,\delta_1[$ and every $\eps>0$, $S_{f,a}(\delta,\eps)=B_\delta(a)$, so $V^{\mu,\varphi}_{f,a}(\delta,\eps)=\int^-_{B_\delta(a)}\varphi^+\,d\mu$. For $\delta$ small this is a real number in $[\frac{1}{2}\varphi(a)\mu(B_\delta(a)),\,2\varphi(a)\mu(B_\delta(a))]\subseteq\,]0,+\infty[$ (Step~2 of the proof of Lemma~\ref{lem:prior}), so~\eqref{eq:order} holds with $(\lambda,m)=(0,1)$.
\end{proof}

The volume functions of $f$ at $a$ coincide with those of $\abs{f-f(a)}$; this is how Lemmas~\ref{lem:compare} and~\ref{lem:homog-order} apply to functions that do not vanish at the base point. Below, all losses vanish at the base point.
\section{Orbits of chains}\label{sec:quiver}

Throughout this section $D\in\N^*$, $n\in\N^{D+1}$ is a width vector, and $h$ is a barcode of length $D$ with $\dim(h)=n$. A total order on $\Sl(h)$ is fixed, so that the interval chain $\I(h)\in\Rep(n)$ and the bases $(\epsilon^v_s)_{s\in\Sl_v(h)}$ of the spaces $\R^{n_v}$ are defined (Definition~\ref{def:interval-chain}). We denote by
\[
O_h=G(n)\cdot\I(h)=\{g\cdot\I(h):\ g\in G(n)\}
\]
the orbit of the interval chain $\I(h)$ under the action of the group $G(n)$. By Lemma~\ref{lem:barcode-of-chain} and Definition~\ref{def:barcode-of-chain}, $O_h$ does not depend on the total order on $\Sl(h)$, and
\begin{equation}\label{eq:orbit-ranks}
O_h=\{A\in\Rep(n):\ R(A)=\rho(h)\}=\{A\in\Rep(n):\ \beta(A)=h\}.
\end{equation}

\subsection{Orbits are submanifolds}
Let $\partial_h:\g(n)\to\Rep(n)$ be the derivative at the identity of the orbit map $G(n)\to\Rep(n)$, $g\mapsto g\cdot\I(h)$, where $\g(n)$ is the Lie algebra of $G(n)$ (Section~\ref{sec:chains}). It is the linear map
\[
(\partial_h\xi)_k=\xi_{k+1}\,\I(h)_k-\I(h)_k\,\xi_k\qquad(\xi=(\xi_v)_v\in\g(n),\ k\in\ii0{D-1}).
\]
Its kernel is the Lie algebra of the stabilizer $\{g\in G(n):\ g\cdot\I(h)=\I(h)\}$ of $\I(h)$.

\begin{lemma}\label{lem:orbit}
The orbit $O_h$ is an embedded real-analytic submanifold of $\Rep(n)$, with tangent space $\partial_h(\g(n))$ at $\I(h)$, and $\dim O_h=\dim\g(n)-\dim\ker\partial_h$. If $\g'\subseteq\g(n)$ is a linear complement of $\ker\partial_h$, then $\xi\mapsto(\exp\xi)\cdot\I(h)$ maps a neighbourhood of $0$ in $\g'$ diffeomorphically onto a neighbourhood of $\I(h)$ in $O_h$.
\end{lemma}
\begin{proof}
By~\eqref{eq:orbit-ranks}, $O_h$ is the finite intersection of the sets $\{A:\ \rk P_{ij}(A)=\rho_{ij}(h)\}$. Each of them is the intersection of the open set $\{\rk P_{ij}\ge\rho_{ij}(h)\}$ with the closed set $\{\rk P_{ij}\le\rho_{ij}(h)\}$, so $O_h$ is locally closed in $\Rep(n)$, hence locally compact. The Lie group $G(n)$ has finitely many connected components and acts transitively on $O_h$. By~\cite[Chapter~II, Theorem~3.2]{helgason}, the orbit map induces a homeomorphism from the homogeneous space of $G(n)$ modulo the stabilizer of $\I(h)$ onto $O_h$. This map is also an injective immersion~\cite[Chapter~21]{lee}. Hence it is a diffeomorphism onto an embedded submanifold, with the stated tangent space and dimension; the last assertion follows from the inverse function theorem. The submanifold is real-analytic because the action is algebraic.
\end{proof}

\subsection{Linked pairs and the normal space}
For a linked pair $\pi=(s,s')\in\Lk(h)$, let $J_\pi\in\Rep(n)$ be the chain whose only nonzero entry is a $1$ in the matrix of the arrow $b_{s'}-1$, in the column of $\epsilon^{b_{s'}-1}_s$ and the row of $\epsilon^{b_{s'}}_{s'}$. This makes sense: $b_s\le b_{s'}-1\le e_s$, so $s\in\Sl_{b_{s'}-1}(h)$, and $s'\in\Sl_{b_{s'}}(h)$. For $x=(x_\pi)_{\pi\in\Lk(h)}\in\R^{\Lk(h)}$ put $Jx=\sum_{\pi\in\Lk(h)}x_\pi J_\pi\in\Rep(n)$.

\begin{lemma}\label{lem:ext}
$\Rep(n)=\partial_h(\g(n))\oplus\bigoplus_{\pi\in\Lk(h)}\R J_\pi$. Consequently $\codim O_h=\#\Lk(h)$.
\end{lemma}
\begin{proof}
Write elements of $\g(n)$ and of $\Rep(n)$ in the bases $(\epsilon^v_s)$. An entry of $\xi_v$ is indexed by a vertex $v$ and an ordered pair of slots $(s,s')$ with $v\in I_s\cap I_{s'}$ (column $s$, row $s'$). An entry of $Y\in\Rep(n)$ is indexed by an arrow $k$ and an ordered pair of slots $(s,s')$ with $k\in I_s$ and $k+1\in I_{s'}$. Since every matrix $\I(h)_k$ maps basis vectors to basis vectors of the same slot or to $0$, the map $\partial_h$ respects the decomposition according to the pair $(s,s')$. On the summand of $(s,s')$ it is the map
\[
\partial_{s,s'}:\R^{I_s\cap I_{s'}}\to\R^{\mathrm{Ar}(s,s')},\qquad(\partial_{s,s'}u)_k=u_{k+1}-u_k,
\]
where $\mathrm{Ar}(s,s')=\{k\in\ii0{D-1}:\ k\in I_s,\ k+1\in I_{s'}\}$, and $u_v$ is read as $0$ for $v\notin I_s\cap I_{s'}$. Write $I_s=\ii be$ and $I_{s'}=\ii{b'}{e'}$.

\emph{Linked pairs.} If $b<b'\le e+1\le e'$, then $I_s\cap I_{s'}=\ii{b'}e$ (empty if $b'=e+1$) and $\mathrm{Ar}(s,s')=\ii{b'-1}e$. In this case $\partial_{s,s'}$ is injective: the coordinate $k=b'-1$ gives $u_{b'}=0$, and inductively $u=0$. Its image lies in the hyperplane $\{y:\ \sum_ky_k=0\}$, by telescoping, and has dimension $e-b'+1=\#\mathrm{Ar}(s,s')-1$. So the image is exactly this hyperplane, and $J_\pi$, whose coordinate sum is $1$, spans a complement.

\emph{Non-linked pairs.} We show that $\partial_{s,s'}$ is surjective. If $\mathrm{Ar}(s,s')=\emptyset$ there is nothing to prove; otherwise $b\le e'-1$ and $b'\le e+1$.
\begin{itemize}
\item If $b'\le b$ and $e'\le e$, then $I_s\cap I_{s'}=\ii b{e'}$ and $\mathrm{Ar}(s,s')=\ii b{e'-1}$. For any $y\in\R^{\mathrm{Ar}(s,s')}$, set $u_b=0$ and $u_{v+1}=u_v+y_v$ for $v\in\ii b{e'-1}$. Then $\partial_{s,s'}u=y$, so the map is surjective.
\item If $b'\le b$ and $e'>e$, then both index sets equal $\ii be$, and $\partial_{s,s'}$ is triangular with diagonal entries $-1$, since $(\partial_{s,s'}u)_e=-u_e$.
\item If $b'>b$, then non-linkedness and $b'\le e+1$ force $e'\le e$. So $I_s\cap I_{s'}=\ii{b'}{e'}$ and $\mathrm{Ar}(s,s')=\ii{b'-1}{e'-1}$. Then $(\partial_{s,s'}u)_{b'-1}=u_{b'}$, and $\partial_{s,s'}$ is triangular with diagonal entries $1$.
\end{itemize}
Summing over all pairs gives the decomposition, and $\codim O_h=\dim\Rep(n)-\dim\partial_h(\g(n))=\#\Lk(h)$ by Lemma~\ref{lem:orbit}.
\end{proof}

The normal-space calculation above is the interval-chain instance of Voigt's lemma~\cite{voigt}. Its proof here uses only the displayed difference maps.

\subsection{The slice chart}
Fix a linear complement $\g'_h$ of $\ker\partial_h$ in $\g(n)$ and define
\[
\Psi_h:\g'_h\times\R^{\Lk(h)}\to\Rep(n),\qquad\Psi_h(\xi,x)=(\exp\xi)\cdot\bigl(\I(h)+Jx\bigr).
\]

\begin{lemma}[Slice chart]\label{lem:slice}
There is an open neighbourhood $U_h$ of $(0,0)$ in $\g'_h\times\R^{\Lk(h)}$ such that $\Psi_h$ maps $U_h$ real-analytically and diffeomorphically onto an open neighbourhood of $\I(h)$ in $\Rep(n)$, and $\Psi_h(U_h)\cap O_h=\Psi_h\bigl(U_h\cap(\g'_h\times\{0\})\bigr)$. In particular $\I(h)+Jx\notin O_h$ for every $x\neq0$ with $(0,x)\in U_h$.
\end{lemma}
\begin{proof}
The differential $d\Psi_h$ at $(0,0)$ is $(\xi,x)\mapsto\partial_h\xi+Jx$, which is bijective by Lemma~\ref{lem:ext}, since $\partial_h$ is injective on $\g'_h$. The inverse function theorem gives a neighbourhood $U_1$ of $(0,0)$ on which $\Psi_h$ is a diffeomorphism onto an open set. By Lemma~\ref{lem:orbit} there are an open neighbourhood $\mathcal V$ of $\I(h)$ in $\Rep(n)$ and a neighbourhood $\mathcal U$ of $0$ in $\g'_h$ with $\mathcal U\times\{0\}\subseteq U_1$ and $O_h\cap\mathcal V=\{(\exp\xi)\cdot\I(h):\ \xi\in\mathcal U\}$. Let $U_h\subseteq U_1$ be an open neighbourhood of $(0,0)$ with $\Psi_h(U_h)\subseteq\mathcal V$ and $U_h\cap(\g'_h\times\{0\})\subseteq\mathcal U\times\{0\}$. If $\Psi_h(\xi,x)\in O_h$ with $(\xi,x)\in U_h$, then $\Psi_h(\xi,x)=\Psi_h(\xi',0)$ for some $\xi'\in\mathcal U$, and injectivity of $\Psi_h$ on $U_1$ gives $x=0$.
\end{proof}

\begin{lemma}[Path expansion]\label{lem:path}
Let $x\in\R^{\Lk(h)}$, let $s$ be an input slot and $s'$ an output slot of $h$. The entry of $P(\I(h)+Jx)$ in the column of $\epsilon^0_s$ and the row of $\epsilon^D_{s'}$ equals
\[
\sum_{s=s_0\lin s_1\lin\cdots\lin s_p=s'}\ \prod_{l=1}^p x_{(s_{l-1},s_l)},
\]
the sum over all $p\in\N$ and all sequences $(s_0,\dots,s_p)$ of slots with $s_0=s$, $s_p=s'$ and $(s_{l-1},s_l)\in\Lk(h)$ for $l\in\ii1p$. The sequence with $p=0$ occurs only if $s=s'$, which is then a full-span slot.
\end{lemma}
\begin{proof}
At the arrow $k$, the matrix $(\I(h)+Jx)_k$ sends $\epsilon^k_u$ to $\epsilon^{k+1}_u$ if $k+1\in I_u$, and to $0$ otherwise, plus $\sum x_{(u,u')}\epsilon^{k+1}_{u'}$ over the linked pairs $(u,u')$ with $b_{u'}=k+1$. Expanding the product $(\I(h)+Jx)_{D-1}\cdots(\I(h)+Jx)_0$, the entry in question is a sum over the sequences $(u_0,\dots,u_D)$ of slots with $u_0=s$, $u_D=s'$ and $k\in I_{u_k}$ for all $k$. At each arrow $k$, either $u_{k+1}=u_k$ (weight $1$, possible only if $k+1\in I_{u_k}$), or $u_k\lin u_{k+1}$ with $b_{u_{k+1}}=k+1$ (weight $x_{(u_k,u_{k+1})}$). The distinct slots visited form a sequence $s=s_0\lin\cdots\lin s_p=s'$ of linked pairs with $b_{s_1}<\dots<b_{s_p}$. Conversely, such a sequence determines $(u_k)$ uniquely: $u_k=s_l$ for $b_{s_l}\le k<b_{s_{l+1}}$, with the conventions $b_{s_0}=0$ and $b_{s_{p+1}}=D+1$. This sequence is admissible, since $b_{s_{l+1}}\le e_{s_l}+1$ means that $s_l$ is present until the jump, and $e_{s'}=D$. So each sequence of linked pairs contributes exactly the product of its coordinates.
\end{proof}

\emph{Full-span slots.} Let $r=h(\ii0D)$ be the number of full-span slots of $h$, and let $h^\circ$ be the zero-output part of $h$, so that $\Sl(h^\circ)\subseteq\Sl(h)$ consists of the slots that are not full-span and $\dim(h^\circ)=n-r$ by~\eqref{eq:zero-output-part}. Since no linked pair involves a full-span slot, $\Lk(h)=\Lk(h^\circ)$. Suppose that the full-span slots come first in the total order on $\Sl(h)$, and give $\Sl(h^\circ)$ the induced order. Then at each vertex $v$ the first $r$ basis vectors $\epsilon^v_s$ belong to the full-span slots, and in the decomposition $\R^{n_v}=\R^r\oplus\R^{n_v-r}$ given by this splitting of the basis,
\begin{equation}\label{eq:blocks}
\I(h)+Jx=(\id_r,\dots,\id_r)\oplus\bigl(\I(h^\circ)+J^\circ x\bigr),\qquad P(\I(h)+Jx)=\id_r\oplus P\bigl(\I(h^\circ)+J^\circ x\bigr)
\end{equation}
for all $x\in\R^{\Lk(h)}$, where $J^\circ$ is the map $J$ defined for the barcode $h^\circ$, and $(\id_r,\dots,\id_r)\in\Rep(r,\dots,r)$.

\subsection{Switches and the closure order}
\begin{lemma}[Switch]\label{lem:switch}
Let $\pi=(s,s')\in\Lk(h)$, and define the barcode
\[
h^\pi=h-\one_{I_s}-\one_{I_{s'}}+\one_{\ii{b_s}{e_{s'}}}+\one_{\ii{b_{s'}}{e_s}},
\]
where the last term is omitted if $b_{s'}=e_s+1$. Then $\dim(h^\pi)=n$, and $\I(h)+\tau J_\pi\in O_{h^\pi}$ for every real number $\tau\neq0$. Hence $O_h\subseteq\overline{O_{h^\pi}}$. Moreover, for $0\le i\le j\le D$, $\rho_{ij}(h^\pi)=\rho_{ij}(h)+1$ if $b_s\le i<b_{s'}$ and $e_s<j\le e_{s'}$, and $\rho_{ij}(h^\pi)=\rho_{ij}(h)$ otherwise.
\end{lemma}
\begin{proof}
Since $b_s<b_{s'}$, the intervals $I_s$ and $I_{s'}$ are distinct and $h(I_s),h(I_{s'})\ge1$, so $h^\pi$ has values in $\N$. A vertex $v$ lies in as many of the intervals $\ii{b_s}{e_{s'}}$, $\ii{b_{s'}}{e_s}$ as of the intervals $I_s$, $I_{s'}$, because $b_s<b_{s'}\le e_s+1\le e_{s'}$; hence $\dim(h^\pi)=\dim(h)$. Fix $\tau\neq0$ and consider the vectors
\[
w^v=\begin{cases}\epsilon^v_s&\text{if }v\in\ii{b_s}{b_{s'}-1},\\ \epsilon^v_s+\tau\epsilon^v_{s'}&\text{if }v\in\ii{b_{s'}}{e_s},\\ \tau\epsilon^v_{s'}&\text{if }v\in\ii{e_s+1}{e_{s'}}.\end{cases}
\]
Put $Y=\I(h)+\tau J_\pi$. We group vectors into \emph{strings}: a string with interval $\ii bc\in\Int_D$ is a family $(z^v)_{v\in\ii bc}$ with $z^v\in\R^{n_v}$, such that $Y_v z^v=z^{v+1}$ for $b\le v<c$ and $Y_cz^c=0$ if $c<D$. The following are strings:
\begin{itemize}
\item $(w^v)_{v\in\ii{b_s}{e_{s'}}}$, with interval $\ii{b_s}{e_{s'}}$: at the arrow $b_{s'}-1$ the term $\tau J_\pi$ adds $\tau\epsilon^{b_{s'}}_{s'}$, and at the other arrows $Y$ acts as $\I(h)$;
\item $(\epsilon^v_s)_{v\in\ii{b_{s'}}{e_s}}$, with interval $\ii{b_{s'}}{e_s}$ (if $b_{s'}\le e_s$), since $Y$ acts as $\I(h)$ at the arrows $v\ge b_{s'}$;
\item $(\epsilon^v_t)_{v\in I_t}$ for every slot $t\notin\{s,s'\}$, with interval $I_t$.
\end{itemize}
At each vertex $v$, the vectors of these strings present at $v$ form a basis of $\R^{n_v}$, since $\tau\neq0$. For $i\le j$, the map $P_{ij}(Y)$ therefore sends the basis vectors at $i$ of the strings whose interval contains $\ii ij$ to distinct basis vectors at $j$, and the other basis vectors at $i$ to $0$. Hence $\rk P_{ij}(Y)$ is the number of strings whose interval contains $\ii ij$. The intervals of the strings are those of $h^\pi$, counted with multiplicity, so $R(Y)=\rho(h^\pi)$ and $Y\in O_{h^\pi}$ by~\eqref{eq:orbit-ranks}. Letting $\tau\to0$ gives $\I(h)\in\overline{O_{h^\pi}}$, hence $O_h\subseteq\overline{O_{h^\pi}}$, because $\overline{O_{h^\pi}}$ is $G(n)$-invariant. For the ranks: the interval $\ii ij$ is contained in $\ii{b_s}{e_{s'}}$ or $\ii{b_{s'}}{e_s}$ as often as in $I_s$ or $I_{s'}$, except when $b_s\le i<b_{s'}$ and $e_s<j\le e_{s'}$, in which case it is contained in exactly one of the new intervals and in neither of the old ones.
\end{proof}

\begin{definition}\label{def:geom-order}
For barcodes $h,k$ of length $D$ with $\dim(h)=\dim(k)=n$, we write $h\sqsubseteq k$ if $O_h\subseteq\overline{O_k}$. Equivalently, $\I(h)\in\overline{O_k}$, since $\overline{O_k}$ is $G(n)$-invariant.
\end{definition}

\begin{proposition}[Closure order]\label{prop:closure}
Let $h,k$ be barcodes of length $D$ with $\dim(h)=\dim(k)=n$.
\begin{enumerate}[label=\textup{(\alph*)}]
\item If $h\sqsubseteq k$, then $h\preceq k$. Hence $\sqsubseteq$ is a partial order on the barcodes of length $D$ with dimension vector $n$.
\item If $h\preceq k$, then $h\sqsubseteq k$.
\end{enumerate}
\end{proposition}
\begin{proof}
(a) The set $\{A\in\Rep(n):\ \rk P_{ij}(A)\le\rho_{ij}(k)\text{ for all }0\le i\le j\le D\}$ is closed and contains $O_k$, hence it contains $\overline{O_k}\ni\I(h)$; so $\rho(h)=R(\I(h))\le\rho(k)$. Reflexivity and transitivity of $\sqsubseteq$ are clear, and antisymmetry follows from that of $\preceq$ (Fact~\ref{fact:gabriel}(c)).

(b) By Proposition~\ref{prop:generation} (Appendix~\ref{app:generation}), if $h\preceq k$ there is a sequence $h=h_0,h_1,\dots,h_\ell=k$ in which each $h_{t+1}$ is obtained from $h_t$ by a switch along a linked pair of $h_t$ (Lemma~\ref{lem:switch}). By Lemma~\ref{lem:switch}, $h_t\sqsubseteq h_{t+1}$ for each $t$, and $\sqsubseteq$ is transitive.
\end{proof}
Part (b) is classical over algebraically closed fields~\cite{abeasis-delfra-kraft}; over $\R$ we give a complete argument in Appendix~\ref{app:generation}. The stratification of the fibres of $P$ over $\R$ by the ranks of partial products, the dimensions of its strata and their closure relations are studied in detail by Shewchuk and Bhattacharya~\cite{shewchuk-bhattacharya}. Only part (a) is used in the proof of Theorem~\ref{thm:A}. Part (b) is used to express the multiplicity and the codimension through the computable order $\preceq$ in Theorem~\ref{thm:B}.

\begin{lemma}[Strict monotonicity of dimension]\label{lem:strict}
Let $k$ be a barcode of length $D$ with $\dim(k)=n$. If $h\sqsubseteq k$ and $h\neq k$, then $\dim O_h<\dim O_k$, i.e.\ $\#\Lk(k)<\#\Lk(h)$.
\end{lemma}
\begin{proof}
For a real number $t>0$ let $\Gamma(t)\in G(n)$ act at each vertex $v$ by $\epsilon^v_s\mapsto t^{b_s}\epsilon^v_s$ ($s\in\Sl_v(h)$), and let $\Delta_t$ be the linear automorphism of $\R^{\Lk(h)}$ given by $(\Delta_tx)_\pi=t^{b_{s'}-b_s}x_\pi$ for $\pi=(s,s')$. The matrices $\I(h)_j$, $j\in\ii0{D-1}$, join two basis vectors of the same slot, so $\Gamma(t)\cdot\I(h)=\I(h)$, and $\Gamma(t)\cdot J_\pi=t^{b_{s'}-b_s}J_\pi$. Hence $\Gamma(t)\cdot(\I(h)+Jx)=\I(h)+J\Delta_tx$, with $b_{s'}-b_s\ge1$ for every $\pi\in\Lk(h)$. Since $\I(h)\in\overline{O_k}$ and $\Psi_h(U_h)$ is a neighbourhood of $\I(h)$, there is $(\xi,x)\in U_h$ with $\Psi_h(\xi,x)\in O_k$, hence $\I(h)+Jx=(\exp(-\xi))\cdot\Psi_h(\xi,x)\in O_k$. Here $x\neq0$: otherwise $\Psi_h(\xi,x)\in O_h$, and $O_h\cap O_k=\emptyset$ since $h\neq k$. Replacing $x$ by $\Delta_tx$ for a small $t$ (which keeps $\I(h)+Jx$ in $O_k$), we may assume that $\{(\xi',\Delta_tx):\ \xi'\in\mathcal U',\ t\in\,]0,1[\}\subseteq U_h$ for a small open ball $\mathcal U'\ni0$ of $\g'_h$. The map $(\xi',t)\mapsto\Psi_h(\xi',\Delta_tx)=(\exp\xi')\cdot\Gamma(t)\cdot(\I(h)+Jx)$ then takes values in $O_k$. It is injective and immersive, because $\Psi_h$ is a diffeomorphism on $U_h$ and $t\mapsto\Delta_tx$ is an injective immersion (its derivative has the coordinates $(b_{s'}-b_s)t^{b_{s'}-b_s-1}x_\pi$, not all zero). Since $O_k$ is an embedded submanifold (Lemma~\ref{lem:orbit}), this is an injective immersion of a manifold of dimension $\dim\g'_h+1=\dim O_h+1$ into $O_k$, so $\dim O_k\ge\dim O_h+1$. The last assertion follows from Lemma~\ref{lem:ext}.
\end{proof}

\subsection{Strata of a fibre}
\begin{lemma}\label{lem:fibre-strata}
Let $\Phi\in\R^{n_D\times n_0}$ be of rank $r$. Then $F_\Phi$ is the disjoint union, over the barcodes $k$ of length $D$ with $\dim(k)=n$ and $k(\ii0D)=r$, of the sets $F_\Phi\cap O_k$. Each of them is empty or an embedded submanifold of $\Rep(n)$ of dimension $\dim O_k-\creg(n,r)$.
\end{lemma}
\begin{proof}
The orbits $O_k$ are pairwise disjoint and cover $\Rep(n)$ by~\eqref{eq:orbit-ranks}, and a point $A\in O_k$ has $\rk P(A)=R_{0D}(A)=k(\ii0D)$; this gives the decomposition. Let $M_r\subseteq\R^{n_D\times n_0}$ be the manifold of matrices of rank $r$, of dimension $\creg(n,r)$. The map $P|_{O_k}:O_k\to M_r$ satisfies $P(g\cdot A)=g_DP(A)g_0^{-1}$. Since $D\ge1$, the vertices $0$ and $D$ are distinct, so $g\mapsto(g_D,g_0)$ maps $G(n)$ onto $\GL_{n_D}(\R)\times\GL_{n_0}(\R)$, which acts transitively on $M_r$. Thus $P|_{O_k}$ is an equivariant map between homogeneous spaces of $G(n)$, so its rank is constant. It is surjective onto $M_r$ (when $O_k\neq\emptyset$), so by Sard's theorem its rank is $\dim M_r$: it is a submersion. Its fibres are therefore embedded submanifolds of $O_k$ of codimension $\creg(n,r)$.
\end{proof}
For $D=0$ the lemma fails as stated: the action on the target is then by conjugation $\Phi\mapsto g_0\Phi g_0^{-1}$, which is not transitive on $M_r$.
\section{Reduction to the zero target}\label{sec:reduction}

\subsection{A Schur complement estimate}
Let $p,q,r\in\N$ with $r\le\min(p,q)$, and let $\Phi_0=\begin{pmatrix}\id_r&0\\0&0\end{pmatrix}\in\R^{p\times q}$. For $M\in\R^{p\times q}$ write $M=\begin{pmatrix}M_{11}&M_{12}\\M_{21}&M_{22}\end{pmatrix}$ with $M_{11}\in\R^{r\times r}$, $M_{12}\in\R^{r\times(q-r)}$, $M_{21}\in\R^{(p-r)\times r}$ and $M_{22}\in\R^{(p-r)\times(q-r)}$. On the open set of matrices $M$ with $M_{11}$ invertible put
\begin{gather*}
\omega(M)=(M_{11}-\id_r,\ M_{12},\ M_{21})\in\R^{r\times r}\times\R^{r\times(q-r)}\times\R^{(p-r)\times r},\\
\Sch(M)=M_{22}-M_{21}M_{11}^{-1}M_{12}\in\R^{(p-r)\times(q-r)}.
\end{gather*}
The space of values of $\omega$ has dimension $r^2+r(q-r)+r(p-r)=r(p+q-r)$; we give it the Euclidean norm $\Fr{\omega(M)}=(\Fr{M_{11}-\id_r}^2+\Fr{M_{12}}^2+\Fr{M_{21}}^2)^{1/2}$.

\begin{lemma}\label{lem:schur}
\begin{enumerate}[label=\textup{(\alph*)}]
\item On a neighbourhood of $\Phi_0$ in $\R^{p\times q}$,
\[
\tfrac13\bigl(\Fr{\omega(M)}^2+\Fr{\Sch(M)}^2\bigr)\le\Fr{M-\Phi_0}^2\le3\bigl(\Fr{\omega(M)}^2+\Fr{\Sch(M)}^2\bigr),
\]
and $M=\Phi_0$ if and only if $\omega(M)=0$ and $\Sch(M)=0$.
\item There is a neighbourhood of $(\id_p,\id_q,0)$ in $\GL_p(\R)\times\GL_q(\R)\times\R^{(p-r)\times(q-r)}$ on which
\[
\tfrac12\Fr X\le\Fr{\Sch\bigl(U(\id_r\oplus X)V\bigr)}\le2\Fr X .
\]
\end{enumerate}
\end{lemma}
\begin{proof}
(a) Restrict to the matrices $M$ with $\norm{M_{11}-\id_r}_{\mathrm{op}}<\frac{1}{2}$ and $\Fr{\omega(M)}\le\frac{1}{2}$. Then $\norm{M_{11}^{-1}}_{\mathrm{op}}\le2$, and
\[
\Fr{M_{21}M_{11}^{-1}M_{12}}\le2\Fr{M_{21}}\Fr{M_{12}}\le\Fr{\omega(M)}^2\le\tfrac12\Fr{\omega(M)}.
\]
Since $\Fr{M-\Phi_0}^2=\Fr{\omega(M)}^2+\Fr{M_{22}}^2$ and $M_{22}=\Sch(M)+M_{21}M_{11}^{-1}M_{12}$, the inequality $(\alpha+\beta)^2\le2\alpha^2+2\beta^2$, applied to $\Fr{M_{22}}\le\Fr{\Sch(M)}+\frac{1}{2}\Fr{\omega(M)}$ and to $\Fr{\Sch(M)}\le\Fr{M_{22}}+\frac{1}{2}\Fr{\omega(M)}$, gives both bounds. The last assertion follows.

(b) For $U$ near $\id_p$ and $V$ near $\id_q$ we have the block factorizations
\[
U=\begin{pmatrix}\id&0\\U_{21}U_{11}^{-1}&\id\end{pmatrix}\begin{pmatrix}U_{11}&0\\0&S_U\end{pmatrix}\begin{pmatrix}\id&y_U\\0&\id\end{pmatrix},\qquad
V=\begin{pmatrix}\id&0\\y_V&\id\end{pmatrix}\begin{pmatrix}V_{11}&0\\0&S_V\end{pmatrix}\begin{pmatrix}\id&V_{11}^{-1}V_{12}\\0&\id\end{pmatrix},
\]
with $S_U=\Sch(U)$, $S_V=\Sch(V)$ (for the block sizes $(r,p-r)$ and $(r,q-r)$), $y_U=U_{11}^{-1}U_{12}$ and $y_V=V_{21}V_{11}^{-1}$. If $\mathcal L$ is block lower triangular and $\mathcal R$ is block upper triangular with invertible diagonal blocks, and $\mathcal M_{11}$ is invertible, a direct computation gives $\Sch(\mathcal L\mathcal M\mathcal R)=\mathcal L_{22}\Sch(\mathcal M)\mathcal R_{22}$. Applying this,
\[
\Sch\bigl(U(\id_r\oplus X)V\bigr)=S_U\,\Sch(\mathcal M)\,S_V,
\]
where
\[
\mathcal M=\begin{pmatrix}\id&y_U\\0&\id\end{pmatrix}\begin{pmatrix}\id&0\\0&X\end{pmatrix}\begin{pmatrix}\id&0\\y_V&\id\end{pmatrix}=\begin{pmatrix}\id+y_UXy_V&y_UX\\Xy_V&X\end{pmatrix}.
\]
By the identity $(\id+y_Vy_UX)^{-1}=\id-y_V(\id+y_UXy_V)^{-1}y_UX$, valid for small $X$,
\[
\Sch(\mathcal M)=X-Xy_V(\id+y_UXy_V)^{-1}y_UX=X(\id+y_Vy_UX)^{-1}.
\]
Hence $\Sch(U(\id_r\oplus X)V)=S_U\,X\,(\id+y_Vy_UX)^{-1}S_V$. As $(U,V,X)\to(\id_p,\id_q,0)$, the three matrices $S_U$, $(\id+y_Vy_UX)^{-1}$, $S_V$ and their inverses tend to identity matrices; so near $(\id_p,\id_q,0)$ they all have operator norm at most $2^{1/3}$. Since $\Fr{MXN}\le\norm M_{\mathrm{op}}\Fr X\norm N_{\mathrm{op}}$, this gives $\Fr{\Sch(U(\id_r\oplus X)V)}\le2\Fr X$, and, writing $X=S_U^{-1}\,\Sch(U(\id_r\oplus X)V)\,S_V^{-1}(\id+y_Vy_UX)$, also $\Fr X\le2\Fr{\Sch(U(\id_r\oplus X)V)}$.
\end{proof}

\subsection{The reduction}
In this subsection $D\in\N^*$, $n\in\N^{D+1}$, $\Phi\in\R^{n_D\times n_0}$, $A\in F_\Phi$, and $r=\rk\Phi$. By Lemma~\ref{lem:barcode-of-chain}(i), $\beta(A)$ has exactly $r$ full-span slots, its zero-output part is the residual barcode $h_A\in\Zb(n-r)$, and $\Lk(\beta(A))=\Lk(h_A)$ (Section~\ref{sec:quiver}). We fix a total order on $\Sl(\beta(A))$ in which the $r$ full-span slots come first, and give $\Sl(h_A)\subseteq\Sl(\beta(A))$ the induced order. This defines the interval chains $\I(\beta(A))\in\Rep(n)$ and $\I(h_A)\in\Rep(n-r)$.

\begin{proposition}\label{prop:reduction}
For all $(\lambda',m')\in\R\times\N$: if $L_{n-r,0}$ has local volume order $(\lambda',m')$ at $\I(h_A)$, then $L_{n,\Phi}$ has local volume order $(\tfrac12\creg(n,r)+\lambda',\,m')$ at $A$. Moreover
\[
\codim_AF_{n,\Phi}=\creg(n,r)+\codim_{\I(h_A)}Z(n-r).
\]
\end{proposition}

\begin{proof}
\emph{Step 1: normal form.} By Lemma~\ref{lem:barcode-of-chain}(i) there is $g\in G(n)$ with $g\cdot A=\I(\beta(A))$. Put $\Phi_0=P(\I(\beta(A)))=g_D\Phi g_0^{-1}$; by~\eqref{eq:blocks} (with $x=0$), $\Phi_0=\id_r\oplus0$ in the chosen bases. The linear automorphism $T:Y\mapsto g^{-1}\cdot Y$ of $\Rep(n)$ maps $\I(\beta(A))$ to $A$ and $F_{n,\Phi_0}$ onto $F_{n,\Phi}$, and
\[
L_{n,\Phi}(T(Y))=\tfrac12\Fr{g_D^{-1}(P(Y)-\Phi_0)g_0}^2\qquad(Y\in\Rep(n)),
\]
which lies between two positive multiples of $L_{n,\Phi_0}(Y)$, since $\Fr{MXN}\le\norm M_{\mathrm{op}}\Fr X\norm N_{\mathrm{op}}$. By Lemmas~\ref{lem:diffeo} and~\ref{lem:compare}, $L_{n,\Phi}$ at $A$ and $L_{n,\Phi_0}$ at $\I(\beta(A))$ have the same local volume orders, and by~\ref{H:lip} $\codim_AF_{n,\Phi}=\codim_{\I(\beta(A))}F_{n,\Phi_0}$. So it suffices to prove the proposition for the chain $\I(\beta(A))$ and the target $\Phi_0$.

\emph{Step 2: slices.} Let $\Psi_{\beta(A)}$ and $\Psi_{h_A}$ be the slice charts of $\I(\beta(A))$ and $\I(h_A)$ (Lemma~\ref{lem:slice}), with complements $\g'_{\beta(A)}\subseteq\g(n)$ and $\g'_{h_A}\subseteq\g(n-r)$ and domains $U_{\beta(A)}$ and $U_{h_A}$. Both have the same second factor $\R^{\Lk(h_A)}$. Put
\[
X(x)=P\bigl(\I(h_A)+J^\circ x\bigr)\in\R^{(n_D-r)\times(n_0-r)}\qquad(x\in\R^{\Lk(h_A)}),
\]
a polynomial map with $X(0)=0$, where $J^\circ$ is the map $J$ of $h_A$. By~\eqref{eq:blocks},
\begin{equation}\label{eq:M}
M(\xi,x):=P\bigl(\Psi_{\beta(A)}(\xi,x)\bigr)=e^{\xi_D}\bigl(\id_r\oplus X(x)\bigr)e^{-\xi_0}\qquad\bigl((\xi,x)\in\g'_{\beta(A)}\times\R^{\Lk(h_A)}\bigr).
\end{equation}
For the chain $\I(h_A)$ (the case $r=0$ of~\eqref{eq:M}), $L_{n-r,0}(\Psi_{h_A}(\xi,x))=\frac{1}{2}\Fr{e^{\xi_D}X(x)e^{-\xi_0}}^2$ lies between two positive multiples of $\Fr{X(x)}^2$ near $(0,0)$. By Lemmas~\ref{lem:diffeo},~\ref{lem:compare} and~\ref{lem:spectator}, $L_{n-r,0}$ at $\I(h_A)$ has the same local volume orders as the function
\[
x\mapsto\Fr{X(x)}^2\quad\text{at }0\in\R^{\Lk(h_A)},
\]
and $Z(n-r)\cap\Psi_{h_A}(U_{h_A})=\Psi_{h_A}\bigl((\g'_{h_A}\times Y_A)\cap U_{h_A}\bigr)$, where $Y_A=\{x\in\R^{\Lk(h_A)}:\ X(x)=0\}$.

\emph{Step 3: Schur coordinates.} Let $\omega$ be as in Lemma~\ref{lem:schur} for $(p,q)=(n_D,n_0)$; we identify its space of values, of dimension $\creg(n,r)$, with $\R^{\creg(n,r)}$ by listing the entries of the three blocks. This is an isometry for the Euclidean norm given to it before Lemma~\ref{lem:schur}, so that $\norm u=\Fr{\omega(M)}$ below. Consider the map $\theta(\xi)=\omega(e^{\xi_D}\Phi_0e^{-\xi_0})$, defined on the neighbourhood of $0$ in $\g'_{\beta(A)}$ where the upper left $r\times r$ block of $e^{\xi_D}\Phi_0e^{-\xi_0}$ is invertible, with values in $\R^{\creg(n,r)}$. Its differential at $0$ is $\xi\mapsto\omega'(\xi_D\Phi_0-\Phi_0\xi_0)$, where $\omega'(M)=(M_{11},M_{12},M_{21})$ is the linear part of $\omega$. Writing $\xi_D=\begin{pmatrix}\alpha_1&\alpha_2\\\alpha_3&\alpha_4\end{pmatrix}$ and $\xi_0=\begin{pmatrix}\gamma_1&\gamma_2\\\gamma_3&\gamma_4\end{pmatrix}$ in blocks of sizes $(r,n_D-r)$ and $(r,n_0-r)$, we get $\omega'(\xi_D\Phi_0-\Phi_0\xi_0)=(\alpha_1-\gamma_1,\,-\gamma_2,\,\alpha_3)$. On all of $\g(n)$ this linear map is onto $\R^{\creg(n,r)}$. It vanishes on $\ker\partial_{\beta(A)}$: indeed $\xi_D\Phi_0-\Phi_0\xi_0$ is the derivative at $t=0$ of $t\mapsto P(\exp(t\xi)\cdot\I(\beta(A)))$, so it equals $dP(\partial_{\beta(A)}\xi)$, where $dP$ is the differential of $P$ at $\I(\beta(A))$. Since $\g(n)=\g'_{\beta(A)}\oplus\ker\partial_{\beta(A)}$, the map is onto already on $\g'_{\beta(A)}$. Choose a decomposition $\g'_{\beta(A)}=\g'_1\oplus\g'_2$ such that $d\theta_0$ maps $\g'_1$ isomorphically onto $\R^{\creg(n,r)}$ and vanishes on $\g'_2$, and define
\[
\Theta(\xi,x)=\bigl(\omega(M(\xi,x)),\ \xi_2,\ x\bigr),\qquad\xi=\xi_1+\xi_2\text{ with }\xi_1\in\g'_1,\ \xi_2\in\g'_2 .
\]
We have $\omega(M(0,x))=\omega(\id_r\oplus X(x))=0$ for all $x$, and $\omega(M(\xi,0))=\theta(\xi)$. Hence $d\Theta_{(0,0)}(\xi,x)=(d\theta_0(\xi_1),\xi_2,x)$ is bijective, and $\Theta$ is a diffeomorphism from a neighbourhood of $(0,0)$ onto a neighbourhood of $0$. In the coordinates $(u,\xi_2,x)=\Theta(\xi,x)$, Lemma~\ref{lem:schur} and~\eqref{eq:M} give, near $0$,
\begin{equation}\label{eq:reduced}
\tfrac1{24}\bigl(\norm u^2+\Fr{X(x)}^2\bigr)\le L_{n,\Phi_0}\bigl(\Psi_{\beta(A)}(\Theta^{-1}(u,\xi_2,x))\bigr)\le6\bigl(\norm u^2+\Fr{X(x)}^2\bigr):
\end{equation}
by Lemma~\ref{lem:schur}(a), $L_{n,\Phi_0}=\frac{1}{2}\Fr{M-\Phi_0}^2$ is between $\frac{1}{6}$ and $\frac{3}{2}$ times $\Fr{\omega(M)}^2+\Fr{\Sch(M)}^2$, and by Lemma~\ref{lem:schur}(b) with $U=e^{\xi_D}$, $V=e^{-\xi_0}$, $\frac{1}{4}\Fr{X(x)}^2\le\Fr{\Sch(M)}^2\le4\Fr{X(x)}^2$. Moreover, for a small open neighbourhood $U_*$ of $(0,0)$,
\[
F_{n,\Phi_0}\cap\Psi_{\beta(A)}(U_*)=\Psi_{\beta(A)}\Theta^{-1}\bigl(\{0\}\times\g'_2\times Y_A\bigr)\cap\Psi_{\beta(A)}(U_*).
\]

\emph{Step 4: conclusion.} By~\eqref{eq:reduced} and Lemmas~\ref{lem:diffeo},~\ref{lem:compare},~\ref{lem:spectator} and~\ref{lem:gauss} (applied with $k=\creg(n,r)$), $L_{n,\Phi_0}$ has local volume order $(\frac{1}{2}\creg(n,r)+\lambda',m')$ at $\I(\beta(A))$ whenever $x\mapsto\Fr{X(x)}^2$ has local volume order $(\lambda',m')$ at $0$. By Step~2, this is the case if $L_{n-r,0}$ has local volume order $(\lambda',m')$ at $\I(h_A)$. For the codimension, by~\ref{H:lip} and the product formula~\ref{H:prod} below,
\[
d_{\I(\beta(A))}(F_{n,\Phi_0})=\dim\g'_2+d_0(Y_A),\qquad d_{\I(h_A)}(Z(n-r))=\dim\g'_{h_A}+d_0(Y_A).
\]
By Lemma~\ref{lem:ext}, $N(n)=\dim\g'_{\beta(A)}+\#\Lk(h_A)=\creg(n,r)+\dim\g'_2+\#\Lk(h_A)$ and $N(n-r)=\dim\g'_{h_A}+\#\Lk(h_A)$. Hence
\[
\codim_{\I(\beta(A))}F_{n,\Phi_0}=\creg(n,r)+\#\Lk(h_A)-d_0(Y_A)=\creg(n,r)+\codim_{\I(h_A)}Z(n-r).\qedhere
\]
\end{proof}

The product formula used above is the following standard fact~\cite[Product formulae~7.2 and~7.3]{falconer}.
\begin{enumerate}[label=(H\arabic*),leftmargin=3em,start=4]
\item\label{H:prod} For every nonempty Borel set $S\subseteq\R^b$ and every nonempty open ball (or open cube) $Q\subseteq\R^a$, $\dimH(Q\times S)=a+\dimH S$.
\end{enumerate}
Consequently, for a closed set $S\subseteq\R^b$ and a point $y\in S$, the local Hausdorff dimension at $(0,y)$ of $\R^a\times S$ is $a+d_y(S)$ (the sets $S\cap B_\delta(y)$ are nonempty Borel sets), and the local Hausdorff dimension at $(0,y)$ of $\{0\}\times S$ is $d_y(S)$. In Step~4 these facts are applied to the closed set $S=Y_A$ and the point $y=0\in Y_A$, after the changes of coordinates given by the charts, which are bi-Lipschitz on compact neighbourhoods (see~\ref{H:lip}).
\section{The zero target: a recursion over orbits}\label{sec:zero}

Throughout this section $D\in\N^*$ and $n\in\N^{D+1}$ is a width vector. Recall the function $K_n=\Fr P^2=2L_{n,0}$ on $\Rep(n)$ (Section~\ref{sec:chains}); by Lemma~\ref{lem:compare}, $K_n$ and $L_{n,0}$ have the same local volume orders. The zero fibre $Z(n)$ is the disjoint union of the orbits $O_h$, $h\in\Zb(n)$, by~\eqref{eq:orbit-ranks} and Lemma~\ref{lem:barcode-of-chain}(i): a chain $A$ with $P(A)=0$ has a barcode $\beta(A)$ without full-span slot. Since $K_n(g\cdot Y)=\Fr{g_DP(Y)g_0^{-1}}^2$ lies between two positive multiples of $K_n(Y)$, Lemmas~\ref{lem:diffeo} and~\ref{lem:compare} show that $K_n$ has the same local volume orders at all points of an orbit. For $h\in\Zb(n)$ we write $(\lambda_h,m_h)$ for the local volume order of $K_n$ at the points of $O_h$, when it exists.

In the rest of this section we fix $h\in\Zb(n)$, a total order on $\Sl(h)$, and the slice chart $\Psi_h$ of Lemma~\ref{lem:slice}, with domain $U_h$. Since $h$ is zero-output, every slot of $h$ is of exactly one of the kinds input, output, interior. A linked pair $(s,s')\in\Lk(h)$ is \emph{direct} if $s$ is an input slot and $s'$ an output slot, and \emph{residual} otherwise. We write $\Lk(h)=\Lk_{\mathrm{dir}}(h)\sqcup\Lk_{\mathrm{res}}(h)$ accordingly, and $x=(y,z)$ for the corresponding decomposition of $x\in\R^{\Lk(h)}=\R^{\Lk_{\mathrm{dir}}(h)}\times\R^{\Lk_{\mathrm{res}}(h)}$. Put $q_h=\tfrac12\#\Lk_{\mathrm{dir}}(h)$.

\subsection{The shear}
The rows of $P(\I(h)+Jx)$ are indexed by the output slots (the slots present at the vertex $D$) and its columns by the input slots (the slots present at the vertex $0$).

\begin{lemma}\label{lem:shear}
For an input slot $s$ and an output slot $s'$ of $h$,
\[
P(\I(h)+Jx)_{s',s}=[s\lin s']\,y_{(s,s')}+Q_{s',s}(z)\qquad\bigl(x=(y,z)\in\R^{\Lk(h)}\bigr),
\]
where $[s\lin s']$ is $1$ if $(s,s')\in\Lk(h)$ and $0$ otherwise, and $Q_{s',s}$ is a polynomial in the residual coordinates $z$ only: the sum in Lemma~\ref{lem:path} over the sequences of length $p\ge2$. Consequently the map $\mathrm{Sh}_h(y,z)=(y+Q_{\mathrm{dir}}(z),z)$, where $Q_{\mathrm{dir}}(z)=(Q_{s',s}(z))_{(s,s')\in\Lk_{\mathrm{dir}}(h)}$, is a polynomial automorphism of $\R^{\Lk(h)}$ with polynomial inverse and Jacobian determinant $1$, and
\[
K_n(\I(h)+Jx)=\norm{y+Q_{\mathrm{dir}}(z)}^2+H_h(z),\qquad H_h(z)=\sum_{(s,s')}Q_{s',s}(z)^2,
\]
the last sum being over the pairs (input slot $s$, output slot $s'$) that are not linked.
\end{lemma}
\begin{proof}
There are no full-span slots, so the term $p=0$ of Lemma~\ref{lem:path} does not occur, and the term $p=1$ is the coordinate of the linked pair $(s,s')$ if there is one. In a sequence $s=s_0\lin s_1\lin\cdots\lin s_p=s'$ with $p\ge2$, the slots $s_1,\dots,s_{p-1}$ are interior: they are not input slots because $b_{s_l}>b_s\ge0$, and not output slots because an output slot is never the first slot of a linked pair. So every pair of such a sequence starts or ends at an interior slot, i.e.\ is residual. The rest follows by expanding $\Fr{P(\I(h)+Jx)}^2$ entry by entry.
\end{proof}

By Lemma~\ref{lem:slice}, $K_n(\Psi_h(\xi,x))=\Fr{e^{\xi_D}P(\I(h)+Jx)e^{-\xi_0}}^2$, which lies between two positive multiples of $K_n(\I(h)+Jx)$ for $(\xi,x)$ near $(0,0)$. Applying Lemmas~\ref{lem:diffeo},~\ref{lem:compare},~\ref{lem:spectator} and~\ref{lem:gauss} at a point $(0,\mathrm{Sh}_h^{-1}(0,z_0))$ of the chart, we obtain:

\begin{corollary}\label{cor:H}
Let $z_0\in\R^{\Lk_{\mathrm{res}}(h)}$ with $H_h(z_0)=0$, and put
\[
p_h(z_0)=\I(h)+J\,\mathrm{Sh}_h^{-1}(0,z_0)=\I(h)+J\bigl(-Q_{\mathrm{dir}}(z_0),z_0\bigr).
\]
Then $p_h(z_0)\in Z(n)$. If $z_0$ is small enough that $(0,\mathrm{Sh}_h^{-1}(0,z_0))\in U_h$, then for all $(\lambda,m)\in\R\times\N$: $K_n$ has local volume order $(\lambda,m)$ at $p_h(z_0)$ if and only if $H_h$ has local volume order $(\lambda-q_h,m)$ at $z_0$. This includes $z_0=0$, for which $p_h(0)=\I(h)$.
\end{corollary}
Here the case $\Lk_{\mathrm{res}}(h)=\emptyset$ is included: then $H_h$ is the zero function on the point $\R^0$, with local volume order $(0,1)$.

\subsection{Potentials and a dilation}
\begin{definition}\label{def:potential}
A (rational) \emph{admissible potential} for $h$ is a function $u:\Sl(h)\to\Q$ with $u_s=0$ for the input slots $s$, $u_s=1$ for the output slots $s$, and $u_s<u_{s'}$ for every linked pair $(s,s')\in\Lk(h)$. Its \emph{threshold} is
\[
T(u)=\tfrac12\sum_{(s,s')\in\Lk(h)}(u_{s'}-u_s).
\]
\end{definition}
Let $\nu(b,e)=b/(b+D-e)$ for integers $0\le b\le e\le D$ with $(b,e)\neq(0,D)$. The \emph{canonical potential} is $u_s=\nu(b_s,e_s)$; its denominator is positive because $h$ has no full-span slot. It is admissible: it is $0$ on input slots and $1$ on output slots, and the function $\nu$ is nondecreasing in $b$, and nondecreasing in $e$, strictly when $b>0$. If $(s,s')\in\Lk(h)$ and $b_s>0$, then $u_s=\nu(b_s,e_s)<\nu(b_s,e_{s'})\le\nu(b_{s'},e_{s'})=u_{s'}$, since $e_s<e_{s'}$ and $b_s<b_{s'}$; if $b_s=0$, then $u_s=0<u_{s'}$ because $b_{s'}>0$. The canonical potential is the one used in the formal definition of the program (Appendix~\ref{app:lean}).

\begin{lemma}[Dilation]\label{lem:dilation}
Let $u$ be an admissible potential for $h$, and $M\in\N^*$ with $Mu_s\in\Z$ for all $s\in\Sl(h)$. For $\pi=(s,s')\in\Lk(h)$ put $w_\pi=M(u_{s'}-u_s)\in\N^*$. For a real number $r>0$, let $\Gamma(r)\in G(n)$ act at each vertex $v$ by $\epsilon^v_s\mapsto r^{Mu_s}\epsilon^v_s$ ($s\in\Sl_v(h)$), and let $\Delta_r$ be the linear automorphism of $\R^{\Lk(h)}$ with $(\Delta_rx)_\pi=r^{w_\pi}x_\pi$. Then for all $x\in\R^{\Lk(h)}$ and $z\in\R^{\Lk_{\mathrm{res}}(h)}$:
\begin{gather*}
\Gamma(r)\cdot(\I(h)+Jx)=\I(h)+J\Delta_rx,\qquad P(\I(h)+J\Delta_rx)=r^M\,P(\I(h)+Jx),\\
Q_{\mathrm{dir}}(\Delta_rz)=r^M\,Q_{\mathrm{dir}}(z),\qquad H_h(\Delta_rz)=r^{2M}H_h(z).
\end{gather*}
Here $\Delta_r$ also denotes its restriction to the coordinate subspace $\R^{\Lk_{\mathrm{res}}(h)}$, which it preserves.
\end{lemma}
\begin{proof}
The first identity is checked arrow by arrow, as in the proof of Lemma~\ref{lem:strict}: $\Gamma(r)\cdot\I(h)=\I(h)$, and $\Gamma(r)\cdot J_\pi=r^{w_\pi}J_\pi$. Every slot present at the vertex $0$ is an input slot and every slot present at the vertex $D$ is an output slot, so $\Gamma(r)_0=\id$ and $\Gamma(r)_D=r^M\id$. Hence $P(\Gamma(r)\cdot Y)=r^MP(Y)$ for every chain $Y$, which gives the second identity. Direct pairs have weight $w_\pi=M$. Comparing the entries of Lemma~\ref{lem:shear} on both sides of the second identity, as polynomial identities in $x$, gives the other two.
\end{proof}

\begin{lemma}[Zeros of $H_h$ and higher orbits]\label{lem:higher}
\begin{enumerate}[label=\textup{(\alph*)}]
\item For every $z\neq0$ with $H_h(z)=0$, the point $p_h(z)$ of Corollary~\ref{cor:H} lies in an orbit $O_k$ with $k\in\Zb(n)$, $h\sqsubseteq k$ and $k\neq h$, and so do the points $p_h(\Delta_rz)$ for all $r>0$.
\item Conversely, if $k\in\Zb(n)$, $h\sqsubseteq k$ and $k\neq h$, then for every $\eta>0$ there is $z\in\R^{\Lk_{\mathrm{res}}(h)}$ with $0<\norm z<\eta$, $H_h(z)=0$ and $p_h(z)\in O_k$.
\end{enumerate}
\end{lemma}
\begin{proof}
(a) $p_h(z)\in Z(n)$ by Corollary~\ref{cor:H}, so $p_h(z)\in O_k$ for some $k\in\Zb(n)$. By Lemma~\ref{lem:dilation}, $\Gamma(r)\cdot p_h(z)=p_h(\Delta_rz)$, so all the points $p_h(\Delta_rz)$ lie in $O_k$, and $p_h(\Delta_rz)\to\I(h)$ as $r\to0$. Hence $\I(h)\in\overline{O_k}$, i.e.\ $h\sqsubseteq k$. Finally $k\neq h$: for small $r$, the point $p_h(\Delta_rz)=\I(h)+Jx$ has $x\neq0$ small, so it is not in $O_h$ by Lemma~\ref{lem:slice}.

(b) Since $\I(h)\in\overline{O_k}$ and $\Psi_h(U_h)$ is a neighbourhood of $\I(h)$, there are points $\Psi_h(\xi,x)\in O_k$ with $(\xi,x)\in U_h$ arbitrarily small. Then $\I(h)+Jx=(\exp(-\xi))\cdot\Psi_h(\xi,x)\in O_k\subseteq Z(n)$. By Lemma~\ref{lem:shear}, $K_n(\I(h)+Jx)=0$ forces $y=-Q_{\mathrm{dir}}(z)$ and $H_h(z)=0$, so $\I(h)+Jx=p_h(z)$. If $z=0$, then $x=0$ and $p_h(z)=\I(h)\notin O_k$; hence $z\neq0$.
\end{proof}

\subsection{The recursion}
\begin{proposition}\label{prop:recursion}
Suppose that $K_n$ has a local volume order $(\lambda_k,m_k)$ along every orbit $O_k$ with $k\in\Zb(n)$, $h\sqsubseteq k$ and $k\neq h$. Put
\begin{gather*}
a_*=\min\{\lambda_k:\ k\in\Zb(n),\ h\sqsubseteq k,\ k\neq h\},\\
m_*=\max\{m_k:\ k\in\Zb(n),\ h\sqsubseteq k,\ k\neq h,\ \lambda_k=a_*\},
\end{gather*}
with $a_*=+\infty$ (and $m_*$ undefined) if there is no such $k$. Then for every admissible potential $u$ for $h$, $K_n$ has at $\I(h)$ the local volume order $(\lambda_h,m_h)$ given by
\[
\lambda_h=\min\bigl(T(u),a_*\bigr),\qquad m_h=\begin{cases}1&\text{if }T(u)<a_*,\\ m_*&\text{if }T(u)>a_*,\\ m_*+1&\text{if }T(u)=a_*.\end{cases}
\]
\end{proposition}

\begin{proof}
Fix $u$, an integer $M$ as in Lemma~\ref{lem:dilation}, the weights $w_\pi$ and the dilations $\Delta_r$. Each direct pair contributes $\frac{1}{2}$ to $T(u)$, so
\[
T(u)=q_h+T_H,\qquad T_H=\tfrac12\sum_{(s,s')\in\Lk_{\mathrm{res}}(h)}(u_{s'}-u_s)=\frac{\abs w}{2M},\quad\text{where }\abs w=\sum_{\pi\in\Lk_{\mathrm{res}}(h)}w_\pi .
\]
By Corollary~\ref{cor:H} with $z_0=0$, it suffices to prove that $H_h$ has at $0\in\R^{\Lk_{\mathrm{res}}(h)}$ the local volume order $(\min(T_H,a_*-q_h),m_h)$, with $m_h$ as in the statement. If $\Lk_{\mathrm{res}}(h)=\emptyset$, then $H_h$ is the zero function on a point, with order $(0,1)$; in that case Lemma~\ref{lem:higher}(b) shows that no $k$ as in the statement exists, so $a_*=+\infty$, and the claim reads $(T_H,1)=(0,1)$. Assume from now on that $d:=\#\Lk_{\mathrm{res}}(h)\ge1$, and put $a=a_*-q_h\in\,]-\infty,+\infty]$.

\emph{Step 1: weighted polar coordinates.} Let $M_0\in\N^*$ be a common multiple of the weights $w_\pi$, $\pi\in\Lk_{\mathrm{res}}(h)$, and
\[
\varrho(z)=\sum_{\pi\in\Lk_{\mathrm{res}}(h)}z_\pi^{2M_0/w_\pi}\qquad\bigl(z\in\R^{\Lk_{\mathrm{res}}(h)}\bigr).
\]
This is a polynomial, positive on $\R^{\Lk_{\mathrm{res}}(h)}\setminus\{0\}$, with $\varrho(\Delta_rz)=r^{2M_0}\varrho(z)$. The set $\mathbb S=\varrho^{-1}(1)$ is compact. The vector field $\mathcal E(z)=(w_\pi z_\pi)_\pi$ generates the dilations ($\frac{d}{dr}\Delta_rz=r^{-1}\mathcal E(\Delta_rz)$), and Euler's identity gives $d\varrho_z(\mathcal E(z))=2M_0\varrho(z)$. Hence $\mathbb S$ is a smooth hypersurface transverse to $\mathcal E$, and
\[
\Pi:\,]0,+\infty[\,\times\mathbb S\to\R^{\Lk_{\mathrm{res}}(h)}\setminus\{0\},\qquad(r,\vartheta)\mapsto\Delta_r\vartheta,
\]
is a diffeomorphism, with inverse $z\mapsto(\varrho(z)^{1/2M_0},\Delta_{\varrho(z)^{-1/2M_0}}z)$. For a Borel set $S\subseteq\mathbb S$ put $\sigma(S)=\vol\,\Pi(]0,1[\,\times S)$. Since $\Delta_r$ is linear with determinant $r^{\abs w}$, we get $\vol\,\Pi(]0,r_0[\,\times S)=r_0^{\abs w}\sigma(S)$ for $r_0>0$. Hence $\Pi^{-1}$ carries Lebesgue measure to the product of the measure $\abs w\,r^{\abs w-1}dr$ on $]0,+\infty[$ with $\sigma$, by uniqueness of measures that agree on products. In a chart $\chi:V_\chi\to\mathbb S$ ($V_\chi\subseteq\R^{d-1}$ open), the map $(r,v)\mapsto\Delta_r\chi(v)$ has Jacobian determinant $r^{\abs w-1}j(v)$, with $j(v)=\det(\mathcal E(\chi(v)),d\chi(v))$ continuous and nowhere zero (by transversality). Therefore $\sigma=\abs w^{-1}\abs{j(v)}\,dv$ on $\chi(V_\chi)$: the measure $\sigma$ is finite and has a continuous positive density in every chart.

\emph{Step 2: the angular function.} For $\eta>0$ put $\Xi(\eta)=\sigma\{\vartheta\in\mathbb S:\ H_h(\vartheta)<\eta\}$. For $r_0>0$ consider the window $W_{r_0}=\{z:\ \varrho(z)<r_0^{2M_0}\}=\Pi(]0,r_0[\,\times\mathbb S)\cup\{0\}$. The polar decomposition and $H_h(\Delta_r\vartheta)=r^{2M}H_h(\vartheta)$ give
\begin{equation}\label{eq:polar}
V_{r_0}(\eps):=\vol\{z\in W_{r_0}:\ H_h(z)<\eps\}=\abs w\int_0^{r_0}r^{\abs w-1}\,\Xi(\eps r^{-2M})\,dr .
\end{equation}

\emph{Step 3: angular asymptotics.} Let $Z_{\mathbb S}=\{\vartheta\in\mathbb S:\ H_h(\vartheta)=0\}$, a compact set.

\emph{Case $Z_{\mathbb S}=\emptyset$.} Then $0<\eta_1\le H_h\le\eta_2$ on $\mathbb S$ for some real numbers $\eta_1,\eta_2$, so $\Xi(\eta)=0$ for $\eta\le\eta_1$ and $\Xi(\eta)=\sigma(\mathbb S)$ for $\eta>\eta_2$. No $k$ as in the statement exists, so $a_*=+\infty$: otherwise Lemma~\ref{lem:higher}(b) would give a zero $z\neq0$ of $H_h$, and then $\Delta_{\varrho(z)^{-1/2M_0}}z\in Z_{\mathbb S}$ by Lemma~\ref{lem:dilation}.

\emph{Case $Z_{\mathbb S}\neq\emptyset$.} Let $\vartheta\in Z_{\mathbb S}$. For small $r>0$, $\Delta_r\vartheta$ is a small zero of $H_h$, and $p_h(\Delta_r\vartheta)\in O_k$ for some $k\in\Zb(n)$ with $h\sqsubseteq k\neq h$, independent of $r$ (Lemma~\ref{lem:higher}(a)). By Corollary~\ref{cor:H} and Lemma~\ref{lem:homog-order} (applied to the linear map $\Delta_r$ and the function $H_h$, with the constant $r^{2M}$), $H_h$ has local volume order $(\lambda_k-q_h,m_k)$ at $\Delta_r\vartheta$, hence at $\vartheta$. Near $\vartheta$, use the coordinates $(r,v)\mapsto\Delta_r\chi(v)$ around $(1,v_\vartheta)$, where $\chi(v_\vartheta)=\vartheta$. In them $H_h(\Delta_r\chi(v))=r^{2M}H_h(\chi(v))$, which lies between two positive multiples of $H_h(\chi(v))$ for $r\in\,]\frac{1}{2},2[$. By Lemmas~\ref{lem:diffeo},~\ref{lem:compare} and~\ref{lem:spectator}, the function $v\mapsto H_h(\chi(v))$ has local volume order $(\lambda_k-q_h,m_k)$ at $v_\vartheta$. By Lemma~\ref{lem:prior}, this holds for the measure $\sigma$ (transported by $\chi$) as well, since $\sigma$ has a continuous positive density.

By compactness, finitely many chart balls $U_1,\dots,U_N$ cover $Z_{\mathbb S}$, each centred at a zero $\vartheta_i$ with associated barcode $k_i$, and each of radius below the radius $\delta_0$ of the corresponding local volume order. On the compact set $\mathbb S\setminus\bigcup_iU_i$ we have $H_h\ge\eta_3$ for some $\eta_3>0$. Hence for $\eta<\eta_3$
\[
\Xi(\eta)\le\sum_{i=1}^N\sigma\{\vartheta\in U_i:\ H_h(\vartheta)<\eta\}\le C\,\eta^{a}L(\eta)^{m_*-1}\qquad(\eta\text{ small}),
\]
Indeed, $\lambda_{k_i}-q_h\ge a$. If $\lambda_{k_i}-q_h=a$, then $m_{k_i}\le m_*$ and $L(\eta)>1$ for $\eta<e^{-1}$. If $\lambda_{k_i}-q_h>a$, then
\[
\frac{\eta^{\lambda_{k_i}-q_h}L(\eta)^{m_{k_i}-1}}{\eta^aL(\eta)^{m_*-1}}=\eta^{\lambda_{k_i}-q_h-a}L(\eta)^{m_{k_i}-m_*}\longrightarrow0\qquad(\eta\to0).
\]
Since the cover is finite, a single positive constant and a single positive cutoff suffice for the displayed upper bound. Conversely, by Lemma~\ref{lem:higher}(b) and Lemma~\ref{lem:dilation}, every orbit $O_k$ with $k\in\Zb(n)$, $h\sqsubseteq k\neq h$, contains $p_h(\vartheta)$ for some $\vartheta\in Z_{\mathbb S}$: take $z$ from Lemma~\ref{lem:higher}(b) and $\vartheta=\Delta_{\varrho(z)^{-1/2M_0}}z$. The local volume order at that $\vartheta$ has exponent $\lambda_k-q_h\ge0$, by Definition~\ref{def:order}. This applies to every $k$ in the minimum defining $a_*$; in particular $a=a_*-q_h\ge0$. Choosing $k$ with $(\lambda_k,m_k)=(a_*,m_*)$ and a small chart ball around the corresponding $\vartheta$ gives the matching lower bound. Thus there is $\eta_0\in\,]0,e^{-1}[$ with
\[
\Xi(\eta)\asymp\eta^aL(\eta)^{m_*-1}\ \text{ on }]0,\eta_0[,\qquad\Xi(\eta_0)\le\Xi(\eta)\le\sigma(\mathbb S)\ \text{ for }\eta\ge\eta_0,
\]
and $\Xi(\eta_0)>0$.

\emph{Step 4: radial integration.} First suppose $Z_{\mathbb S}\neq\emptyset$. Fix $r_0>0$ and split~\eqref{eq:polar} at $r_\eps=(\eps/\eta_0)^{1/2M}$, which is $<r_0$ for small $\eps$. For $r<r_\eps$ we have $\eps r^{-2M}>\eta_0$, so the inner part satisfies
\[
\abs w\int_0^{r_\eps}r^{\abs w-1}\,\Xi(\eps r^{-2M})\,dr\asymp r_\eps^{\abs w}\asymp\eps^{T_H}.
\]
In the outer part substitute $r=(\eps e^\tau)^{1/2M}$, so that $\eps r^{-2M}=e^{-\tau}$ and $dr=r\,d\tau/2M$:
\[
\abs w\int_{r_\eps}^{r_0}r^{\abs w-1}\,\Xi(\eps r^{-2M})\,dr\asymp\eps^{T_H}\int_{\tau_0}^{\tau_\eps}e^{(T_H-a)\tau}\,\tau^{m_*-1}\,d\tau,
\]
with $\tau_0=\log(1/\eta_0)>1$ and $\tau_\eps=\log(r_0^{2M}/\eps)\to+\infty$. If $T_H>a$ and $\tau_\eps\ge\max(\tau_0+1,2)$, then
\begin{align*}
e^{(T_H-a)(\tau_\eps-1)}(\tau_\eps-1)^{m_*-1}
&\le\int_{\tau_0}^{\tau_\eps}e^{(T_H-a)\tau}\tau^{m_*-1}\,d\tau\\
&\le\frac{e^{(T_H-a)\tau_\eps}\tau_\eps^{m_*-1}}{T_H-a}.
\end{align*}
The lower bound uses the interval $[\tau_\eps-1,\tau_\eps]$; the upper bound uses $\tau^{m_*-1}\le\tau_\eps^{m_*-1}$. Since $\tau_\eps-1\ge\tau_\eps/2$, these inequalities give an integral $\asymp e^{(T_H-a)\tau_\eps}\tau_\eps^{m_*-1}$. If $T_H=a$, the integral equals $(\tau_\eps^{m_*}-\tau_0^{m_*})/m_*$ and is $\asymp\tau_\eps^{m_*}$. If $T_H<a$, the integral converges to a strictly positive finite number: it is bounded below by its integral over $[\tau_0,\tau_0+1]$, and for sufficiently large $\tau$ its integrand is at most $e^{(T_H-a)\tau/2}$, because $\tau^{m_*-1}\le e^{(a-T_H)\tau/2}$. Thus it is $\asymp1$. Since $\tau_\eps=L(\eps)+2M\log r_0$:
\begin{itemize}
\item if $T_H>a$: $V_{r_0}(\eps)\asymp\eps^{T_H}+\eps^aL(\eps)^{m_*-1}\asymp\eps^aL(\eps)^{m_*-1}$;
\item if $T_H=a$: $V_{r_0}(\eps)\asymp\eps^{T_H}+\eps^aL(\eps)^{m_*}\asymp\eps^aL(\eps)^{m_*}$;
\item if $T_H<a$: $V_{r_0}(\eps)\asymp\eps^{T_H}$.
\end{itemize}
In the case $Z_{\mathbb S}=\emptyset$, $\Xi(\eps r^{-2M})$ vanishes for $r>(\eps/\eta_1)^{1/2M}$ and equals $\sigma(\mathbb S)$ for $r<(\eps/\eta_2)^{1/2M}$, so again $V_{r_0}(\eps)\asymp\eps^{T_H}$. In all cases $H_h$ has the claimed local volume order with respect to the window system $(W_{r_0})_{r_0>0}$, for Lebesgue measure and the prior $1$. This window system is equivalent to the ball system of the Euclidean norm: the sets $W_{r_0}=\Delta_{r_0}(W_1)$ are bounded open neighbourhoods of $0$ that shrink to $0$ as $r_0\to0$, and $B_\delta(0)\subseteq W_{r_0}$ as soon as $r_0^{2M_0}>\max\{\varrho(z):\ \norm z\le\delta\}$. Lemma~\ref{lem:sandwich} concludes. Adding $q_h$ (Corollary~\ref{cor:H}) gives $\lambda_h=\min(q_h+T_H,\,q_h+a)=\min(T(u),a_*)$ and the stated $m_h$.
\end{proof}

\begin{corollary}[Existence]\label{cor:existence}
For every $h\in\Zb(n)$, $K_n$ has a local volume order $(\lambda_h,m_h)$ at the points of $O_h$. Moreover $\lambda_h\le T(u)$ for every admissible potential $u$ for $h$, and $\lambda_h\le\lambda_k$ for every $k\in\Zb(n)$ with $h\sqsubseteq k$.
\end{corollary}
\begin{proof}
$(\Zb(n),\sqsubseteq)$ is a finite partially ordered set (Proposition~\ref{prop:closure}(a)). We argue by descending induction: for a maximal element, Proposition~\ref{prop:recursion} applies with $a_*=+\infty$; in general it applies once all strictly larger elements have been treated. The inequalities follow from $\lambda_h=\min(T(u),a_*)$.
\end{proof}
\section{The exponent and the codimension}\label{sec:exponent}

We keep the notation of Section~\ref{sec:zero}: $D\in\N^*$, $n\in\N^{D+1}$, and for $h\in\Zb(n)$, $(\lambda_h,m_h)$ is the local volume order of $K_n=\Fr P^2$ at the points of $O_h$, which exists by Corollary~\ref{cor:existence}. In this section the order $\sqsubseteq$ of Definition~\ref{def:geom-order} is used, and for $h\in\Zb(n)$ we put
\[
C^\sqsubseteq(h)=\min\{\cout(k):\ k\in\Zb(n),\ h\sqsubseteq k,\ k\text{ balanced}\}\in\N\cup\{+\infty\},
\]
with $\min\emptyset=+\infty$. By Proposition~\ref{prop:closure}, $C^\sqsubseteq=C$, but the proof of Theorem~\ref{thm:A} does not need this.

\subsection{Stationarity}
For a slot $s$ of $h$ let $\mathrm{in}(s)$ and $\mathrm{out}(s)$ be as in Definition~\ref{def:links}. Input slots are never the second slot of a linked pair, and output slots never the first one. Hence, for every admissible potential $u$ for $h$ (Definition~\ref{def:potential}),
\begin{equation}\label{eq:telescope}
2T(u)=\sum_{(s,s')\in\Lk(h)}(u_{s'}-u_s)=\sum_{s\in\Sl(h)}\bigl(\mathrm{in}(s)-\mathrm{out}(s)\bigr)u_s=\cout(h)+\sum_{s\text{ interior}}\bigl(\mathrm{in}(s)-\mathrm{out}(s)\bigr)u_s ,
\end{equation}
using $u_s=0$ on input slots and $u_s=1$ on output slots. In particular, if $h$ is balanced then $2T(u)=\cout(h)$ for every admissible potential $u$.

\begin{lemma}[Stationarity]\label{lem:stationarity}
Let $h\in\Zb(n)$. If $\lambda_h=T(u_0)$ for some admissible potential $u_0$ for $h$, then $h$ is balanced and $2\lambda_h=\cout(h)$. Equivalently: if $h$ is not balanced, then $\lambda_h<T(u)$ for every admissible potential $u$ for $h$.
\end{lemma}
\begin{proof}
By Corollary~\ref{cor:existence}, $\lambda_h\le T(u)$ for every admissible potential $u$. Suppose $\lambda_h=T(u_0)$. Let $\Omega$ be the set of real functions $u$ on $\Sl(h)$ with $u_s=0$ on input slots, $u_s=1$ on output slots, and $u_s<u_{s'}$ for every $(s,s')\in\Lk(h)$. It is an open subset of an affine space (the values on interior slots are free), the admissible potentials are its rational points, and $u_0\in\Omega$. If some interior slot $s$ had $\mathrm{in}(s)\neq\mathrm{out}(s)$, then by~\eqref{eq:telescope} a small rational change of $u_0$ at $s$ alone, in the appropriate direction, would give an admissible potential $u$ with $T(u)<T(u_0)=\lambda_h$, a contradiction. Hence $h$ is balanced, and $2\lambda_h=2T(u_0)=\cout(h)$ by~\eqref{eq:telescope}.
\end{proof}

\subsection{The exponent}
\begin{theorem}\label{thm:coefficient}
For every $h\in\Zb(n)$, the set minimized in $C^\sqsubseteq(h)$ is nonempty, and $2\lambda_h=C^\sqsubseteq(h)$.
\end{theorem}
\begin{proof}
Descending induction on the finite partially ordered set $(\Zb(n),\sqsubseteq)$. Fix $h$ and assume the claim for all $k\in\Zb(n)$ with $h\sqsubseteq k\neq h$. Let $a_*$ be as in Proposition~\ref{prop:recursion}. The balanced elements $k''$ with $h\sqsubseteq k''\neq h$ are exactly the balanced elements above some $k$ with $h\sqsubseteq k\neq h$ (take $k=k''$). So, by the induction hypothesis,
\[
\min\{\cout(k''):\ k''\in\Zb(n),\ h\sqsubseteq k''\neq h,\ k''\text{ balanced}\}=\min_{h\sqsubseteq k\neq h}C^\sqsubseteq(k)=\min_{h\sqsubseteq k\neq h}2\lambda_k=2a_*,
\]
with the convention $\min\emptyset=+\infty$ on both sides. Fix an admissible potential $u$ for $h$; by Proposition~\ref{prop:recursion}, $\lambda_h=\min(T(u),a_*)$.

If $\lambda_h<a_*$, then $\lambda_h=T(u)$, so $h$ is balanced with $\cout(h)=2\lambda_h$ (Lemma~\ref{lem:stationarity}). Hence $h$ itself belongs to the minimized set, and $C^\sqsubseteq(h)=\min(2\lambda_h,2a_*)=2\lambda_h$.

If $\lambda_h=a_*$, then $a_*<+\infty$, and either $h$ is not balanced, or it is balanced with $\cout(h)=2T(u)\ge2\lambda_h=2a_*$. In both cases the minimized set is nonempty and $C^\sqsubseteq(h)=2a_*=2\lambda_h$.

Maximal elements fall under the first case, since $a_*=+\infty$ for them.
\end{proof}

\subsection{The codimension}
\begin{proposition}[Output cut]\label{prop:codim}
For every $h\in\Zb(n)$ and every $X\in O_h$,
\[
\codim_XZ(n)=\min\{\#\Lk(k):\ k\in\Zb(n),\ h\sqsubseteq k\}=C^\sqsubseteq(h).
\]
\end{proposition}
\begin{proof}
\emph{The first equality.} $Z(n)$ is the finite disjoint union of the embedded submanifolds $O_k$, $k\in\Zb(n)$ (Lemma~\ref{lem:orbit}). We have $X\in\overline{O_k}$ if and only if $O_h\subseteq\overline{O_k}$, i.e.\ $h\sqsubseteq k$, since $\overline{O_k}$ is $G(n)$-invariant. By~\eqref{eq:localdim-strata}, $d_X(Z(n))=\max\{\dim O_k:\ k\in\Zb(n),\ h\sqsubseteq k\}$, and Lemma~\ref{lem:ext} gives the first equality. In particular $\codim_XZ(n)\le\codim_{X'}Z(n)$ for $X'\in O_k$ with $h\sqsubseteq k$.

\emph{Upper bound by $C^\sqsubseteq(h)$.} Let $k\in\Zb(n)$, and let $\Lk_{\to\mathrm{out}}(k)\subseteq\Lk(k)$ be the set of linked pairs $(s,s')$ with $s'$ an output slot; $\#\Lk_{\to\mathrm{out}}(k)=\cout(k)$. Every sequence of linked pairs from an input slot to an output slot ends with a pair in $\Lk_{\to\mathrm{out}}(k)$. By Lemma~\ref{lem:path} (applied to $k$, which has no full-span slot), $P(\I(k)+Jx)=0$ whenever $x_\pi=0$ for all $\pi\in\Lk_{\to\mathrm{out}}(k)$. Since $Z(n)$ is $G(n)$-invariant, the slice chart $\Psi_k$ maps the intersection of its domain with the subspace $\{(\xi,x):\ x_\pi=0\text{ for all }\pi\in\Lk_{\to\mathrm{out}}(k)\}$ into $Z(n)$. The image is an embedded submanifold of codimension $\cout(k)$ through $\I(k)$. By~\ref{H:mono} and~\ref{H:mfd}, $\codim_{\I(k)}Z(n)\le\cout(k)$. With the monotonicity above, $\codim_XZ(n)\le\cout(k)$ for every $k\in\Zb(n)$ with $h\sqsubseteq k$, in particular for the balanced ones: $\codim_XZ(n)\le C^\sqsubseteq(h)$.

\emph{Lower bound.} Choose $k^*\in\Zb(n)$ with $h\sqsubseteq k^*$ and $\#\Lk(k^*)$ minimal, so that $\codim_XZ(n)=\#\Lk(k^*)$. Then $k^*$ is maximal in $(\Zb(n),\sqsubseteq)$: if $k^*\sqsubseteq k''\neq k^*$, Lemma~\ref{lem:strict} would give $\#\Lk(k'')<\#\Lk(k^*)$ with $h\sqsubseteq k''$. Suppose $k^*$ had a residual linked pair $\pi$. By Lemma~\ref{lem:path} (applied to $k^*$, which has no full-span slot, so that sequences with $p=0$ do not occur), the only sequence of linked pairs that uses no pair other than $\pi$ is $\pi$ itself, and it does not go from an input slot to an output slot. Hence $P(\I(k^*)+\tau J_\pi)=0$, i.e.\ $\I(k^*)+\tau J_\pi\in Z(n)$, for every real $\tau$. By Lemma~\ref{lem:switch}, $\I(k^*)+\tau J_\pi\in O_{k''}$ for every $\tau\neq0$, where $k''=(k^*)^\pi$; so $O_{k''}\subseteq Z(n)$, i.e.\ $k''\in\Zb(n)$, and $k^*\sqsubseteq k''$ (Lemma~\ref{lem:switch}). Moreover $k''\neq k^*$, since $\rho_{b_s,e_{s'}}(k'')=\rho_{b_s,e_{s'}}(k^*)+1$ for $\pi=(s,s')$. This contradicts maximality. So every linked pair of $k^*$ is direct. Then no interior slot of $k^*$ belongs to a linked pair, so $k^*$ is balanced, and $\cout(k^*)=\#\Lk(k^*)$. Therefore $\codim_XZ(n)=\cout(k^*)\ge C^\sqsubseteq(h)$.
\end{proof}

Combining Theorem~\ref{thm:coefficient} and Proposition~\ref{prop:codim}:
\begin{corollary}\label{cor:zero-A}
For every $h\in\Zb(n)$ and every $X\in O_h$, $K_n$ has local volume order $(\lambda_h,m_h)$ at $X$, and $2\lambda_h=\codim_XZ(n)=C^\sqsubseteq(h)$.
\end{corollary}
\section{The multiplicity}\label{sec:multiplicity}

We keep the notation of Sections~\ref{sec:zero} and~\ref{sec:exponent}. Call $h\in\Zb(n)$ \emph{$\sqsubseteq$-critical} if it is balanced and $\cout(h)=C^\sqsubseteq(h)$. Let $F^\sqsubseteq(h)$ be the maximum of $\#\{i\in\ii0{\ell-1}:\ h_i\text{ is }\sqsubseteq\text{-critical}\}$ over all $\ell\in\N$ and all finite sequences $h=h_0\sqsubseteq h_1\sqsubseteq\cdots\sqsubseteq h_\ell$ in $\Zb(n)$ with $h_i\neq h_{i+1}$ and $C^\sqsubseteq(h_i)=C^\sqsubseteq(h)$ for all $i$. By Proposition~\ref{prop:closure}, these notions agree with ``critical'' and $F$ of Definition~\ref{def:invariants}.

\begin{lemma}\label{lem:critical}
Let $h\in\Zb(n)$ and let $u$ be an admissible potential for $h$. Then $h$ is $\sqsubseteq$-critical if and only if $\lambda_h=T(u)$.
\end{lemma}
\begin{proof}
If $\lambda_h=T(u)$, then $h$ is balanced and $\cout(h)=2\lambda_h$ by Lemma~\ref{lem:stationarity}, and $2\lambda_h=C^\sqsubseteq(h)$ by Theorem~\ref{thm:coefficient}. Conversely, if $h$ is balanced, then $2T(u)=\cout(h)$ by~\eqref{eq:telescope}; if moreover $\cout(h)=C^\sqsubseteq(h)=2\lambda_h$, then $T(u)=\lambda_h$.
\end{proof}

\begin{theorem}\label{thm:mult}
For every $h\in\Zb(n)$, $m_h=1+F^\sqsubseteq(h)$.
\end{theorem}
\begin{proof}
Let $\mathrm{Succ}(h)=\{k\in\Zb(n):\ h\sqsubseteq k\neq h,\ C^\sqsubseteq(k)=C^\sqsubseteq(h)\}$. Splitting off the first step of a sequence gives
\begin{equation}\label{eq:F-rec}
F^\sqsubseteq(h)=\max\Bigl(0,\ \max_{k\in\mathrm{Succ}(h)}\bigl([h\text{ is }\sqsubseteq\text{-critical}]+F^\sqsubseteq(k)\bigr)\Bigr),
\end{equation}
where $\max\emptyset=0$ and $[\cdot]\in\{0,1\}$ is the truth value. We argue by descending induction on $(\Zb(n),\sqsubseteq)$, and compare~\eqref{eq:F-rec} with Proposition~\ref{prop:recursion} for a fixed admissible potential $u$. By Theorem~\ref{thm:coefficient}, $k\in\mathrm{Succ}(h)$ if and only if $h\sqsubseteq k\neq h$ and $\lambda_k=\lambda_h$. Since $\lambda_h=\min(T(u),a_*)\le a_*\le\lambda_k$ for all such $k$, the set $\mathrm{Succ}(h)$ is nonempty exactly when $\lambda_h=a_*$, and then it consists of the $k$ with $h\sqsubseteq k\neq h$ and $\lambda_k=a_*$.
\begin{itemize}
\item If $T(u)<a_*$: $h$ is $\sqsubseteq$-critical (Lemma~\ref{lem:critical}) and $\mathrm{Succ}(h)=\emptyset$, so $F^\sqsubseteq(h)=0$ and $m_h=1$.
\item If $T(u)>a_*$: $h$ is not $\sqsubseteq$-critical, and by the induction hypothesis
\[
m_h=m_*=\max_{k\in\mathrm{Succ}(h)}m_k=1+\max_{k\in\mathrm{Succ}(h)}F^\sqsubseteq(k)=1+F^\sqsubseteq(h).
\]
\item If $T(u)=a_*$: $h$ is $\sqsubseteq$-critical, $\mathrm{Succ}(h)\neq\emptyset$, and
\[
m_h=m_*+1=1+\max_{k\in\mathrm{Succ}(h)}\bigl(1+F^\sqsubseteq(k)\bigr)=1+F^\sqsubseteq(h).\qedhere
\]
\end{itemize}
\end{proof}

\begin{remark}[Critical strata]\label{rem:critical-geom}
By Corollary~\ref{cor:zero-A} and Proposition~\ref{prop:codim}, $C^\sqsubseteq(k)$ is the local codimension of $Z(n)$ at the points of $O_k$. The sequences in the definition of $F^\sqsubseteq(h)$ are therefore chains of orbits $O_h\subseteq\overline{O_{h_1}}$, $O_{h_1}\subseteq\overline{O_{h_2}},\dots$ along which the local codimension of $Z(n)$ is constant. An orbit $O_k$ is critical if and only if the polynomial $H_k$ of Section~\ref{sec:zero} has local volume exponent exactly $T_H$, the exponent ``expected'' from the weights; equivalently, the radial integral in the proof of Proposition~\ref{prop:recursion} contributes at the leading order. Each critical orbit along a chain contributes one factor $L(\eps)$. Combinatorial formulas of a similar kind, for the real log canonical threshold and its multiplicity, are known for real hyperplane arrangements~\cite{kosta-windisch}.
\end{remark}

\begin{remark}[Multiplicity and the number of top strata]\label{rem:components}
For $h\in\Zb(n)$, the orbits $O_k\subseteq Z(n)$ whose closures contain $\I(h)$ and whose dimension is largest are precisely those with $h\sqsubseteq k$ and $\#\Lk(k)=\codim_{\I(h)}Z(n)$, by Proposition~\ref{prop:codim}. The top-dimensional components of $Z(n)$ itself are described in~\cite{lehalleur-rimanyi,koncki-rimanyi}. In Example~\ref{ex:depth4} there are four such orbits, whereas $m_h=3$. Thus the multiplicity is not, in general, the number of these orbits. A general relation between the two numbers remains open.
\end{remark}
\section{Proofs of the main theorems}\label{sec:proofs}

\begin{proof}[Proof of Theorem~\ref{thm:A}]
The case $D=0$ was treated after the statement. Let $D\ge1$, $n\in\N^{D+1}$, $\Phi\in\R^{n_D\times n_0}$, $A\in F_\Phi$, $r=\rk\Phi$, and let $h_A\in\Zb(n-r)$ be the residual barcode of $A$. By Corollary~\ref{cor:zero-A} applied to the width vector $n-r$ and the barcode $h_A$, the function $K_{n-r}$ has a local volume order $(\lambda_{h_A},m_{h_A})$ at $\I(h_A)$, with $2\lambda_{h_A}=\codim_{\I(h_A)}Z(n-r)$. By Lemma~\ref{lem:compare}, $L_{n-r,0}=\frac{1}{2}K_{n-r}$ has the same local volume order. By Proposition~\ref{prop:reduction}, $L_\Phi$ has local volume order $(\frac{1}{2}\creg(n,r)+\lambda_{h_A},\,m_{h_A})$ at $A$, and
\[
2\lambda_A(L_\Phi)=\creg(n,r)+2\lambda_{h_A}=\creg(n,r)+\codim_{\I(h_A)}Z(n-r)=\codim_AF_\Phi .
\]
By Lemma~\ref{lem:order-basic}(b) the same pair is the local volume order for every Haar measure, every prior and every norm, and by Lemma~\ref{lem:order-basic}(a) it is unique. Finally $\codim_AF_\Phi=N(n)-d_A(F_\Phi)$ is an integer, since $d_A(F_\Phi)$ is the dimension of a stratum, by~\eqref{eq:localdim-strata} and Lemma~\ref{lem:fibre-strata}.
\end{proof}

\begin{proof}[Proof of Theorem~\ref{thm:B}]
With the notation of the previous proof, Corollary~\ref{cor:zero-A}, Proposition~\ref{prop:codim} and Theorem~\ref{thm:mult}, applied to $h_A\in\Zb(n-r)$, give
\begin{gather*}
2\lambda_{h_A}=\codim_{\I(h_A)}Z(n-r)=C^\sqsubseteq(h_A)=\min\{\#\Lk(k):\ k\in\Zb(n-r),\ h_A\sqsubseteq k\},\\
m_{h_A}=1+F^\sqsubseteq(h_A),
\end{gather*}
and the set defining $C^\sqsubseteq(h_A)$ is nonempty (Theorem~\ref{thm:coefficient}). By Proposition~\ref{prop:closure}, $\sqsubseteq$ coincides with $\preceq$ on $\Zb(n-r)$. Hence $C^\sqsubseteq(h_A)=C(h_A)$, ``$\sqsubseteq$-critical'' means ``critical'', and $F^\sqsubseteq(h_A)=F(h_A)$. The theorem follows from Proposition~\ref{prop:reduction}, as in the proof of Theorem~\ref{thm:A}.
\end{proof}

\begin{remark}[What depends on what]\label{rem:dependencies}
Theorem~\ref{thm:A} uses Sections~\ref{sec:volume}--\ref{sec:exponent}; of Proposition~\ref{prop:closure} it uses only part (a), the easy direction. Theorem~\ref{thm:B} additionally uses Proposition~\ref{prop:closure}(b), proved in Appendix~\ref{app:generation}, and Section~\ref{sec:multiplicity}. Besides standard results of real analysis (the inverse function theorem, the change of variables formula, the Fubini--Tonelli theorem), the external inputs are: Fact~\ref{fact:gabriel}(a) (Gabriel's theorem for $\mathsf A_{D+1}$); the facts~\ref{H:def}--\ref{H:prod} on Hausdorff dimension; the homeomorphism theorem for transitive actions~\cite{helgason} and the fact that an injective immersion which is a homeomorphism onto its image is an embedding~\cite{lee}, both used in Lemma~\ref{lem:orbit}; Sard's theorem, used in Lemma~\ref{lem:fibre-strata}; and, in Appendix~\ref{app:generation}, the rank-matrix criterion for the Bruhat order~\cite{bjorner-brenti}.
\end{remark}

\begin{remark}[Dependence on the data]\label{rem:invariance}
That $(\lambda,m)$ depends only on the barcode of $A$ also follows directly: if $\beta(A)=\beta(A')$ and $A,A'$ are exact fits of $\Phi,\Phi'$, then $A'=g\cdot A$ for some $g\in G(n)$, and Step~1 of the proof of Proposition~\ref{prop:reduction} transports one problem to the other. The content of the theorems is the value of $(\lambda,m)$.
\end{remark}

\begin{remark}[Data on a subspace]\label{rem:covariance}
Assume $D\ge1$. Let $\Sigma\in\R^{n_0\times n_0}$ be symmetric positive semidefinite of rank $n'_0$. The population loss of linear regression with input covariance $\Sigma$ is $B\mapsto\frac{1}{2}\Fr{(P(B)-\Phi)\Sigma^{1/2}}^2$. Choose an orthonormal basis of $\R^{n_0}$ whose first $n'_0$ vectors span the image of $\Sigma$, let $\Sigma'$ be the (invertible) restriction of $\Sigma$ to its image, and for a chain $B=(B_0,\dots,B_{D-1})\in\Rep(n)$ split $B_0=(B'_0\mid B''_0)$ accordingly, with $B'_0\in\R^{n_1\times n'_0}$. Then the loss of $B$ equals $\frac{1}{2}\Fr{(P(B')-\Phi')\Sigma'^{1/2}}^2$, where $B'=(B'_0,B_1,\dots,B_{D-1})\in\Rep(n')$ with $n'=(n'_0,n_1,\dots,n_D)$ and $\Phi'$ is the restriction of $\Phi$ to the image of $\Sigma$; it lies between two positive multiples of $L_{n',\Phi'}(B')$, and the $(n_0-n'_0)n_1$ entries of $B''_0$ are spectator variables. By Lemmas~\ref{lem:compare} and~\ref{lem:spectator}, at a zero $A$ of this loss the local volume order is that of $L_{n',\Phi'}$ at $A'$, given by Theorems~\ref{thm:A} and~\ref{thm:B} for the width vector $n'$. A positive semidefinite weight on the output side is handled in the same way, with the rows of $B_{D-1}$ in place of the columns of $B_0$. A degenerate quadratic form on $\R^{n_D\times n_0}$ that is not of this product form is not covered; for $D=2$ it is studied in~\cite{ferry-monod-yano}.
\end{remark}
\section{The formalization}\label{sec:formal}

\subsection{What is formalized}
The Lean theorem \lean{MAIS.ExactFitStatement.exact_fit_classification}, stated and translated in Appendix~\ref{app:lean}, asserts the following for all $D$, $n$, every Haar measure $\mu$, every exact fit $A$ of $\Phi$, and every function $\varphi$ continuous at $A$ with $\varphi(A)>0$. There are $c,m\in\N$ such that:
\begin{enumerate}[label=(F\arabic*)]
\item the formal version of the program $\Prog$ returns exactly $\{(c,m)\}$ on $(n,R(A))$;
\item $c$ is the local Hausdorff codimension of $F_\Phi$ at $A$;
\item for every $(\lambda,m')$, $L_\Phi$ has local volume order $(\lambda,m')$ at $A$ if and only if $(\lambda,m')=(c/2,m)$.
\end{enumerate}
By Proposition~\ref{prop:formal-implies}, this implies Theorem~\ref{thm:A} (with Lemma~\ref{lem:order-basic}(b) for the independence of the norm) and, with Proposition~\ref{prop:prog-translation} and Lemma~\ref{lem:nonempty}, Theorem~\ref{thm:B} except its last sentence, which is proved only in the text. The proof of Proposition~\ref{prop:prog-translation} uses three results of the text (Proposition~\ref{prop:closure}(a), Theorem~\ref{thm:coefficient} and Lemma~\ref{lem:stationarity}) at one point, as explained after it. The formal proof of this statement only composes three theorems of the formal development:
\begin{itemize}
\item \lean{MAIS.Run2ArithmeticPositive.exactFitPairs_nonempty_singleton} (the program returns one pair);
\item \lean{MAIS.Run2ArithmeticPositive.entry_half_prior_pair_iff} (local volume orders are characterized by the program);
\item \lean{MAIS.Run2ArithmeticPositiveGeometry.exactFitPair_literal_hausdorff_codimension} (the first coordinate is the local codimension);
\end{itemize}
together with the definitional equivalence \lean{MAIS.DirectBandOrder.measuredBandPair_iff}, whose proof is \lean{rfl}. The public file \lean{ExactFitFull_2026_10_28_UTC.lean} in~\cite{feld-deep} concatenates $283$ source modules, including this statement. Its only external import is Mathlib at commit \lean{664cbbe52927ce6acaa39ca35498336116e123b8}. The Lean version is \lean{4.34.0-rc2}. Module names in this section identify declarations within that file; separate module files are not required to check it.

The formal development follows the same strategy as Sections~\ref{sec:quiver}--\ref{sec:multiplicity}, with some technical differences. For instance, the radial integration of Proposition~\ref{prop:recursion} is done with dyadic annuli instead of weighted polar coordinates, and Proposition~\ref{prop:generation} is proved with a different reduction to permutations.

\subsection{What a reader has to trust}
A Lean proof is accepted by Lean's kernel, a small type checker. Besides the kernel, the trusted base consists of the axioms used, the meaning of the Mathlib definitions that occur in the statement, and the translation of the statement into mathematics.
\begin{itemize}
\item \emph{Axioms.} The command \leanhash\lean{print axioms} reports that the theorem depends only on \lean{propext}, \lean{Classical.choice} and \lean{Quot.sound}, the standard axioms of classical mathematics in Lean. A \lean{sorry} or a newly declared axiom would appear in this list. The file checks this list at compile time.
\item \emph{Definitions.} Appendix~\ref{app:lean} quotes the controlling project definitions and the Mathlib definitions that fix the norms, measure and Hausdorff dimension, with a translation; the path product is characterized by the three displayed proved identities. Other standard Mathlib notions, such as rank and finite dimension, retain their usual meanings. The audit in \lean{exact-fit-local/} computes the \emph{definitional closure} of the statement: the constants reachable from its type through types and bodies of definitions, but not through proofs. It prints every repository constant in this closure. This list is what a reader has to check against Appendix~\ref{app:lean} and the public file.
\item \emph{Compilation.} See Sections~\ref{sec:formal-records} and~\ref{sec:formal-local}.
\end{itemize}

\subsection{Deceptive formal statements}\label{sec:formal-security}
A kernel-checked theorem can still fail to mean what its reader thinks. Lean allows user-defined notation, local instances (which can change, for example, which norm $\norm\cdot$ denotes), shadowing of names, metaprograms run at compile time, and identifiers that look alike, such as \lean{rn} and \lean{m}, or Latin and Greek letters. The kernel checks the elaborated proposition, not the reader's interpretation of the source text. The constituent sources and their concatenation were inspected for these mechanisms. The source checks found:
\begin{itemize}
\item no \lean{axiom}, \lean{sorry}/\lean{admit}, \lean{native_decide}, \lean{implemented_by}, \lean{extern}, \lean{unsafe}, or \lean{debug.skipKernelTC};
\item no metaprogramming commands (\lean{macro}, \lean{syntax}, \lean{elab}, \lean{run_cmd}, \leanhash\lean{eval}, \lean{initialize});
\item no global notation (only \lean{local notation} abbreviations inside proofs of the root depth-two file), and no reopening of Mathlib namespaces (in particular none of \lean{Matrix}, \lean{Real}, \lean{MeasureTheory}, \lean{Metric});
\item no instance of a norm, metric, measure or measurable-space class (the instances declared are decidability and finiteness instances, Prop-valued instances such as \lean{BorelSpace}, and the group action on chains);
\item no project declaration whose name could shadow a Mathlib name occurring in the statement;
\item no non-ASCII characters other than standard mathematical symbols and Greek letters, and no invisible or bidirectional characters.
\end{itemize}
Lean does not accept unquoted identifiers in the Cyrillic alphabet, so the most common look-alike letters are excluded automatically. In the final statement, every project name is either fully qualified or defined immediately before the statement. These source checks apply to the released bytes, not to arbitrary later Lean files.

Many files raise the option \lean{maxHeartbeats} (up to $3.2\times10^6$ instead of the default $2\times10^5$). This option bounds the time spent by the \emph{elaborator}, which searches for proofs and fills in implicit arguments. It plays no role in the kernel check, so it cannot make a false statement provable; it indicates slow proofs, which may need maintenance when the dependencies change. The option \lean{backward.isDefEq.respectTransparency false}, also used in many files, likewise only affects elaboration.


\subsection{The public proof file and its checks}\label{sec:formal-records}

The public repository~\cite{feld-deep} contains \lean{ExactFitFull_2026_10_28_UTC.lean} ($63\,475$ lines). Its SHA-256 is
\begin{center}\small
\lean{b362a3ee35df68adad7aa24a1ddd288e63e4e59d074445aaafef4093e8462a2c}.
\end{center}
Its header lists the constituent modules and their source hashes. Their bodies are placed inside sections, with internal imports removed. Concatenation changes which declarations are visible during elaboration; equality of source fragments alone is therefore not a proof that the elaborated statement is unchanged. The fingerprint comparison described below checks the statement and its controlling definitions in addition to compiling the proof.

The author ran the released file through the local harness on 28 October 2026. The published receipt in \lean{exact-fit-local/logs/} records a successful target build, ordinary kernel replay, the axiom list \lean{[propext, Classical.choice, Quot.sound]}, and a fingerprint match. The target build reused cached compiled output in that run; the replay was executed anew. The published replay, audit and fingerprint logs have the hashes listed in the receipt. The scan log was redacted after the run and consequently no longer has the receipt's scan-log hash; the accompanying notice records this limitation. The receipt is unsigned. It is execution evidence, not a certification of the translation into mathematics.

The logs in \lean{verification/2026-10-28-agent-check/} record a separate check on 28 October that compiled the exact same file from source with the pinned toolchain and cached Mathlib, then replayed its compiled module and recomputed the fingerprint. Compilation exited with status $0$ after about $14$ minutes, with one tactic-style warning; replay exited with status $0$, and the fingerprint matched the reference byte for byte. This was an agent-operated execution of Lean, not an independent implementation of its kernel. No fresh replay of all imported Mathlib and Lean-core declarations is claimed for either check.

\subsection{Reproducing the check}\label{sec:formal-local}

The folder \lean{exact-fit-local/} provides a Podman/Docker harness. Its image contains Ubuntu and the Lean version manager. Setup obtains the pinned Lean and Mathlib; subsequent checks run without network access as an unprivileged container user, with the input source mounted read-only. The command \lean{all} performs the following operations.
\begin{enumerate}[label=(\roman*)]
\item A source scan reports proof placeholders, added axioms, native-evaluation escapes, metaprogramming commands, notation declarations and invisible control characters. A lexical scan is a reading aid, not a proof that an arbitrary file cannot execute code.
\item A target build compiles the file, or reuses an unchanged compiled target. Source compilation checks the final \leanhash\lean{guard_msgs} assertion of the theorem's axiom list.
\item \lean{leanchecker} replays the declarations of the compiled module with Lean's kernel. The entire project proof is in that module. Ordinary replay takes Mathlib and Lean core as given; the option \lean{FRESH=1} also replays those imports with the same kernel.
\item The audit prints the statement, once with notation and once with all implicit arguments and instances displayed; it prints the axiom list and a fingerprint of the definitions that determine the statement's meaning.
\item The fingerprint is compared with \lean{exact-fit-local/reference/fingerprint-multi.txt}, and the harness writes a receipt and full logs.
\end{enumerate}
The reference fingerprint has SHA-256
\begin{center}\small
\lean{34e2fa7f3635b4f82eb985509eb7abad5e820413bc0eb34cd06241be172169d6}.
\end{center}
It contains the printed statement, $79$ project definitions, inductive types and constructors, $17$ distinct theorem statements used inside them, and the names of the external constants. For project definitions it records structural hashes of types and bodies. Bound-variable names and metadata are ignored, private declaration names lose their module prefix, and proofs are identified by their propositions. Proof irrelevance justifies the last normalization. The structural hash is not cryptographic and the comparison is not a formally proved equivalence or an adversarial security mechanism. It supplies evidence against accidental changes during extraction; the mathematical interpretation still requires reading the definitions.

The harness README gives the commands and dependency-reuse options. A shared O77 dependency volume remains writable and must not be deleted by the harness's \lean{clean} command. The source file should remain unchanged throughout a run. Preserve the logs together with the receipt, and require an explicit fingerprint match rather than the mere presence of a receipt file.
\section{Disclosure of AI use}\label{sec:ai}

This work was directed by Niels Feld and produced with substantial assistance from AI systems. The following distinctions concern the mathematical text, the formal statements, the formal proofs and their verification.

\subsection*{Origin and contributions}
\begin{itemize}
\item \textbf{Research direction.} Following Lionel Levine's invitation and the Math for AI Safety problems~\cite{levine-mais,mais-repo}, the author directed the Lean formalization of the depth-two problem MAIS-O77~\cite{mais-o77,feld-o77}, then requested the arbitrary-depth exact-fit generalization studied here.
\item \textbf{Formal proofs.} AI agents using OpenAI GPT models through Codex wrote the Lean development under the author's direction. The author did not write these proofs. The checked formal endpoint is (F1)--(F3) of Appendix~\ref{app:lean}. Proposition~\ref{prop:formal-implies} and the discussion after Proposition~\ref{prop:prog-translation} specify the additional mathematical arguments connecting it to Theorems~\ref{thm:A} and~\ref{thm:B}; in particular, the selected endpoint does not itself assert the last sentence of Theorem~\ref{thm:B}.
\item \textbf{Mathematical exposition.} Claude (Anthropic) wrote the natural-language arguments from the formal development and its documentation, then revised them after the author's reading. OpenAI GPT/Codex and Claude agents reviewed successive versions and proposed corrections. GPT/Codex first wrote the depth-four example and its explanation; Claude wrote the proof in Appendix~\ref{app:generation}, which differs from the formal proof. GPT/Codex prepared the revision of 29 October 2026, 00:19:56, including explicit estimates and the public source references; Claude then made the corrections listed in the revision notes accompanying this version. \DIFaddbegin \DIFadd{The present bibliographic correction was prepared with GPT/Codex after the author checked the arXiv record.
}\DIFaddend \item \textbf{Statement and checking tools.} Claude wrote the single-statement wrapper and its translation, the single-file extractor and the checking harness, the latter adapted from~\cite{feld-o77}. The formal wrapper and its proof have been accepted by Lean; its translation into mathematics remains a reading task. Source quotations are copied mechanically and checked against the released Lean file or pinned Mathlib sources.
\item \textbf{Literature.} AI agents assisted with the search and attribution of prior work. References distinguish results used in proofs from related precedents. No claim that the literature search is exhaustive is made.
\end{itemize}

\subsection*{Human review and execution}
The author has read the mathematical exposition and the formal statement with its controlling definitions. He found no error in those readings, but has not checked every algebraic calculation by hand and does not claim a complete certification of the translation. He has not read all the formal proofs. His expertise is in algebraic geometry, motivic homotopy theory and K-theory, rather than singular learning theory. He ran the local Lean check described in Section~\ref{sec:formal-records}; this is distinct from the source compilations and checks performed by AI agents.

The present version incorporates further AI-proposed revisions. The author's readings of earlier versions should not be interpreted as approval of every later edit. The red-and-blue comparison files accompanying this version identify the mathematical and editorial changes for review.

\subsection*{Scope of machine verification and responsibility}
Lean's kernel checks proof terms for the elaborated formal statement, subject to the axioms and imported declarations described in Section~\ref{sec:formal}. It does not check whether the statement expresses the intended mathematics, whether the prose proof is correct, or whether a bibliographic attribution is complete. Kernel replay uses the same kernel; agreement among AI reviewers is not independent human refereeing. No referee report is claimed. The author is responsible for the version he chooses to circulate and for its mathematical claims.

\section*{Acknowledgements}
The author acknowledges Lionel Levine and the Math for AI Safety contributors for the questions that initiated this work; Rob Sneiderman and Svyatoslav Novikov, with OpenAI Codex, for the posted depth-two candidate solutions~\cite{mais-issue-3,mais-issue-12}; Simon Pepin Lehalleur and Richárd Rimányi for the geometric study of the fibres, the global half-codimension theorem and the pointwise conjecture~\cite{lehalleur-rimanyi}; and Miki Aoyagi and Sumio Watanabe for the global learning-coefficient calculations~\cite{aoyagi-watanabe,aoyagi-2024}. Related contributions, including precedents not used in the proof, are identified in Section~\ref{sec:intro} and where they enter the argument. The formal development depends on Lean and the work of the Mathlib contributors. The substantial contributions of the AI systems are described above rather than presented as personal communications or human refereeing.
\appendix
\section{The formal statement and its translation}\label{app:lean}

This appendix is for a reader who wants to check that the Lean theorem says what Theorems~\ref{thm:A} and~\ref{thm:B} say. Each project quotation below is copied mechanically from the public file \lean{ExactFitFull_2026_10_28_UTC.lean} identified in Section~\ref{sec:formal-records}; its displayed line numbers refer to that file. Mathlib quotations are from the pinned public Mathlib source. Each block gives its source and the beginning of its SHA-256 hash.

\subsection{Reading conventions}\label{app:conventions}
\begin{itemize}
\item \lean{ℕ}, \lean{ℝ}, \lean{ℚ} are $\N$, $\R$, $\Q$. \lean{ℝ≥0∞} (also \lean{ENNReal}) is $[0,+\infty]$. \lean{Fin k} is $\{0,\dots,k-1\}$. For $e$ in \lean{Fin D}, \lean{e.castSucc} and \lean{e.succ} are the vertices $e$ and $e+1$, and \lean{Fin.last D} is $D$.
\item Subtraction on $\N$ is truncated at $0$; so is subtraction on $[0,+\infty]$, with $(+\infty)-x=+\infty$ for finite $x$. \lean{ENNReal.ofReal x} is $\max(x,0)\in[0,+\infty[$.
\item \lean{Matrix (Fin p) (Fin q) ℝ} is $\R^{p\times q}$, as functions of (row, column). \lean{M.rank} is the rank.
\item \lean{Real.rpow ε λ} is $\eps^\lambda$ (used only for $\eps>0$), \lean{Real.log} is the natural logarithm, \lean{Real.exp} the exponential, \lean{Real.sqrt} the nonnegative square root.
\item \lean{∃ c m : ℕ, P} means ``there exist $c,m\in\N$ with $P$''. \lean{{(c, m)}} is a one-element finite set, and \lean{s = {(c, m)}} says that the finite set $s$ is exactly $\{(c,m)\}$.
\end{itemize}
The following Mathlib definitions control the meaning of the statement.

\emph{Norms.} After \lean{open scoped Matrix.Norms.Frobenius}, the norm of a matrix is the Frobenius norm:
\begin{LeanCode}
theorem frobenius_norm_def (A : Matrix m n α) :
    ‖A‖ = (∑ i, ∑ j, ‖A i j‖ ^ (2 : ℝ)) ^ (1 / 2 : ℝ) :=
  (congr_arg ((↑) : ℝ≥0 → ℝ) (frobenius_nnnorm_def A)).trans <| by simp [NNReal.coe_sum]
\end{LeanCode}
\leansource{1}{Mathlib 664cbbe5, Mathlib/Analysis/Matrix/Normed.lean}{561}{563}{e430f9176d36}
\begin{LeanCode}
namespace Norms.Frobenius
attribute [scoped instance]
  Matrix.frobeniusSeminormedAddCommGroup
  Matrix.frobeniusNormedAddCommGroup
  Matrix.frobeniusNormedSpace
  Matrix.frobeniusNormedRing
  Matrix.frobeniusNormedAlgebra
  Matrix.frobeniusIsBoundedSMul
  Matrix.frobeniusNormSMulClass
end Norms.Frobenius
\end{LeanCode}
\leansource{2}{Mathlib 664cbbe5, Mathlib/Analysis/Matrix/Normed.lean}{681}{690}{e430f9176d36}
\emph{Translation.} $\norm A=\Fr A$ for every real matrix $A$. A chain is an element of a finite product, which Mathlib normalizes with the sup norm:
\begin{LeanCode}
/-- Finite product of seminormed groups, using the sup norm. -/
@[to_additive /-- Finite product of seminormed groups, using the sup norm. -/]
instance Pi.seminormedGroup : SeminormedGroup (∀ i, G i) where
  norm f := ↑(Finset.univ.sup fun b => ‖f b‖₊)
\end{LeanCode}
\leansource{3}{Mathlib 664cbbe5, Mathlib/Analysis/Normed/Group/Constructions.lean}{301}{304}{ab6a98350718}
\emph{Translation.} For a chain $B$, $\norm B=\max_e\Fr{B_e}$, so \lean{Metric.ball A r}$=B^\infty_r(A):=\{B:\max_e\Fr{B_e-A_e}<r\}$. This is \emph{not} the Euclidean ball; it is used only in the codimension conjunct, where the choice of norm does not matter (Definition~\ref{def:codim}).

\emph{Measures.} A measure on the chain space is defined on its Borel $\sigma$-algebra. An additive Haar measure is:
\begin{LeanCode}
/-- A measure on an additive group is an additive Haar measure if it is left-invariant, and
gives finite mass to compact sets and positive mass to open sets.

Textbooks generally require an additional regularity assumption to ensure nice behavior on
arbitrary locally compact groups. Use `[IsAddHaarMeasure μ] [Regular μ]` or
`[IsAddHaarMeasure μ] [InnerRegular μ]` in these situations. Note that a Haar measure in our
sense is automatically regular and inner regular on second countable locally compact groups, as
checked just below this definition. -/
class IsAddHaarMeasure {G : Type*} [AddGroup G] [TopologicalSpace G] [MeasurableSpace G]
    (μ : Measure G) : Prop
    extends IsFiniteMeasureOnCompacts μ, IsAddLeftInvariant μ, IsOpenPosMeasure μ
\end{LeanCode}
\leansource{4}{Mathlib 664cbbe5, Mathlib/MeasureTheory/Group/Measure.lean}{764}{774}{cba62e69e14a}
\emph{Translation.} A translation-invariant measure, finite on compact sets and positive on nonempty open sets. On a finite-dimensional real vector space these are exactly the positive multiples of Lebesgue measure.
\begin{LeanCode}
/-- Given a measure `μ : Measure α` and a function `f : α → ℝ≥0∞`, `μ.withDensity f` is the
measure such that for a measurable set `s` we have `μ.withDensity f s = ∫⁻ a in s, f a ∂μ`. -/
noncomputable
def Measure.withDensity {m : MeasurableSpace α} (μ : Measure α) (f : α → ℝ≥0∞) : Measure α :=
  Measure.ofMeasurable (fun s _ => ∫⁻ a in s, f a ∂μ) (by simp) fun _ hs hd =>
    lintegral_iUnion hs hd _
\end{LeanCode}
\leansource{5}{Mathlib 664cbbe5, Mathlib/MeasureTheory/Measure/WithDensity.lean}{36}{41}{975237c8f778}
\emph{Translation.} $(\mu.\mathrm{withDensity}\,f)(S)=\int^-_Sf\,d\mu$ for measurable $S$, where $\int^-$ (Mathlib's \lean{lintegral}) is the lower integral of Section~\ref{sec:measures}. The sets $S$ to which it is applied below are open, hence measurable.

\emph{Hausdorff dimension.}
\begin{LeanCode}
/-- Hausdorff dimension of a set in an (e)metric space. -/
@[irreducible] noncomputable def dimH (s : Set X) : ℝ≥0∞ := by
  borelize X; exact ⨆ (d : ℝ≥0) (_ : @hausdorffMeasure X _ _ ⟨rfl⟩ d s = ∞), d
\end{LeanCode}
\leansource{6}{Mathlib 664cbbe5, Mathlib/Topology/MetricSpace/HausdorffDimension.lean}{96}{98}{d525b021e2f4}
\emph{Translation.} $\dimH S=\sup\{d\ge0:\ \mathcal H^d(S)=\infty\}$, computed with the metric of the ambient normed space. This agrees with Section~\ref{sec:hausdorff}.

\subsection{The statement}
The theorem and the six definitions it introduces occur near the end of the public file:
\begin{LeanCode}
open MeasureTheory
open scoped ENNReal Matrix.Norms.Frobenius

namespace MAIS.ExactFitStatement

/-- Coordinate index family: the vertex `v` carries the coordinate set `Fin (n v)`. -/
abbrev idx {D : ℕ} (n : Fin (D + 1) → ℕ) : Fin (D + 1) → Type :=
  fun v => Fin (n v)

/-- A chain `A = (A_0, …, A_{D-1})`, where `A e` is a real `n (e+1) × n e` matrix.
It unfolds to `(e : Fin D) → Matrix (Fin (n e.succ)) (Fin (n e.castSucc)) ℝ`. -/
abbrev Chain {D : ℕ} (n : Fin (D + 1) → ℕ) : Type :=
  MAIS.QuiverRealOrbit.Rep (idx n)

/-- The product `A_{D-1} ⋯ A_0` (the identity matrix when `D = 0`). -/
def product {D : ℕ} (n : Fin (D + 1) → ℕ) (A : Chain n) :
    Matrix (Fin (n (Fin.last D))) (Fin (n 0)) ℝ :=
  MAIS.MatrixChainIsomorphism.pathMatrix (ι := idx n) A 0 (Fin.last D) (Nat.zero_le D)

/-- `rankTable n A i j` is the rank of `A_{j-1} ⋯ A_i` when `i ≤ j`, and `0` otherwise. -/
def rankTable {D : ℕ} (n : Fin (D + 1) → ℕ) (A : Chain n) (i j : Fin (D + 1)) : ℕ :=
  MAIS.Run2BarcodeEnumeration.pathRankInput (idx n) A i j

/-- The loss `B ↦ ½ ‖P(B) - Φ‖²`, with the Frobenius norm. -/
def halfLoss {D : ℕ} (n : Fin (D + 1) → ℕ)
    (Φ : Matrix (Fin (n (Fin.last D))) (Fin (n 0)) ℝ) (B : Chain n) : ℝ :=
  (1 / 2 : ℝ) * ‖product n B - Φ‖ ^ 2

/-- The Euclidean ball of radius `δ` around `A`, in all matrix entries. -/
def entryBall {D : ℕ} (n : Fin (D + 1) → ℕ) (A : Chain n) (δ : ℝ) : Set (Chain n) :=
  {B | Real.sqrt (∑ e : Fin D, ∑ r : Fin (n e.succ), ∑ s : Fin (n e.castSucc),
      (B e r s - A e r s) ^ 2) < δ}
\end{LeanCode}
\leansource{7}{ExactFitFull_2026_10_28_UTC.lean}{63365}{63396}{b362a3ee35df}
\emph{Translations.}
\begin{itemize}
\item \lean{Chain n} is $\Rep(n)$. It unfolds to the type of tuples whose component $e$ is a real $n_{e+1}\times n_e$ matrix, by the definition
\begin{LeanCode}
abbrev Rep (ι : Fin (D + 1) → Type*) :=
  (e : Fin D) → Matrix (ι e.succ) (ι e.castSucc) ℝ
\end{LeanCode}
\leansource{8}{ExactFitFull_2026_10_28_UTC.lean}{788}{789}{b362a3ee35df}
\item \lean{product n A} is $P(A)=A_{D-1}\cdots A_0$. The function \lean{pathMatrix A i j} is defined through Mathlib's category of composable arrows, but it is pinned down by the three proved lemmas
\begin{LeanCode}
noncomputable def pathMatrix (A : Rep ι) (i j : Fin (d + 1)) (hij : i ≤ j) :
    Matrix (ι j) (ι i) ℝ := LinearMap.toMatrix' ((matrixChain A).map i j hij)

lemma pathMatrix_self (A : Rep ι) (i : Fin (d + 1)) : pathMatrix A i i le_rfl = 1 := by
  rw [pathMatrix, (matrixChain A).map_id]
  exact LinearMap.toMatrix'_id

lemma pathMatrix_comp (A : Rep ι) (i j k : Fin (d + 1)) (hij : i ≤ j) (hjk : j ≤ k) :
    pathMatrix A j k hjk * pathMatrix A i j hij = pathMatrix A i k (hij.trans hjk) := by
  rw [pathMatrix, pathMatrix, ← LinearMap.toMatrix'_comp, (matrixChain A).map_comp]
  rfl

lemma pathMatrix_arrow (A : Rep ι) (e : Fin d) :
    pathMatrix A e.castSucc e.succ (by exact Nat.le_succ e.val) = A e := by
  rw [pathMatrix, matrixChain, ofArrows_map_succ]
  exact LinearMap.toMatrix'_toLin' _
\end{LeanCode}
\leansource{9}{ExactFitFull_2026_10_28_UTC.lean}{4551}{4566}{b362a3ee35df}
which say $M_{ii}=\id$, $M_{jk}M_{ij}=M_{ik}$ and $M_{e,e+1}=A_e$ for $M_{ij}=$ \lean{pathMatrix A i j}. By induction on $j-i$, every family with these three properties equals $(A_{j-1}\cdots A_i)_{i\le j}$. Hence \lean{pathMatrix A i j}$=P_{ij}(A)$ and \lean{product n A}$=P(A)$; for $D=0$ this is $\id_{n_0}$.
\item \lean{rankTable n A i j} is $R_{ij}(A)$ for $i\le j$ and $0$ for $i>j$, by
\begin{LeanCode}
def pathRankInput (A : Rep ι) (i j : Fin (D + 1)) : ℕ :=
  if hij : i ≤ j then (pathMatrix A i j hij).rank else 0
\end{LeanCode}
\leansource{10}{ExactFitFull_2026_10_28_UTC.lean}{56724}{56725}{b362a3ee35df}
\item \lean{halfLoss n Φ B} is $L_\Phi(B)=\frac{1}{2}\Fr{P(B)-\Phi}^2$, and \lean{entryBall n A δ} is the Euclidean ball $B_\delta(A)$ in all entries.
\end{itemize}
The theorem:
\begin{LeanCode}
/-- **Exact-fit classification.**  Let `A` be a chain with product `Φ`, let `μ` be an
additive Haar measure and `φ` a real function, continuous at `A`, with `φ A > 0`.
Then there are natural numbers `c` and `m` such that

1. the finite program returns exactly `{(c, m)}`;
2. `c` is the local Hausdorff codimension of the fibre `{B | P(B) = Φ}` at `A`
   (ambient dimension minus the infimum over all radii `r > 0` of the Hausdorff
   dimension of the fibre inside the open ball of radius `r`; the balls and the
   Hausdorff dimension are those of the norm `max_e ‖B e‖_F`);
3. for every real `lam` and natural `m'`, the band volume
   `(δ, ε) ↦ ∫_{B ∈ entryBall A δ, |L(B) - L(A)| < ε} φ⁺ dμ` has the two-sided
   power-log order `(lam, m')` (`HasOrder`) if and only if `lam = c / 2` and `m' = m`. -/
theorem exact_fit_classification
    {D : ℕ} (n : Fin (D + 1) → ℕ)
    (μ : Measure (Chain n)) [μ.IsAddHaarMeasure]
    (A : Chain n) (Φ : Matrix (Fin (n (Fin.last D))) (Fin (n 0)) ℝ)
    (hfit : product n A = Φ)
    (φ : Chain n → ℝ) (hφ : ContinuousAt φ A) (hφA : 0 < φ A) :
    ∃ c m : ℕ,
      MAIS.Run2ArithmeticPositiveData.exactFitPairs n (rankTable n A) = {(c, m)} ∧
      (Module.finrank ℝ (Chain n) : ℝ≥0∞) -
          (⨅ (r : ℝ) (_ : 0 < r),
            dimH ({B : Chain n | product n B = Φ} ∩ Metric.ball A r)) =
        (c : ℝ≥0∞) ∧
      ∀ (lam : ℝ) (m' : ℕ),
        MAIS.DirectBandOrder.HasOrder
            (fun δ ε => μ.withDensity (fun B => ENNReal.ofReal (φ B))
              {B | B ∈ entryBall n A δ ∧ |halfLoss n Φ B - halfLoss n Φ A| < ε})
            lam m' ↔
          (lam = (c : ℝ) / 2 ∧ m' = m) := by
\end{LeanCode}
\leansource{11}{ExactFitFull_2026_10_28_UTC.lean}{63398}{63427}{b362a3ee35df}
The predicate \lean{HasOrder} is defined in one block:
\begin{LeanCode}
/-- Two-sided small-loss order at every sufficiently small fixed radius. -/
def HasOrder (volume : _root_.Real → _root_.Real → ENNReal)
    (exponent : _root_.Real) (multiplicity : Nat) : Prop :=
  0 ≤ exponent ∧ 1 ≤ multiplicity ∧
  ∃ radius : _root_.Real, 0 < radius ∧
    ∀ delta : _root_.Real, 0 < delta → delta < radius →
      ∃ lower upper cutoff : _root_.Real,
        0 < lower ∧ lower ≤ upper ∧ 0 < cutoff ∧ cutoff < Real.exp (-1) ∧
        ∀ epsilon : _root_.Real, 0 < epsilon → epsilon < cutoff →
          ENNReal.ofReal lower *
            ENNReal.ofReal (Real.rpow epsilon exponent *
              (Real.log (1 / epsilon)) ^ (multiplicity - 1)) ≤ volume delta epsilon ∧
          volume delta epsilon ≤ ENNReal.ofReal upper *
            ENNReal.ofReal (Real.rpow epsilon exponent *
              (Real.log (1 / epsilon)) ^ (multiplicity - 1))
\end{LeanCode}
\leansource{12}{ExactFitFull_2026_10_28_UTC.lean}{63297}{63311}{b362a3ee35df}
\emph{Translation of \lean{HasOrder V λ m}.} For a function $V:\R\times\R\to[0,+\infty]$: $\lambda\ge0$, $m\ge1$, and there is $\delta_0>0$ such that for every $\delta\in\,]0,\delta_0[$ there are real numbers $c,C,\eps_0$ with $0<c\le C$ and $\eps_0\in\,]0,e^{-1}[$ such that
\[
c\,\eps^\lambda L(\eps)^{m-1}\le V(\delta,\eps)\le C\,\eps^\lambda L(\eps)^{m-1}\qquad\text{for all }\eps\in\,]0,\eps_0[ .
\]
Here $L(\eps)=\log(1/\eps)$ as in~\eqref{eq:L}. Indeed, for $\eps\in\,]0,e^{-1}[$ the real number $\eps^\lambda L(\eps)^{m-1}$ is positive, the natural subtraction $m-1$ is ordinary because $m\ge1$, and $\mathrm{ofReal}(c)\cdot\mathrm{ofReal}(t)=ct$ for $c,t\ge0$.

\emph{Translation of the theorem.} For all $D\in\N$ and $n\in\N^{D+1}$, every additive Haar measure $\mu$ on $\Rep(n)$, every $A\in\Rep(n)$ and $\Phi\in\R^{n_D\times n_0}$ with $P(A)=\Phi$, and every $\varphi:\Rep(n)\to\R$ continuous at $A$ with $\varphi(A)>0$, there are $c,m\in\N$ such that:
\begin{enumerate}[label=(F\arabic*)]
\item\label{F:prog} the formal program returns $\{(c,m)\}$ on the input $(n,R(A))$;
\item\label{F:codim} $N(n)\ \dot-\ \inf_{r>0}\dimH\bigl(F_\Phi\cap B^\infty_r(A)\bigr)=c$ in $[0,+\infty]$, where $\dot-$ is truncated subtraction and $\dimH$ uses the norm $\max_e\Fr{\cdot}$;
\item\label{F:order} for all $\lambda\in\R$ and $m'\in\N$: $L_\Phi$ has local volume order $(\lambda,m')$ at $A$, for the measure $\mu$, the prior $\varphi$ and the Euclidean balls, if and only if $\lambda=c/2$ and $m'=m$.
\end{enumerate}
Here~\ref{F:order} is Definition~\ref{def:order} verbatim, with $V^{\mu,\varphi}_{L_\Phi,A}(\delta,\eps)=(\mu.\mathrm{withDensity}\,\varphi^+)(S_{L_\Phi,A}(\delta,\eps))$ and $S_{L_\Phi,A}(\delta,\eps)=\{B\in B_\delta(A):\ \abs{L_\Phi(B)-L_\Phi(A)}<\eps\}$.

\begin{proposition}\label{prop:formal-implies}
The statement~\ref{F:prog}--\ref{F:order}, together with Lemma~\ref{lem:order-basic}(b), implies Theorem~\ref{thm:A}. Together moreover with Proposition~\ref{prop:prog-translation} below and Lemma~\ref{lem:nonempty}, it implies Theorem~\ref{thm:B} except its last sentence.
\end{proposition}
\begin{proof}
By~\ref{F:order} with $(\lambda,m')=(c/2,m)$, $L_\Phi$ has a local volume order at $A$ for $(\mu,\varphi)$, and by the converse direction it is unique. In~\ref{F:codim}, every subset of $\Rep(n)$ has Hausdorff dimension at most $N(n)$, so the truncated subtraction is ordinary. The infimum is $d_A(F_\Phi)$ (Definition~\ref{def:codim}, which allows any norm). Hence $c=\codim_AF_\Phi$ and $\lambda=c/2=\frac{1}{2}\codim_AF_\Phi$; since $c\in\N$, $2\lambda\in\N$. Independence of $(\mu,\varphi)$ is part of the quantification. Independence of the norm is not: \ref{F:order} uses Euclidean balls, and the other norms are covered by Lemma~\ref{lem:order-basic}(b). For Theorem~\ref{thm:B}, Proposition~\ref{prop:prog-translation} identifies the formal program with $\Prog$, and the nonemptiness asserted in Theorem~\ref{thm:B} is Lemma~\ref{lem:nonempty}. The last sentence of Theorem~\ref{thm:B}, $C(h_A)=\min\{\#\Lk(k):\ k\in\Zb(n-r),\ h_A\preceq k\}$, does not follow from~\ref{F:prog}--\ref{F:order}; it is proved in Section~\ref{sec:proofs} from Propositions~\ref{prop:codim} and~\ref{prop:closure}.
\end{proof}

The reader who only wants Theorem~\ref{thm:A} can stop here: \ref{F:codim} and~\ref{F:order} determine $c$ and $m$ without reference to the program.

\subsection{The program}
The formal program is \lean{exactFitPairs}. We quote its definitions in dependency order, then translate.
\begin{LeanCode}
def residualDimensions (dims : Fin (D + 1) → ℕ) (r : ℕ) : Fin (D + 1) → ℕ :=
  fun v => dims v - r

def residualRanks (ranks : Fin (D + 1) → Fin (D + 1) → ℕ) (r : ℕ) :
    Fin (D + 1) → Fin (D + 1) → ℕ := fun i j => ranks i j - r

abbrev ResidualCarrier (dims : Fin (D + 1) → ℕ) (r : ℕ) (v : Fin (D + 1)) :=
  Fin (residualDimensions dims r v)

/-- Cardinalities of the literal canonical finite residual carriers. These
are proved equal to the supplied residual dimensions below. -/
def residualArchitecture (dims : Fin (D + 1) → ℕ) (r : ℕ) : Fin (D + 1) → ℕ :=
  fun v => Fintype.card (ResidualCarrier dims r v)

def residualBudget (dims : Fin (D + 1) → ℕ) (r : ℕ) : ℕ :=
  ∑ v, residualArchitecture dims r v

def regularCoefficient (dims : Fin (D + 1) → ℕ) (r : ℕ) : ℕ :=
  r * (dims 0 + dims (Fin.last D) - r)

def residualNodes (dims : Fin (D + 1) → ℕ) (r : ℕ) :
    Finset (Multiplicities D (residualBudget dims r)) :=
  fastCandidates (residualArchitecture dims r) (fun _ _ => 0)

theorem residualArchitecture_eq (dims : Fin (D + 1) → ℕ) (r : ℕ) :
    residualArchitecture dims r = residualDimensions dims r := by
  funext v
  exact Fintype.card_fin _

theorem residualBudget_eq (dims : Fin (D + 1) → ℕ) (r : ℕ) :
    residualBudget dims r = ∑ v, (dims v - r) := by
  simp only [residualBudget, residualArchitecture_eq, residualDimensions]

def residualPairs (dims : Fin (D + 1) → ℕ)
    (ranks : Fin (D + 1) → Fin (D + 1) → ℕ) (r : ℕ) : Finset (ℕ × ℕ) :=
  computedPairs (residualNodes dims r) (residualRanks ranks r)

/-- The first returned integer is twice the actual exponent, and the second
is its logarithmic multiplicity. Actual correctness and singleton output
are proved separately; no desired measured value enters this program. -/
def exactFitPairs (dims : Fin (D + 1) → ℕ)
    (ranks : Fin (D + 1) → Fin (D + 1) → ℕ) : Finset (ℕ × ℕ) :=
  if D = 0 then {(0, 1)} else
    (residualPairs dims ranks (ranks 0 (Fin.last D))).image
      (fun p => (regularCoefficient dims (ranks 0 (Fin.last D)) + p.1, p.2))
\end{LeanCode}
\leansource{13}{ExactFitFull_2026_10_28_UTC.lean}{5560}{5604}{b362a3ee35df}
\emph{Translation.} For $D=0$ the output is $\{(0,1)\}$. Otherwise, with $r=R_{0D}$, the output is the image of the \emph{residual program} under $(c_0,m)\mapsto(\creg(n,r)+c_0,m)$, where $\creg$ is computed with truncated subtraction. On actual inputs $r\le n_0,n_D$, so it is ordinary. The residual program runs on the widths $n\dot-r=(n_v\dot-r)_v$ and the ranks $(R_{ij}\dot-r)_{i\le j}$ (exact on actual inputs, since $r\le n_v$ and $r\le R_{ij}$ for $i\le j$), with the slot bound $M_{\mathrm{slot}}=\sum_v(n_v\dot-r)$. Its candidate set is \lean{fastCandidates} for the widths $n\dot-r$ with zero lower bounds on the ranks. The theorems \lean{residualArchitecture_eq} and \lean{residualBudget_eq} only record that the cardinalities are the intended numbers.

\emph{Barcodes.} A barcode is encoded by multiplicities of pairs of endpoints:
\begin{LeanCode}
abbrev Endpoint (D : ℕ) := Fin (D + 2)
abbrev IntervalKey (D : ℕ) := Endpoint D × Endpoint D
abbrev Multiplicities (D M : ℕ) := IntervalKey D → Fin (M + 1)
abbrev Slot {D M : ℕ} (h : Multiplicities D M) := Σ p : IntervalKey D, Fin (h p).val

def birth {D M : ℕ} {h : Multiplicities D M} (s : Slot h) : ℕ := s.1.1.val
def death {D M : ℕ} {h : Multiplicities D M} (s : Slot h) : ℕ := s.1.2.val

def Valid {D M : ℕ} (h : Multiplicities D M) : Prop :=
  ∀ s : Slot h, birth s < death s

def vertexCount {D M : ℕ} (h : Multiplicities D M) (v : Fin (D + 1)) : ℕ :=
  Fintype.card {s : Slot h // birth s ≤ v.val ∧ v.val < death s}

def pathCount {D M : ℕ} (h : Multiplicities D M) (i j : Fin (D + 1)) : ℕ :=
  Fintype.card {s : Slot h // birth s ≤ i.val ∧ j.val < death s}

def Linked {D M : ℕ} {h : Multiplicities D M} (r s : Slot h) : Prop :=
  birth s < birth r ∧ birth r ≤ death s ∧ death s < death r
instance linkedDecidable {D M : ℕ} {h : Multiplicities D M} (r s : Slot h) :
    Decidable (Linked r s) := inferInstanceAs (Decidable (_ ∧ _ ∧ _))

def incoming {D M : ℕ} {h : Multiplicities D M} (r : Slot h) : ℕ :=
  Fintype.card {s : Slot h // Linked r s}
def outgoing {D M : ℕ} {h : Multiplicities D M} (s : Slot h) : ℕ :=
  Fintype.card {r : Slot h // Linked r s}
def Balanced {D M : ℕ} (h : Multiplicities D M) : Prop :=
  ∀ s : Slot h, birth s ≠ 0 → death s ≠ D + 1 → incoming s = outgoing s

def outputCount {D M : ℕ} (h : Multiplicities D M) : ℕ :=
  Fintype.card {p : Slot h × Slot h // Linked p.1 p.2 ∧ death p.1 = D + 1}

/-- Actual finite arithmetic feasibility for a prescribed architecture and
prescribed integer path-rank input. -/
def Feasible {D M : ℕ} (dims : Fin (D + 1) → ℕ)
    (ranks : Fin (D + 1) → Fin (D + 1) → ℕ) (h : Multiplicities D M) : Prop :=
  Valid h ∧ (∀ v, vertexCount h v = dims v) ∧
  pathCount h 0 (Fin.last D) = 0 ∧
  ∀ i j, i ≤ j → ranks i j ≤ pathCount h i j
\end{LeanCode}
\leansource{14}{ExactFitFull_2026_10_28_UTC.lean}{307}{345}{b362a3ee35df}
\emph{Dictionary.} A key \lean{(b, d)} with $b<d$ (\lean{Valid}) is the interval $\ii b{d-1}$, i.e.\ \lean{birth}$=b_s$ and \lean{death}$=e_s+1$. A \lean{Slot h} is a pair (key, copy index), i.e.\ an element of $\Sl(h)$. Then:
\begin{itemize}
\item \lean{vertexCount h v} is $\dim_v(h)$ and \lean{pathCount h i j} is $\rho_{ij}(h)$ (Definition~\ref{def:barcode-dims}), since $b\le i$ and $j<d$ means $\ii ij\subseteq\ii b{d-1}$.
\item \lean{Linked r s} is $s\lin r$, with the first argument the \emph{target}: $b_s<b_r\le e_s+1\le e_r$.
\item \lean{incoming r} is $\mathrm{in}(r)$ and \lean{outgoing s} is $\mathrm{out}(s)$.
\item \lean{Balanced} quantifies over interior slots (\lean{birth}$\ne0$ and \lean{death}$\ne D+1$, i.e.\ $b_s\ne0$ and $e_s\ne D$).
\item \lean{outputCount h} counts pairs $(t,s)$ with $s\lin t$ and $e_t=D$, i.e.\ $\cout(h)$.
\item \lean{Feasible dims ranks h} says: valid, dimension vector \lean{dims}, no full-span slot ($\rho_{0D}(h)=0$), and \lean{ranks}$_{ij}\le\rho_{ij}(h)$ for $i\le j$.
\end{itemize}
Multiplicities are bounded by $M_{\mathrm{slot}}$ (\lean{Fin (M + 1)}, where \lean{M} is $M_{\mathrm{slot}}$). This loses nothing: a barcode with dimension vector $n\dot-r$ has at most $\sum_v(n_v\dot-r)=M_{\mathrm{slot}}$ slots.

\emph{Candidates.}
\begin{LeanCode}
abbrev ValidKey (D : ℕ) := {p : IntervalKey D // p.1.val < p.2.val}
abbrev Assignment (D M : ℕ) := Fin M → Option (ValidKey D)

def histogram {D M : ℕ} (f : Assignment D M) : Multiplicities D M := fun p =>
  ⟨Fintype.card {i : Fin M // Option.map Subtype.val (f i) = some p}, by
    have h := Fintype.card_le_of_injective
      (Subtype.val : {i : Fin M // Option.map Subtype.val (f i) = some p} → Fin M)
      Subtype.val_injective
    simpa only [Fintype.card_fin] using Nat.lt_succ_of_le h⟩

def generated {D M : ℕ} : Finset (Multiplicities D M) :=
  Finset.univ.image histogram

def fastCandidates {D M : ℕ} (dims : Fin (D + 1) → ℕ)
    (ranks : Fin (D + 1) → Fin (D + 1) → ℕ) : Finset (Multiplicities D M) :=
  (Finset.univ.filter (fun f : Assignment D M => Feasible dims ranks (histogram f))).image histogram
\end{LeanCode}
\leansource{15}{ExactFitFull_2026_10_28_UTC.lean}{5478}{5493}{b362a3ee35df}
\emph{Translation.} The candidate set is the set of histograms of maps $\ii0{M_{\mathrm{slot}}-1}\to\{\text{valid keys}\}\cup\{\text{none}\}$ that are feasible. Every barcode with at most $M_{\mathrm{slot}}$ slots is such a histogram. Hence, with zero lower bounds, the candidate set is exactly $\Zb(n-r)$.

\emph{The coefficient, the threshold, the weight and the edges.}
\begin{LeanCode}
def rankInput (h : Multiplicities D M) (i j : Fin (D + 1)) : ℕ :=
  if i ≤ j then pathCount h i j else 0

def Matches (h : Multiplicities D M)
    (ranks : Fin (D + 1) → Fin (D + 1) → ℕ) : Prop :=
  ∀ i j : Fin (D + 1), i ≤ j → pathCount h i j = ranks i j

instance matchesDecidable (h : Multiplicities D M)
    (ranks : Fin (D + 1) → Fin (D + 1) → ℕ) : Decidable (Matches h ranks) := by
  unfold Matches
  infer_instance

def RankLE (h k : Multiplicities D M) : Prop :=
  ∀ i j : Fin (D + 1), i ≤ j → pathCount h i j ≤ pathCount k i j

def RankLT (h k : Multiplicities D M) : Prop :=
  RankLE h k ∧ ∃ i j : Fin (D + 1), i ≤ j ∧ pathCount h i j < pathCount k i j

instance rankLEDecidable (h k : Multiplicities D M) : Decidable (RankLE h k) := by
  unfold RankLE
  infer_instance
instance rankLTDecidable (h k : Multiplicities D M) : Decidable (RankLT h k) := by
  unfold RankLT
  infer_instance

/-- Independently enumerated balanced-output objective above a given array.
The supplied list may be replaced by any proved-equal faster enumeration. -/
def higherOutputValues (ν : Finset (Multiplicities D M)) (h : Multiplicities D M) : Finset ℕ :=
  (ν.filter (fun k => RankLE h k ∧ Balanced k)).image outputCount

/-- Total arithmetic coefficient. Off the feasible domain an empty list
returns zero; the actual measured adapter proves nonemptiness on-domain. -/
def coefficient (ν : Finset (Multiplicities D M)) (h : Multiplicities D M) : ℕ :=
  (higherOutputValues ν h).min.untopD 0

/-- Exact rational canonical potential threshold, before any real cast. -/
def threshold (h : Multiplicities D M) : ℚ :=
  (∑ p : MAIS.IntervalCochainProduct.NormalIndex
      (@birth D M h) (@death D M h),
    MAIS.IntervalPositivePotentials.edgeWeight D (@birth D M h) (@death D M h) p) / 2

/-- A critical source contributes one precisely at exact rational equality.
This equality test never compares arbitrary real numbers. -/
def weight (ν : Finset (Multiplicities D M)) (h : Multiplicities D M) : ℕ :=
  if (coefficient ν h : ℚ) / 2 = threshold h then 1 else 0

def Edge (ν : Finset (Multiplicities D M)) (h k : Multiplicities D M) : Prop :=
  RankLT h k ∧ coefficient ν h = coefficient ν k
\end{LeanCode}
\leansource{16}{ExactFitFull_2026_10_28_UTC.lean}{5328}{5375}{b362a3ee35df}
\begin{LeanCode}
def Linked (r s : S) : Prop :=
  birth s < birth r ∧ birth r ≤ death s ∧ death s < death r

instance linked_decidable (r s : S) : Decidable (Linked birth death r s) := by
  unfold Linked
  infer_instance

abbrev NormalIndex := {p : S × S // Linked birth death p.1 p.2}
abbrev NormalParams := NormalIndex birth death → ℝ
\end{LeanCode}
\leansource{17}{ExactFitFull_2026_10_28_UTC.lean}{4817}{4825}{b362a3ee35df}
\begin{LeanCode}
variable {S : Type*} (D : ℕ) (birth death : S → ℕ)

def denominator (s : S) : ℕ := birth s + (D + 1 - death s)

/-- Rational interval potential, before clearing any finite denominators. -/
def potential (s : S) : ℚ := (birth s : ℚ) / (denominator D birth death s : ℚ)

def edgeWeight (p : NormalIndex birth death) : ℚ :=
  potential D birth death p.val.1 - potential D birth death p.val.2
\end{LeanCode}
\leansource{18}{ExactFitFull_2026_10_28_UTC.lean}{5009}{5017}{b362a3ee35df}
\emph{Translation.} For $h,k$ in the candidate set $\nu$:
\begin{itemize}
\item \lean{RankLE h k} is $h\preceq k$ (all candidates have the same dimension vector), and \lean{RankLT h k} is $h\prec k$.
\item \lean{coefficient ν h} is $\min\{\cout(k):k\in\nu,\ h\preceq k,\ k\text{ balanced}\}$, with the value $0$ if this set is empty.
\item \lean{potential} is the canonical potential $u_s=b_s/(b_s+D-e_s)$, since $D+1-\text{death}=D-e_s$. \lean{edgeWeight} of a linked pair $s\lin s'$ is $u_{s'}-u_s$, so \lean{threshold h} is $T(u)$ for the canonical potential, as an exact rational.
\item \lean{weight ν h} is $1$ if $\text{coefficient}/2=T(u)$, and $0$ otherwise.
\item \lean{Edge ν h k} means $h\prec k$ and equal coefficients.
\end{itemize}

\emph{The longest weighted path and the output.}
\begin{LeanCode}
abbrev Node (ν : Finset (Multiplicities D M)) := ↥ν

def nodeEdge (ν : Finset (Multiplicities D M)) (h k : Node ν) : Prop := Edge ν h.val k.val
instance nodeEdgeDecidable (ν : Finset (Multiplicities D M)) : DecidableRel (nodeEdge ν) := by
  intro h k
  exact edgeDecidable ν h.val k.val

def nodeWeight (ν : Finset (Multiplicities D M)) (h : Node ν) : ℕ := weight ν h.val

abbrev PathIndex (D : ℕ) := {p : Fin (D + 1) × Fin (D + 1) // p.1 ≤ p.2}

/-- Strict rank-order edges strictly increase this literal finite sum. -/
def height (h : Multiplicities D M) : ℕ :=
  ∑ p : PathIndex D, pathCount h p.val.1 p.val.2

def capacity (D M : ℕ) : ℕ := Fintype.card (PathIndex D) * M

def defect (h : Multiplicities D M) : ℕ := capacity D M - height h

/-- A fully executable finite-arithmetic longest-path candidate. Its fuel
is an architecture-independent finite bound for all feasible arrays. -/
def criticalDepth (ν : Finset (Multiplicities D M)) (h : Node ν) : ℕ :=
  MAIS.Run2ArithmeticCriticalPathDynamic.depthFuel (nodeEdge ν) (nodeWeight ν)
    (capacity D M) h

/-- Complete proof-free algorithm output from a supplied integer rank
table: all coefficient/multiplicity pairs at matching finite sources.
Actual correctness later proves this output is a nonempty singleton. -/
def computedPairs (ν : Finset (Multiplicities D M))
    (ranks : Fin (D + 1) → Fin (D + 1) → ℕ) : Finset (ℕ × ℕ) :=
  (ν.attach.filter (fun h => Matches h.val ranks)).image
    (fun h => (coefficient ν h.val, 1 + criticalDepth ν h))
\end{LeanCode}
\leansource{19}{ExactFitFull_2026_10_28_UTC.lean}{5384}{5415}{b362a3ee35df}
\begin{LeanCode}
/-- Maximum source-weight sum over paths using at most the supplied number
of edges. Stopping at the current vertex always has cost zero. -/
def depthFuel (edge : J → J → Prop) [DecidableRel edge]
    (weight : J → ℕ) : ℕ → J → ℕ
  | 0, _ => 0
  | n + 1, i => Finset.univ.sup (fun j =>
      if edge i j then weight i + depthFuel edge weight n j else 0)
\end{LeanCode}
\leansource{20}{ExactFitFull_2026_10_28_UTC.lean}{5226}{5232}{b362a3ee35df}
\emph{Translation.} \lean{depthFuel edge w k i} is the largest value of $w(i_0)+\dots+w(i_{l-1})$ over paths $i=i_0\to i_1\to\dots\to i_l$ along edges with $l\le k$. Stopping contributes $0$, so the last vertex is not weighted. \lean{computedPairs ν R'} is the set of pairs $(\text{coefficient}(h),\,1+\text{criticalDepth}(h))$ over the $h\in\nu$ whose rank table is $R'$.

\begin{proposition}\label{prop:prog-translation}
Let $D\ge1$ and $A\in\Rep(n)$. The formal program on the input $(n,R(A))$ returns $\Prog(n,R(A))=\{(\creg(n,r)+C(h_A),\,1+F(h_A))\}$ (Section~\ref{sec:statements}).
\end{proposition}
\begin{proof}
Put $r=R_{0D}(A)$. By the translations above, the candidate set is $\Zb(n-r)$, and the only candidate $h$ with $\rho(h)=(R_{ij}(A)-r)_{i\le j}$ is $h_A$, by~\eqref{eq:zero-output-part} and Fact~\ref{fact:gabriel}(c). The coefficient is $C$, because the minimized set is nonempty (Lemma~\ref{lem:nonempty}, a combinatorial statement). The weight of $h$ is $1$ if and only if $h$ is critical: if $h$ is balanced then $2T(u)=\cout(h)$ by~\eqref{eq:telescope}; if not, then $C(h)\le C^\sqsubseteq(h)=2\lambda_h<2T(u)$. Here the inequality holds because the set minimized in $C^\sqsubseteq(h)$ is contained in the one minimized in $C(h)$, by Proposition~\ref{prop:closure}(a); the equality is Theorem~\ref{thm:coefficient}; and the strict inequality is Lemma~\ref{lem:stationarity}, since the canonical potential is admissible (Section~\ref{sec:zero}). Along an edge, $\sum_{i\le j}\rho_{ij}$ strictly increases and stays at most $\#\{(i,j):\ 0\le i\le j\le D\}\cdot M_{\mathrm{slot}}$, so the fuel bounds the length of every path and \lean{criticalDepth} is the unbounded maximum $F(h_A)$.
\end{proof}

Besides the translations, the proof of Proposition~\ref{prop:prog-translation} uses combinatorial facts proved in the text (Lemma~\ref{lem:nonempty}, \eqref{eq:zero-output-part}, Fact~\ref{fact:gabriel}(c), \eqref{eq:telescope} and the bound on the fuel), and, in one place, three geometric results of the text: Proposition~\ref{prop:closure}(a), Theorem~\ref{thm:coefficient} and Lemma~\ref{lem:stationarity}. They serve only to see that the formal weight agrees with ``critical'' for unbalanced barcodes. A reader who prefers not to rely on them can read Theorem~\ref{thm:B} with ``critical'' replaced by the formal test $C(h)=2T(u)$, where $u$ is the canonical potential of $h$. Then \ref{F:prog}--\ref{F:order}, the translations of this appendix and these combinatorial facts certify the displayed formulas of Theorem~\ref{thm:B}.
\section{The rank order is generated by switches}\label{app:generation}

\begin{proposition}\label{prop:generation}
Let $D\in\N^*$, and let $h\preceq k$ be barcodes of length $D$. There is a finite sequence $h=h_0,h_1,\dots,h_N=k$ such that each $h_{t+1}$ is obtained from $h_t$ by the switch of Lemma~\ref{lem:switch} along a linked pair of $h_t$.
\end{proposition}

The statement is purely combinatorial; for its use see Proposition~\ref{prop:closure}(b). Over algebraically closed fields the equality of the rank order and the degeneration order is due to Abeasis, Del Fra and Kraft~\cite{abeasis-delfra-kraft}. The proof below reduces to the Bruhat order of the symmetric group by ``standardization''.

\subsection*{Padded matrices}
Fix $D\in\N^*$ and $N_0\in\N$. For a barcode $h$ of length $D$ and $a\in\ii0D$ let $\mathrm{st}_a(h)=\#\{s\in\Sl(h):\ b_s=a\}$ be the number of slots starting at $a$, and assume $\mathrm{st}_a(h)\le N_0$ for all $a$. Define the $(D+2)\times(D+2)$ matrix $X=X(h)$, with rows and columns indexed by $\ii0{D+1}$, by
\[
X(b,c)=h(\ii b{c-1})\ (b<c),\qquad X(a,a)=N_0-\mathrm{st}_a(h)\ (a\le D),\qquad X(D+1,D+1)=N_0,
\]
and $X(b,c)=0$ for $c<b$. So the number of slots with interval $\ii b{c-1}$ is recorded in position $(b,c)$, and the diagonal entries are ``padding''. Every row sum of $X$ is $N_0$. Let $n=\dim(h)$ and put $n_{-1}=n_{D+1}=0$. The column sums are
\[
\kappa_c=N_0-(n_c-n_{c-1})\quad(c\in\ii0D),\qquad\kappa_{D+1}=N_0+n_D,
\]
because $\mathrm{st}_c(h)-\#\{s\in\Sl(h):\ e_s=c-1\}=n_c-n_{c-1}$. For $p,q\in\ii0{D+1}$ put
\[
\mathrm{NE}_M(p,q)=\sum_{b\le p,\ c\ge q}M(b,c)
\]
for any matrix $M$. Then:
\begin{enumerate}[label=\textup{(\alph*)}]
\item $\mathrm{NE}_X(i,j+1)=\rho_{ij}(h)$ for $0\le i\le j\le D$. The diagonal entries do not contribute, since $a\le i<j+1\le a$ is impossible.
\item $\mathrm{NE}_X(p,q)=\sum_{c\ge q}\kappa_c-\sum_{b>p}N_0$ for $p\ge q$, since every entry with $b>p\ge q$ has $c\ge b>p$.
\end{enumerate}
Hence, for barcodes $h,k$ of length $D$ with the same dimension vector (and $N_0\ge\mathrm{st}_a(h),\mathrm{st}_a(k)$ for all $a$), $X(h)$ and $X(k)$ have the same row and column sums, and $h\preceq k$ if and only if $\mathrm{NE}_{X(h)}\le\mathrm{NE}_{X(k)}$ entrywise.

\subsection*{Moves}
For $b,c\in\ii0{D+1}$ let $E_{bc}$ be the $(D+2)\times(D+2)$ matrix with a single nonzero entry, equal to $1$, at $(b,c)$. For a $(D+2)\times(D+2)$ matrix $M$ with nonnegative integer entries, indices $b_1<b_2$, $c_1<c_2$ with $M(b_1,c_1)\ge1$ and $M(b_2,c_2)\ge1$, the \emph{move}
\[
M'=M-E_{b_1c_1}-E_{b_2c_2}+E_{b_1c_2}+E_{b_2c_1}
\]
satisfies $\mathrm{NE}_{M'}(p,q)=\mathrm{NE}_M(p,q)+1$ if $b_1\le p<b_2$ and $c_1<q\le c_2$, and $\mathrm{NE}_{M'}(p,q)=\mathrm{NE}_M(p,q)$ otherwise. This follows from $([b_1\le p]-[b_2\le p])([c_2\ge q]-[c_1\ge q])\ge0$.

\begin{lemma}\label{lem:move-switch}
Let $M=X(h)$, and let $M'$ be the move of $M$ with indices $b_1<b_2$ and $c_1<c_2$. Suppose that $M'$ is upper triangular ($M'(b,c)=0$ for $c<b$). Then $b_1<b_2\le c_1<c_2$, the intervals $I=\ii{b_1}{c_1-1}$ and $I'=\ii{b_2}{c_2-1}$ satisfy $h(I)\ge1$ and $h(I')\ge1$, every pair of slots $(s,s')$ with $I_s=I$ and $I_{s'}=I'$ is linked, and $M'=X(h^\pi)$ for $\pi=(s,s')$ and the switched barcode $h^\pi$ of Lemma~\ref{lem:switch}.
\end{lemma}
\begin{proof}
Since $M'(b_2,c_1)\ge1$ and $M'$ is upper triangular, $b_2\le c_1$. Hence $b_1<b_2\le c_1<c_2$, so $(b_1,c_1)$ and $(b_2,c_2)$ are strictly above the diagonal, and $h(I)=M(b_1,c_1)\ge1$, $h(I')=M(b_2,c_2)\ge1$. For slots $s,s'$ with $I_s=I$ and $I_{s'}=I'$, the linking condition $b_s<b_{s'}\le e_s+1\le e_{s'}$ reads $b_1<b_2\le c_1\le c_2-1$, which holds. The move adds $1$ at $(b_1,c_2)$, i.e.\ to the multiplicity of $\ii{b_1}{c_2-1}=\ii{b_s}{e_{s'}}$, and $1$ at $(b_2,c_1)$: this is the multiplicity of $\ii{b_2}{c_1-1}=\ii{b_{s'}}{e_s}$ if $b_2<c_1$, and the padding entry at $(b_2,b_2)$ if $b_2=c_1$, i.e.\ if $b_{s'}=e_s+1$. In the latter case $\mathrm{st}_{b_2}$ decreases by one, consistently with the padding. Comparing with the definition of $h^\pi$ in Lemma~\ref{lem:switch} gives $M'=X(h^\pi)$.
\end{proof}

\begin{lemma}\label{lem:between}
Let $h\preceq k$ be barcodes of length $D$ with dimension vector $n$, let $N_0\ge\mathrm{st}_a(h),\mathrm{st}_a(k)$ for all $a$, and put $X=X(h)$, $Y=X(k)$. If a $(D+2)\times(D+2)$ matrix $M$ has nonnegative integer entries, the same row and column sums as $X$, and $\mathrm{NE}_X\le\mathrm{NE}_M\le\mathrm{NE}_Y$, then $M=X(h'')$ for a barcode $h''$ with $h\preceq h''\preceq k$.
\end{lemma}
\begin{proof}
For $c\in\ii0D$,
\[
\mathrm{NE}_M(c,c+1)=\sum_{b\le c}N_0-\sum_{c'\le c}\kappa_{c'}+\sum_{b>c,\ c'\le c}M(b,c').
\]
The left side is squeezed between $\mathrm{NE}_X(c,c+1)=n_c=\mathrm{NE}_Y(c,c+1)$, and the last sum vanishes for $X$. Hence $M(b,c')=0$ whenever $b>c'$ (take $c=c'$), so $M$ is upper triangular. Define $h''(\ii b{c-1})=M(b,c)$ for $b<c$. Row sums give $M(a,a)=N_0-\mathrm{st}_a(h'')$, so $M=X(h'')$; the equalities $\mathrm{NE}_M(c,c+1)=n_c$ give $\dim(h'')=n$, and (a) gives $h\preceq h''\preceq k$.
\end{proof}

\subsection*{Permutations}
Let $N_1=(D+2)N_0$. Index rows by pairs $(b,\alpha)$, $b\in\ii0{D+1}$, $\alpha\in\ii1{N_0}$, and columns by pairs $(c,\gamma)$, $\gamma\in\ii1{\kappa_c}$, both ordered lexicographically and identified with $\ii1{N_1}$. A permutation $w$ of $\ii1{N_1}$ is identified with the $0/1$ matrix having ones at $(a,w(a))$. Its \emph{projection} $\mathrm{pr}(w)$ is the $(D+2)\times(D+2)$ matrix counting the ones in each block (rows of $b$)$\times$(columns of $c$). Put $w[p,q]=\#\{a\le p:\ w(a)\ge q\}$. By~\cite[Thm.~2.1.5]{bjorner-brenti}, the Bruhat order satisfies
\begin{equation}\label{eq:tableau}
u\le w\iff u[p,q]\le w[p,q]\text{ for all }p,q .
\end{equation}
By definition~\cite[Def.~2.1.1]{bjorner-brenti}, $u<w$ means that there is a chain $u=u_0,u_1,\dots,u_\ell=w$ where each $u_{t+1}$ is obtained from $u_t$ by exchanging the values at two positions $a<a'$ with $u_t(a)<u_t(a')$. The projection is monotone: $\mathrm{NE}_{\mathrm{pr}(w)}(b,c)=w[p_b,q_c]$ with $p_b$ the last row of block $b$ and $q_c$ the first column of block $c$. Each step $u_t\to u_{t+1}$ projects either to an equality (if $a,a'$ lie in the same row block or $u_t(a),u_t(a')$ in the same column block) or to a move.

\begin{lemma}[Lifting]\label{lem:lifting}
Let $X,Y$ have nonnegative integer entries, row sums $N_0$ and column sums $\kappa$, with $\mathrm{NE}_X\le\mathrm{NE}_Y$. Then there are permutations $u\le w$ with $\mathrm{pr}(u)=X$ and $\mathrm{pr}(w)=Y$.
\end{lemma}
\begin{proof}
\emph{Rows.} Let $X^\uparrow$ be the $N_1\times(D+2)$ matrix with one unit in each row, obtained by giving the $N_0$ rows of block $b$ the $N_0$ units of row $b$ of $X$ in increasing order of column. Let $Y^\downarrow$ be defined likewise from $Y$, in decreasing order. Put $\mathrm{NE}_Z(p,c)=\#\{\text{units at rows}\le p,\text{ columns}\ge c\}$. Fix $c$ and $b$, and for $\alpha\in\ii0{N_0}$ let $p=(b,\alpha)$ denote the $\alpha$-th row of block $b$ ($\alpha=0$: the end of block $b-1$). With $x$ (resp.\ $y$) the number of units in row $b$ of $X$ (resp.\ $Y$) in columns $\ge c$,
\begin{align*}
\mathrm{NE}_{X^\uparrow}((b,\alpha),c)&=\mathrm{NE}_X(b-1,c)+\max(0,\alpha-N_0+x),\\
\mathrm{NE}_{Y^\downarrow}((b,\alpha),c)&=\mathrm{NE}_Y(b-1,c)+\min(\alpha,y).
\end{align*}
The first is convex and the second concave in $\alpha$. At $\alpha=0$ they equal $\mathrm{NE}_X(b-1,c)\le\mathrm{NE}_Y(b-1,c)$, and at $\alpha=N_0$ they equal $\mathrm{NE}_X(b,c)\le\mathrm{NE}_Y(b,c)$. A convex function lies below its chord and a concave one above it, so the first is $\le$ the second for all $\alpha\in\ii0{N_0}$ (the right-hand sides of the two displays make sense for real $\alpha\in[0,N_0]$, where they are respectively convex and concave).

\emph{Columns.} Let $u$ be obtained from $X^\uparrow$ by giving the $\kappa_c$ columns of block $c$ the units of column $c$ in increasing order of row, and $w$ from $Y^\downarrow$ in decreasing order of row. These are permutation matrices with $\mathrm{pr}(u)=X$ and $\mathrm{pr}(w)=Y$. Fix a row $p$ and a column block $c$. With $z$ (resp.\ $z'$) the number of units of column $c$ of $X^\uparrow$ (resp.\ $Y^\downarrow$) in rows $\le p$, and $\gamma\in\ii1{\kappa_c+1}$ (where $(c,\kappa_c+1)$ stands for the first column of block $c+1$),
\begin{align*}
u[p,(c,\gamma)]&=\mathrm{NE}_{X^\uparrow}(p,c+1)+\max(0,z-\gamma+1),\\
w[p,(c,\gamma)]&=\mathrm{NE}_{Y^\downarrow}(p,c+1)+\min(z',\kappa_c-\gamma+1).
\end{align*}
The right-hand sides, extended to real $\gamma\in[1,\kappa_c+1]$, are respectively convex and concave in $\gamma$. At $\gamma=1$ they equal $\mathrm{NE}_{X^\uparrow}(p,c)$ and $\mathrm{NE}_{Y^\downarrow}(p,c)$, and at $\gamma=\kappa_c+1$ they equal the values at $c+1$. The row step gives the inequality at both endpoints, hence everywhere. By~\eqref{eq:tableau}, $u\le w$.
\end{proof}

\begin{proof}[Proof of Proposition~\ref{prop:generation}]
Choose $N_0\ge\max_a\max(\mathrm{st}_a(h),\mathrm{st}_a(k))$, put $X=X(h)$ and $Y=X(k)$, and take $u\le w$ from Lemma~\ref{lem:lifting}. Let $u=u_0,\dots,u_\ell=w$ be a chain as above, and $M_t=\mathrm{pr}(u_t)$. By monotonicity $X=M_0$, $\mathrm{NE}_X\le\mathrm{NE}_{M_t}\le\mathrm{NE}_Y$, and $M_\ell=Y$. By Lemma~\ref{lem:between}, $M_t=X(h_t)$ for barcodes $h_t$ with $h\preceq h_t\preceq k$. Consecutive matrices are equal or differ by a move between upper triangular matrices, which by Lemma~\ref{lem:move-switch} is a switch along a linked pair. Deleting repetitions gives the required sequence.
\end{proof}

The formal development contains an independent Lean proof of the same statement, also by reduction to permutations (\lean{MAIS.LinearChain.Barcode.actual_rank_order_generated}), and of the resulting Euclidean closure order over $\R$ (\lean{MAIS.RealTypeAClosure.mem_closure_orbit_iff_path_rank_le}).
\begin{thebibliography}{99}
\bibitem{abeasis-delfra-kraft} S.~Abeasis, A.~Del~Fra, H.~Kraft, \emph{The geometry of representations of $A_m$}, Math.\ Ann.\ \textbf{256} (1981), 401--418.
\bibitem{aoyagi-watanabe} M.~Aoyagi, S.~Watanabe, \emph{Stochastic complexities of reduced rank regression in Bayesian estimation}, Neural Networks \textbf{18} (2005), 924--933.
\bibitem{aoyagi-2024} M.~Aoyagi, \emph{Consideration on the learning efficiency of multiple-layered neural networks with linear units}, Neural Networks \textbf{172} (2024), 106132.
\bibitem{bierstone-milman} E.~Bierstone, P.~D.~Milman, \emph{Semianalytic and subanalytic sets}, Publ.\ Math.\ IHÉS \textbf{67} (1988), 5--42.
\bibitem{bjorner-brenti} A.~Björner, F.~Brenti, \emph{Combinatorics of Coxeter Groups}, Graduate Texts in Mathematics \textbf{231}, Springer, 2005.
\bibitem{crawley-boevey} W.~Crawley-Boevey, \emph{Decomposition of pointwise finite-dimensional persistence modules}, J.\ Algebra Appl.\ \textbf{14} (2015), 1550066.
\bibitem{deconcini-procesi} C.~De Concini, C.~Procesi, \emph{Wonderful models of subspace arrangements}, Selecta Math.\ (N.S.) \textbf{1} (1995), 459--494.
\bibitem{falconer} K.~Falconer, \emph{Fractal Geometry: Mathematical Foundations and Applications}, 3rd ed., Wiley, 2014.
\bibitem{feld-o77} N.~Feld, \emph{MAIS-O77: learning coefficients on fibers and saddles}, blueprint, lectures and Lean development, October 2026, \url{https://github.com/nfeld/mais-o77-lean}.
\bibitem{feld-deep} N.~Feld, \emph{Deep linear exact fits: Lean development and verification materials}, standalone proof \lean{ExactFitFull_2026_10_28_UTC.lean}, SHA-256 beginning \lean{b362a3ee35df}, October 2026, \url{https://github.com/nfeld/deep-linear-exact-fits}.
\bibitem{ferry-monod-yano} K.~Ferry, A.~Monod, K.~Yano, \emph{Singular geometry of low-rank adaptation}, arXiv:2610.32854 (2026).
\bibitem{furman-lau} Z.~Furman, E.~Lau, \emph{Estimating the local learning coefficient at scale}, arXiv:2402.03698\DIFaddbegin \DIFadd{v1 }\DIFaddend (2024)\DIFaddbegin \DIFadd{; withdrawn, expanded and merged into~\mbox{%DIFAUXCMD
\cite{lau-llc}}\hskip0pt%DIFAUXCMD
}\DIFaddend .
\bibitem{gabriel} P.~Gabriel, \emph{Unzerlegbare Darstellungen I}, Manuscripta Math.\ \textbf{6} (1972), 71--103.
\bibitem{helgason} S.~Helgason, \emph{Differential Geometry, Lie Groups, and Symmetric Spaces}, Academic Press, 1978.
\bibitem{hironaka} H.~Hironaka, \emph{Resolution of singularities of an algebraic variety over a field of characteristic zero: I, II}, Ann.\ of Math.\ \textbf{79} (1964), 109--203, 205--326.
\bibitem{hitchcock-hoogland} R.~Hitchcock, J.~Hoogland, \emph{From global to local: a scalable benchmark for local posterior sampling}, arXiv:2507.21449v2 (2026).
\bibitem{koncki-rimanyi} J.~Koncki, R.~Rimányi, \emph{The main reasons for matrices multiplying to zero}, arXiv:2503.13326 (2025).
\bibitem{kosta-windisch} D.~Kosta, D.~Windisch, \emph{Classification of real hyperplane singularities by real log canonical thresholds}, arXiv:2411.13392v2 (2026).
\bibitem{kurumadani} Y.~Kurumadani, \emph{Real log canonical thresholds at non-singular points}, arXiv:2408.13030 (2024).
\bibitem{kurumadani-upper} Y.~Kurumadani, \emph{Upper bounds for local learning coefficients of three-layer neural networks}, arXiv:2603.12785v2 (2026).
\bibitem{lau-llc} E.~Lau, Z.~Furman, G.~Wang, D.~Murfet, S.~Wei, \emph{The local learning coefficient: a singularity-aware complexity measure}, arXiv:2308.12108v2 (2024).
\bibitem{lee} J.~M.~Lee, \emph{Introduction to Smooth Manifolds}, 2nd ed., Graduate Texts in Mathematics \textbf{218}, Springer, 2013.
\bibitem{lin} S.~Lin, \emph{Ideal-theoretic strategies for asymptotic approximation of marginal likelihood integrals}, J.\ Algebraic Stat.\ \textbf{8} (2017), no.\ 1, 22--55; arXiv:1003.5338v3.
\bibitem{lehalleur-rimanyi} S.~Pepin Lehalleur, R.~Rimányi, \emph{Geometry of fibers of the multiplication map of deep linear neural networks}, arXiv:2411.19920v2 (2024).
\bibitem{levine-mais} L.~Levine, \emph{Math for AI safety: an invitation for mathematicians}, arXiv:2609.15289 (2026). An August 2026 revision was distributed in~\cite{mais-repo} as the paper MAIS-P1.
\bibitem{mais-issue-3} R.~Sneiderman, \emph{MAIS-O70: candidate complete solution, local RLCT table for reduced-rank regression}, issue~\#3 of the MAIS repository, 3 August 2026, \url{https://github.com/Math-for-AI-Safety/MAIS/issues/3}.
\bibitem{mais-issue-12} S.~Novikov (with OpenAI Codex), \emph{MAIS-O77: candidate complete solution --- fiber rank strata and the universal saddle pair}, issue~\#12 of the MAIS repository, 4 August 2026, \url{https://github.com/Math-for-AI-Safety/MAIS/issues/12}.
\bibitem{mais-o70} Math for AI Safety, \emph{Problem MAIS-O70: local learning coefficients of reduced-rank regression}, \url{https://github.com/Math-for-AI-Safety/MAIS/blob/main/open-problems/MAIS-O70.md}.
\bibitem{mais-o77} Math for AI Safety, \emph{Problem MAIS-O77: learning coefficients on the matrix factorization fiber and saddles}, \url{https://github.com/Math-for-AI-Safety/MAIS/blob/main/open-problems/MAIS-O77.md}.
\bibitem{mais-repo} Math for AI Safety, \emph{Open problems and agendas}, \url{https://github.com/Math-for-AI-Safety/MAIS}.
\bibitem{mattila} P.~Mattila, \emph{Geometry of Sets and Measures in Euclidean Spaces}, Cambridge University Press, 1995.
\bibitem{mustata-arrangements} M.~Mustaţă, \emph{Multiplier ideals of hyperplane arrangements}, Trans.\ Amer.\ Math.\ Soc.\ \textbf{358} (2006), 5015--5023.
\bibitem{saito} M.~Saito, \emph{On real log canonical thresholds}, arXiv:0707.2308 (2007).
\bibitem{shewchuk-bhattacharya} J.~R.~Shewchuk, S.~Bhattacharya, \emph{The geometry of the set of equivalent linear neural networks}, arXiv:2404.14855v2 (2026); see also \emph{The hierarchy of manifolds in a stratification of the set of equivalent linear neural networks}, in: 42nd International Symposium on Computational Geometry (SoCG 2026), LIPIcs \textbf{367} (2026), 91:1--91:19.
\bibitem{teitler} Z.~Teitler, \emph{A note on Mustaţă's computation of multiplier ideals of hyperplane arrangements}, Proc.\ Amer.\ Math.\ Soc.\ \textbf{136} (2008), 1575--1579.
\bibitem{voigt} D.~Voigt, \emph{Induzierte Darstellungen in der Theorie der endlichen, algebraischen Gruppen}, Lecture Notes in Mathematics \textbf{592}, Springer, 1977.
\bibitem{watanabe-book} S.~Watanabe, \emph{Algebraic Geometry and Statistical Learning Theory}, Cambridge University Press, 2009.
\bibitem{watanabe-2018} S.~Watanabe, \emph{Mathematical Theory of Bayesian Statistics}, CRC Press, 2018.
\end{thebibliography}
\end{document}
